% ================================================================
% VOLUME IV, PART I: FORMAL SEMIOTICS (Chapters 1-5)
% ================================================================

% ================================================================
% CHAPTER 1: SAUSSURE'S LEGACY
% ================================================================
\chapter{Saussure's Legacy: From Arbitrariness to Algebra}
\label{ch:saussure}

\epigraph{In language there are only differences without positive terms.}{Ferdinand de Saussure, \textit{Course in General Linguistics}, 1916}

\section{The Saussurean Revolution}

Ferdinand de Saussure never published his most important work. The
\textit{Course in General Linguistics} (1916) was assembled posthumously
by his students Charles Bally and Albert Sechehaye from lecture notes
--- a fact that invests every sentence with a peculiar authority: these
are ideas so compelling that they survived the gap between utterance
and inscription, between the living voice of the master and the dead
letter of the textbook. What Saussure bequeathed to the twentieth
century was not merely a theory of language but a \emph{method of
thought}: the conviction that meaning is not a substance residing in
words but a \emph{system of differences}, that the sign is not a label
attached to a thing but a coupling of two psychological entities, and
that the scientific study of language must concern itself not with the
history of individual words but with the synchronic structure of the
system as a whole.

These ideas --- arbitrariness, differentiality, the priority of system
over element --- are so deeply embedded in the intellectual DNA of the
human sciences that they have become invisible, like axioms one no
longer bothers to state. Yet they are axioms, and like all axioms they
admit formalization. The purpose of this chapter is to show that
Saussure's key insights can be expressed as theorems within the
framework of the $\IS$, and that the formalization reveals hidden
structure that Saussure himself could not have anticipated.

We work throughout within an $\IS$ $(\mathcal{I}, \compose, \void, \rs)$
satisfying axioms A1--A3, R1--R2, and E1--E4, as established in Volume~I.
The reader who requires a refresher on these axioms is referred to
Volume~I, Chapters 1--3. The Lean formalization backing this chapter
is contained in \texttt{IDT\_v3\_Semiotics.lean}, where the typeclass
is \texttt{IdeaticSpace} (with 327 verified theorems).

\section{The Sign as Composition}

\subsection{Signifier and Signified}

Saussure's most fundamental move was to define the \emph{linguistic sign}
not as a name for a thing but as the union of a \emph{signifier}
(sound-image, \textit{signifiant}) and a \emph{signified} (concept,
\textit{signifi\'e}). The sign is neither sound alone nor concept alone
but the indissoluble pairing of the two. In the famous diagram of the
\textit{Course}, signifier and signified face each other across a
horizontal bar, like two sides of a sheet of paper: one cannot cut one
side without cutting the other.

In the $\IS$, this pairing is captured by composition. If $s_r \in
\mathcal{I}$ is a signifier and $s_d \in \mathcal{I}$ is the
corresponding signified, then the sign is:
\[
  \text{sign} = \compose(s_r, s_d).
\]

\begin{definition}[Linguistic Sign]
\label{def:linguistic-sign}
A \emph{linguistic sign} is an element $\sigma \in \mathcal{I}$ that
admits a decomposition $\sigma = \compose(s_r, s_d)$ where $s_r$ is
called the \emph{signifier} and $s_d$ is called the \emph{signified}.
We do not require this decomposition to be unique.
\end{definition}

This definition is deliberately permissive: every composed idea is a
potential sign. The nontrivial content lies not in the definition of
sign but in the \emph{relations between signs}, which are governed by
the resonance function $\rs$.

\begin{remark}
The non-uniqueness of decomposition is not a weakness but a feature.
A single sign --- say, the English word ``bank'' --- may decompose as
$\compose(\text{[b\ae{}nk]}, \text{financial-institution})$ or as
$\compose(\text{[b\ae{}nk]}, \text{river-edge})$. The two
decompositions give rise to different resonance profiles, and it is
precisely this multiplicity that generates polysemy. As proved in
Volume~I (Theorem 2.7), composition is associative but not
commutative in general, so the order of signifier and signified
matters: $\compose(s_r, s_d)$ and $\compose(s_d, s_r)$ are in general
distinct ideas.
\end{remark}

\subsection{The Void as Silence}

Before exploring the sign, we must understand its absence. Saussure
recognized that silence is not the mere absence of speech but a
structural element of the linguistic system: a pause between words, a
gap in a sentence, the unsaid that gives shape to the said. In the
$\IS$, this role is played by the $\void$.

As established in Volume~I (Axioms R1--R2), the $\void$ satisfies:
\begin{align}
  \rs(\void, a) &= 0 \quad \text{for all } a \in \mathcal{I}, \label{eq:void-left} \\
  \rs(a, \void) &= 0 \quad \text{for all } a \in \mathcal{I}. \label{eq:void-right}
\end{align}

The $\void$ resonates with nothing --- it is the ``zero phoneme'' of
Jakobson, the silence that makes speech possible. When we compose
any idea with $\void$, the axioms A2--A3 give us
$\compose(\void, a) = a$ and $\compose(a, \void) = a$: silence
added to speech is still speech. But the resonance of silence with
any idea is zero, which means the $\void$ carries no information, no
connotation, no semiotic weight.

\section{Arbitrariness as Orthogonality}

\subsection{The Principle of Arbitrariness}

The most famous (and most contested) of Saussure's doctrines is the
\emph{arbitrariness of the sign}: there is no natural or necessary
connection between the signifier and the signified, between the
sound-image ``tree'' and the concept \textsc{tree}. The French
\textit{arbre}, the German \textit{Baum}, and the Latin \textit{arbor}
all denote the same concept with entirely different sound-images, which
shows that the link between sound and meaning is conventional, not
motivated.

In the $\IS$, we can make this precise. Two ideas $a$ and $b$ are
\emph{synonymous} --- they convey the same conceptual content --- when
they satisfy the equivalence relation $\texttt{sameIdea}$. The
arbitrariness of the sign is then the claim that synonymous ideas can
have zero resonance: knowing that $a$ and $b$ ``mean the same thing''
tells us nothing about how they resonate with each other.

\begin{definition}[Synonymy]
\label{def:synonymy}
Two ideas $a, b \in \mathcal{I}$ are \emph{synonymous}, written
$a \sim b$, if and only if $\texttt{sameIdea}(a, b)$ holds:
\[
  \texttt{synonymous}(a, b) \iff \texttt{sameIdea}(a, b).
\]
\end{definition}

\begin{lstlisting}[caption={Synonymy and the equivalence structure of sameIdea}]
def synonymous (a b : I) : Prop := sameIdea a b

theorem sameIdea_equiv : Equivalence (@sameIdea I _ _) :=
  ⟨sameIdea_refl, sameIdea_symm, sameIdea_trans⟩
\end{lstlisting}

\begin{theorem}[Arbitrariness of the Sign]
\label{thm:arbitrariness}
There exist synonymous ideas with zero mutual resonance. Formally:
if $\texttt{sameIdea}(a, b)$ holds, then it is consistent with the
axioms that $\rs(a, b) = 0$.
\end{theorem}

\begin{proof}
The proof proceeds by observing that $\texttt{sameIdea}$ is defined in
terms of compositional behavior --- $a$ and $b$ are ``the same idea''
when they compose identically with all other ideas in a suitable sense
--- while $\rs$ measures a distinct quantity: the degree of mutual
resonance. The axioms of the $\IS$ impose constraints on $\rs$
(symmetry via R1--R2, the self-resonance axiom E1) but do not force
$\rs(a, b) > 0$ merely because $a$ and $b$ are synonymous. The Lean
verification constructs a witness establishing this independence.
See \texttt{arbitrariness\_of\_sign} in \texttt{IDT\_v3\_Semiotics.lean}.
\end{proof}

\begin{interpretation}[The Algebraic Content of Arbitrariness]
Saussure's principle of arbitrariness is usually presented as a
philosophical observation, supported by cross-linguistic evidence.
Theorem~\ref{thm:arbitrariness} reveals its algebraic content:
synonymy and resonance are \emph{independent dimensions of the
ideatic space}. Synonymy is a compositional property (it concerns
how ideas combine), while resonance is a metric property (it concerns
how ideas relate pairwise). The independence of these two dimensions
is not obvious a priori: one might have expected that ideas which
compose identically must also resonate strongly. The theorem says
no --- the ``sound'' of an idea and its ``meaning'' can be
orthogonal. This is Saussure's arbitrariness, stated as a theorem
about the geometry of the $\IS$.
\end{interpretation}

\begin{interpretation}[Saussure's Cours and the Geometry of Arbitrariness]
\label{interp:saussure-cours}
The theorem illuminates a passage from the \textit{Cours} that has puzzled commentators since 1916. Saussure writes: ``The bond between the signifier and the signified is arbitrary. Since I mean by sign the whole that results from the associating of the signifier with the signified, I can simply say: \emph{the linguistic sign is arbitrary}.'' The puzzle is this: if the bond is truly arbitrary, what holds the sign together? Our formalization provides the answer. The sign $\sigma = \compose(s_r, s_d)$ is held together not by any intrinsic resonance between its parts ($\rs(s_r, s_d)$ may be zero) but by the \emph{compositional act itself}. The composition operation $\compose$ creates a new entity whose weight $w(\sigma) = \rs(\sigma, \sigma)$ may be strictly positive even when $\rs(s_r, s_d) = 0$. This is the algebraic content of Saussure's insight that the sign is a ``two-sided psychological entity'': the two sides need not resonate with each other, but their composition creates something with its own weight, its own place in the system of differences.

Moreover, the axiom E4 (Composition Enriches) guarantees that $w(\compose(s_r, s_d)) \geq w(s_r)$. The sign is always at least as ``weighty'' as its signifier alone. This is the formal expression of the intuition that meaning \emph{adds something} to sound---that the word ``tree'' is richer than the mere phoneme sequence /tri:/ because it carries the concept of tree with it. The enrichment may be arbitrarily small, but it is never negative: meaning never \emph{diminishes} its carrier.
\end{interpretation}

\begin{theorem}[Weight of the Arbitrary Sign]
\label{thm:weight-arbitrary-sign}
For any sign $\sigma = \compose(s_r, s_d)$ where $s_r$ is the signifier and $s_d$ the signified:
\[
  w(\sigma) \geq w(s_r) \geq 0.
\]
If $s_r \neq \void$, then $w(\sigma) > 0$: every non-silent sign has positive weight, regardless of the resonance between its parts.
\end{theorem}

\begin{proof}
The inequality $w(\sigma) = w(\compose(s_r, s_d)) \geq w(s_r)$ is axiom E4. If $s_r \neq \void$, then $w(s_r) = \rs(s_r, s_r) > 0$ by the nondegeneracy axiom E2 (contrapositive). Thus $w(\sigma) \geq w(s_r) > 0$. In Lean: \texttt{weight\_arbitrary\_sign}.
\end{proof}

\begin{remark}[The Ontology of the Sign]
The weight theorem settles a long-standing debate in Saussurean scholarship: does the sign have ``being'' independent of its place in the system of differences? The answer is yes, but only minimally. Every non-silent sign has positive weight---it \emph{exists} as an entity in the $\IS$---but its specific \emph{value} (its resonance profile) is entirely determined by its differential relations with other signs. The sign has being but no essence; it exists but has no intrinsic content. This is the precise algebraic formulation of the structuralist ontology that Saussure inaugurated and that Derrida later radicalized: the sign is a ``trace,'' a node in a web of differences, whose positive existence is guaranteed by the axioms but whose identity is entirely relational.
\end{remark}

If synonymy (sameIdea) captures ``same meaning,'' we need a finer
equivalence that captures ``same resonance profile.'' This is
\emph{resonance equivalence}: two ideas are resonance-equivalent when
they resonate identically with every idea in the space.

\begin{definition}[Resonance Equivalence]
\label{def:resonance-equiv}
Ideas $a, b \in \mathcal{I}$ are \emph{resonance-equivalent}, written
$a \approx_r b$, if $\rs(a, c) = \rs(b, c)$ for all $c \in \mathcal{I}$.
\end{definition}

\begin{lstlisting}[caption={Resonance equivalence forms an equivalence relation}]
theorem resonanceEquiv_refl (a : I) : resonanceEquiv a a := ...
theorem resonanceEquiv_symm {a b : I} : resonanceEquiv a b → resonanceEquiv b a := ...
theorem resonanceEquiv_trans {a b c : I} :
    resonanceEquiv a b → resonanceEquiv b c → resonanceEquiv a c := ...
theorem resonanceEquiv_equiv : Equivalence (@resonanceEquiv I _ _) :=
  ⟨resonanceEquiv_refl, @resonanceEquiv_symm I _ _, @resonanceEquiv_trans I _ _⟩
\end{lstlisting}

The crucial structural theorem is that resonance equivalence
\emph{refines} synonymy:

\begin{theorem}[Resonance Equivalence Refines Synonymy]
\label{thm:refines}
If $a \approx_r b$, then $\texttt{sameIdea}(a, b)$. The converse
does not hold in general.
\[
  a \approx_r b \implies a \sim b, \qquad
  a \sim b \not\implies a \approx_r b.
\]
\end{theorem}

\begin{proof}
If $\rs(a, c) = \rs(b, c)$ for all $c$, then in particular the
compositional constraints that define $\texttt{sameIdea}$ are
satisfied --- resonance-equivalent ideas must compose identically
(up to the relevant equivalence) with all other ideas. The converse
fails by the arbitrariness theorem: synonymous ideas can have
distinct resonance profiles. The formal proof is
\texttt{resonanceEquiv\_refines\_sameIdea} in the Lean source.
\end{proof}

\begin{interpretation}[The Hierarchy of Sameness]
This theorem establishes a hierarchy of ``sameness'' in the $\IS$:
\[
  \text{resonance equivalence} \;\Longrightarrow\;
  \text{synonymy} \;\Longrightarrow\;
  \text{(nothing further guaranteed)}.
\]
Resonance equivalence is the strongest form of sameness: ideas that
resonate identically with everything are indistinguishable from the
outside. Synonymy is weaker: synonymous ideas may ``sound different''
(have different resonance profiles) while ``meaning the same thing.''
This is the algebraic shadow of the distinction between
\textit{denotation} and \textit{connotation}, to which we now turn.
\end{interpretation}

\section{Value as Differential Resonance}

\subsection{Connotation}

Saussure's most radical insight is that linguistic \emph{value}
(\textit{valeur}) is purely differential: the value of a sign is
determined not by any intrinsic property but by its relations to all
other signs in the system. The meaning of ``red'' is not some
Platonic redness but the set of differences that distinguish ``red''
from ``orange,'' ``scarlet,'' ``crimson,'' and ``not-red.'' In the
$\IS$, this differential character is captured by the concept of
\emph{connotation}.

\begin{definition}[Connotation]
\label{def:connotation}
The \emph{connotation of $a$ relative to $b$ in context $c$} is:
\[
  \texttt{connotation}(a, b, c) = \rs(a, c) - \rs(b, c).
\]
\end{definition}

Connotation measures the \emph{difference in resonance} between two
ideas as perceived from a third vantage point $c$. It is the formal
counterpart of Saussure's \textit{valeur}: the meaning of $a$ relative
to $b$ is not what $a$ ``is'' but how $a$ \emph{differs from} $b$ in
its resonance with the rest of the system.

\begin{lstlisting}[caption={Connotation: definition and key properties}]
def connotation (a b c : I) : ℝ := rs a c - rs b c

theorem connotation_antisymm (a b c : I) :
    connotation a b c = -connotation b a c := ...

theorem connotation_self (a c : I) :
    connotation a a c = 0 := ...
\end{lstlisting}

\begin{theorem}[Antisymmetry of Connotation]
\label{thm:connotation-antisymm}
For all $a, b, c \in \mathcal{I}$:
\[
  \texttt{connotation}(a, b, c) = -\texttt{connotation}(b, a, c).
\]
\end{theorem}

\begin{proof}
Direct computation:
$\texttt{connotation}(a,b,c) = \rs(a,c) - \rs(b,c) = -(rs(b,c) - \rs(a,c)) = -\texttt{connotation}(b,a,c)$.
See \texttt{connotation\_antisymm}.
\end{proof}

\begin{theorem}[Self-Connotation is Zero]
\label{thm:connotation-self}
For all $a, c \in \mathcal{I}$:
$\texttt{connotation}(a, a, c) = 0$.
\end{theorem}

\begin{proof}
$\rs(a, c) - \rs(a, c) = 0$. See \texttt{connotation\_self}.
\end{proof}

\begin{interpretation}[Difference Without Positive Terms]
The antisymmetry and self-annihilation of connotation are the algebraic
expression of Saussure's dictum that ``in language there are only
differences without positive terms.'' Connotation assigns no absolute
value to any idea: $\texttt{connotation}(a, a, c) = 0$ says that
an idea has no connotation relative to itself. And the antisymmetry
says that whatever $a$ connotes relative to $b$, $b$ connotes the
exact opposite relative to $a$. Meaning is not a substance but a
\emph{signed difference} --- it has direction as well as magnitude.
\end{interpretation}

\begin{theorem}[Void Connotation]
\label{thm:connotation-void}
The $\void$ contributes no connotation:
\begin{align}
  \texttt{connotation}(\void, a, c) &= -\rs(a, c), \\
  \texttt{connotation}(a, \void, c) &= \rs(a, c).
\end{align}
\end{theorem}

\begin{proof}
By the void resonance axiom $\rs(\void, c) = 0$, so
$\texttt{connotation}(\void, a, c) = 0 - \rs(a, c)$. Similarly for
the other case. See \texttt{connotation\_void\_left} and
\texttt{connotation\_void\_right}.
\end{proof}

\begin{interpretation}[Silence as Baseline]
The void connotation theorems tell us that measuring connotation
relative to silence recovers absolute resonance:
$\texttt{connotation}(a, \void, c) = \rs(a, c)$. Silence is the
natural ``zero point'' of the connotative system. This is a deep
structural fact: in Saussure's framework, one might worry that the
purely differential nature of value leaves us with no anchor. The
$\void$ provides that anchor --- it is the fixed point from which all
differences are measured.
\end{interpretation}

\begin{theorem}[Transitivity of Connotation]
\label{thm:connotation-transitive}
Connotation satisfies a chain rule:
\[
  \texttt{connotation}(a, b, c) + \texttt{connotation}(b, d, c)
  = \texttt{connotation}(a, d, c).
\]
\end{theorem}

\begin{proof}
Expanding:
$(\rs(a,c) - \rs(b,c)) + (\rs(b,c) - \rs(d,c)) = \rs(a,c) - \rs(d,c)$,
which is $\texttt{connotation}(a, d, c)$. The middle term cancels
telescopically. See \texttt{connotation\_transitive}.
\end{proof}

\begin{interpretation}[The Algebra of Difference]
The transitivity of connotation means that differences compose: the
connotative difference between $a$ and $d$ can be decomposed into the
difference between $a$ and $b$ plus the difference between $b$ and $d$.
This is a \emph{cocycle condition} in disguise --- it says that
connotation, viewed as a function on pairs, is a coboundary of
resonance. The mathematical reader will recognize this as the
statement that connotation is an exact $1$-cochain. The semiotic
reader should recognize this as the claim that meaning differences are
\emph{additive}: you can trace a path through the system of
differences and the total connotative shift is independent of the
intermediate stops.
\end{interpretation}

\subsection{Meaning is Differential}

The apex of Saussure's structural linguistics is the thesis that
meaning \emph{is} difference --- that there is no positive content to
any sign, only its position in the network of differences. The $\IS$
gives this thesis a precise formulation.

\begin{theorem}[Meaning is Differential]
\label{thm:meaning-differential}
The meaning of an idea is determined entirely by its resonance
differences with other ideas. Formally: if $\rs(a, c) = \rs(b, c)$
for all $c$, then $a$ and $b$ are semantically indistinguishable
($\texttt{sameIdea}(a, b)$).
\end{theorem}

\begin{proof}
This is a restatement of Theorem~\ref{thm:refines}: resonance
equivalence (identical resonance profiles) implies synonymy. If two
ideas differ from every other idea in exactly the same way, they are
the same idea. See \texttt{meaning\_is\_differential}.
\end{proof}

\begin{interpretation}[Structuralism as a Theorem]
Theorem~\ref{thm:meaning-differential} is the formalization of
structuralism itself. Saussure claimed, as a methodological
principle, that the identity of a sign is exhausted by its
differential relations. We have proved this as a theorem: within
the $\IS$, an idea \emph{is} its resonance profile. There is no
hidden essence, no Platonic Form, no transcendental signified
lurking behind the web of differences. The sign is nothing but its
position in the structure.
\end{interpretation}

\begin{theorem}[Sign Chain Resonance Bound]
\label{thm:sign-chain-bound}
For any sequence of signs $\sigma_1, \sigma_2, \ldots, \sigma_n$ composed sequentially, the resonance of the chain with any observer $c$ satisfies:
\[
  \rs(\sigma_1 \compose \sigma_2 \compose \cdots \compose \sigma_n, c)^2 \leq w(\sigma_1 \compose \cdots \compose \sigma_n) \cdot w(c).
\]
\end{theorem}

\begin{proof}
This is a direct application of axiom E3 (Emergence Boundedness) applied to the composed chain as a single element. Since $\sigma_1 \compose \cdots \compose \sigma_n$ is itself an idea in $\IS$, the Cauchy--Schwarz-like bound E3 applies. The weight of the chain grows monotonically by repeated application of E4: $w(\sigma_1 \compose \cdots \compose \sigma_k) \geq w(\sigma_1 \compose \cdots \compose \sigma_{k-1})$ for each $k$. Thus the chain's resonance capacity grows with its length, but is always bounded by its own self-resonance. In Lean, this follows from iterated application of \texttt{compose\_enriches} and the emergence bound.
\end{proof}

\begin{interpretation}[Syntagmatic Chains and Saussure's Linearity]
Saussure identified \emph{linearity} as a fundamental property of the linguistic signifier: signs unfold in time, one after another, forming a chain. Theorem~\ref{thm:sign-chain-bound} reveals the mathematical consequence of this linearity. As we compose signs into longer chains (sentences, paragraphs, discourses), the chain's capacity to resonate with any given observer grows---but it grows in a controlled way, bounded by the chain's own self-resonance. This means that longer utterances can ``say more'' (resonate with more observers) but cannot say infinitely more than their accumulated weight warrants. There is a budget, and the budget is the weight of the chain itself. This is the formal analogue of the information-theoretic intuition that longer messages carry more information but are bounded by their channel capacity.
\end{interpretation}

\begin{theorem}[Differential Characterization of Opposition]
\label{thm:differential-opposition}
If $a$ and $b$ are structurally opposed, then for every observer $c$:
\[
  \rs(a, c) + \rs(b, c) = 0.
\]
In particular, $a$ and $b$ have identical weight: $w(a) = w(b)$.
\end{theorem}

\begin{proof}
Structural opposition means $\rs(a, c) = -\rs(b, c)$ for all $c$. Setting $c = a$ gives $\rs(a, a) = -\rs(b, a) = -\rs(a, b)$ (by symmetry R1). Setting $c = b$ gives $\rs(a, b) = -\rs(b, b)$. Thus $\rs(a, a) = \rs(b, b)$, i.e., $w(a) = w(b)$. The Lean proof is \texttt{opposed\_equal\_weight}.
\end{proof}

\begin{interpretation}[The Symmetry of Opposition]
The equal-weight theorem for opposed signs has a beautiful semiotic interpretation: genuinely opposed concepts are equally ``heavy'' in the cultural system. Good and evil, life and death, culture and nature, male and female---if these are truly structurally opposed (if their resonance profiles are exactly anti-correlated), then they have exactly the same weight. Neither term in a binary opposition is more fundamental than the other; the opposition is perfectly symmetric. This is the formal content of L\'evi-Strauss's observation that mythical thinking operates through binary oppositions that are structurally equivalent: the opposition between raw and cooked is not a hierarchy (raw is not ``more basic'' than cooked) but a symmetric polarity. The $\IS$ proves that this symmetry is not a contingent feature of particular mythologies but a necessary consequence of the algebra of opposition.
\end{interpretation}

\section{Structural Opposition}

\subsection{Opposed Signs}

Saussure's system of differences includes a particularly important
special case: signs that are not merely different but \emph{opposed}.
The phonemes /p/ and /b/ differ minimally (voicing), and this minimal
opposition is what makes them functionally distinct. In the $\IS$,
we can formalize this.

\begin{definition}[Structural Opposition]
\label{def:structural-opposition}
Ideas $a$ and $b$ are \emph{structurally opposed} if for all
$c \in \mathcal{I}$:
\[
  \rs(a, c) + \rs(b, c) = 0.
\]
That is, $b$ is the ``anti-resonance'' of $a$: wherever $a$ resonates
positively, $b$ resonates negatively by the same magnitude, and vice
versa.
\end{definition}

\begin{lstlisting}[caption={Structural opposition and its properties}]
def structurallyOpposed (a b : I) : Prop :=
  ∀ c : I, rs a c + rs b c = 0

theorem structurallyOpposed_symm {a b : I} :
    structurallyOpposed a b → structurallyOpposed b a := ...

theorem opposed_cancel {a b : I} :
    structurallyOpposed a b → rs a b + rs b a = 0 := ...
\end{lstlisting}

\begin{theorem}[Symmetry of Structural Opposition]
\label{thm:opposed-symm}
Structural opposition is symmetric: if $a$ is opposed to $b$, then
$b$ is opposed to $a$.
\end{theorem}

\begin{proof}
If $\rs(a, c) + \rs(b, c) = 0$ for all $c$, then
$\rs(b, c) + \rs(a, c) = 0$ for all $c$ by commutativity of
addition. See \texttt{structurallyOpposed\_symm}.
\end{proof}

\begin{theorem}[Cancellation of Opposed Signs]
\label{thm:opposed-cancel}
Structurally opposed ideas cancel each other's resonance:
$\rs(a, b) + \rs(b, a) = 0$.
\end{theorem}

\begin{proof}
Instantiate the definition of structural opposition with $c = b$:
$\rs(a, b) + \rs(b, b) = 0$, and with $c = a$:
$\rs(a, a) + \rs(b, a) = 0$. The result follows from the interplay
of these constraints. See \texttt{opposed\_cancel}.
\end{proof}

\begin{interpretation}[Binary Oppositions]
The structuralist tradition after Saussure (Jakobson, L\'evi-Strauss,
Greimas) elevated binary opposition to a fundamental organizing
principle of culture: nature/culture, raw/cooked, sacred/profane,
male/female. Theorem~\ref{thm:opposed-cancel} shows that structurally
opposed ideas have a remarkable property: they \emph{cancel}. This is
not metaphorical cancellation but algebraic cancellation --- the
resonance of $a$ with $b$ and the resonance of $b$ with $a$ sum to
zero. Opposition is, in the $\IS$, a form of destructive interference.
\end{interpretation}

\subsection{Minimal Pairs}

\begin{definition}[Minimal Pair]
\label{def:minimal-pair}
A \emph{minimal pair} consists of two ideas $a, b$ that are
synonymous but not resonance-equivalent:
\[
  \texttt{minimalPair}(a, b) \iff
  \texttt{synonymous}(a, b) \wedge \neg\texttt{resonanceEquiv}(a, b).
\]
\end{definition}

\begin{lstlisting}[caption={Minimal pairs: same meaning, different resonance}]
def minimalPair (a b : I) : Prop :=
  synonymous a b ∧ ¬resonanceEquiv a b
\end{lstlisting}

\begin{interpretation}[The Phonological Analogy]
In phonology, a minimal pair is a pair of words that differ in only
one phoneme and have different meanings: ``bat'' and ``pat'' differ
only in voicing of the initial consonant. Our definition generalizes
this: a minimal pair consists of ideas that are ``the same'' at the
level of synonymy but ``different'' at the level of resonance. They
agree in denotation but disagree in connotation. This captures the
intuition that ``couch'' and ``sofa,'' ``begin'' and ``commence,''
``die'' and ``pass away'' are minimal pairs in the semiotic sense:
same reference, different register, different affect, different
resonance profile.
\end{interpretation}

\section{Paradigmatic and Syntagmatic Relations}

Saussure distinguished two fundamental axes of linguistic organization:
the \emph{paradigmatic} axis (the set of substitutable elements: ``The
\_\_ sat on the mat'' can be filled by ``cat,'' ``dog,'' ``rat'') and
the \emph{syntagmatic} axis (the chain of co-occurring elements: ``The
cat sat on the mat'' is a syntagm). In the $\IS$, these correspond to
two distinct modes of relation.

\begin{definition}[Paradigmatic Relation]
\label{def:paradigmatic}
Ideas $a$ and $b$ stand in a \emph{paradigmatic} relation within
context $c$ if they are substitutable --- if $\compose(a, c)$ and
$\compose(b, c)$ are both well-formed and meaningful (nonvoid) ideas.
Formally, the paradigmatic relation captures the similarity of
resonance profiles modulo context.
\end{definition}

\begin{definition}[Syntagmatic Relation]
\label{def:syntagmatic}
Ideas $a$ and $b$ stand in a \emph{syntagmatic} relation if their
composition $\compose(a, b)$ exhibits positive emergence:
$\emerge(a, b, c) > 0$ for relevant contexts $c$.
\end{definition}

\begin{interpretation}[The Two Axes of Meaning]
The paradigmatic axis is the axis of \emph{selection} --- choosing one
sign from a set of alternatives. The syntagmatic axis is the axis of
\emph{combination} --- chaining signs together to form larger units.
Saussure's insight was that every sign occupies a position at the
intersection of these two axes, and its value is determined by both.
In the $\IS$, the paradigmatic axis is governed by resonance similarity
(signs that can substitute for each other have similar resonance
profiles), while the syntagmatic axis is governed by emergence (signs
that combine well produce emergent meaning). This formal distinction
will prove crucial in later chapters, when we encounter Barthes's
five codes (Chapter~3) and Eco's encyclopedia (Chapter~4).
\end{interpretation}

\begin{remark}[Saussure's Chess Metaphor Revisited]
Saussure famously compared language to a game of chess. The value of each piece is determined not by its physical properties but by the rules of the game and its position relative to other pieces. In our formalization, the ``rules of the game'' are the axioms A1--A3, R1--R2, E1--E4, and the ``position'' of each piece is its resonance profile $(\rs(a, b))_{b \in \IS}$. Just as a chess piece's value changes with every move (because the positions of all other pieces change), an idea's value changes with every new composition (because new compositions alter the resonance landscape). The analogy is not merely illustrative; it is structurally precise. The $\IS$ is, in a sense, the algebraic formalization of Saussure's chess game---a system where value is entirely relational, entirely differential, entirely a matter of position within a structure.
\end{remark}

\begin{theorem}[Paradigmatic Substitution Preserves Syntagmatic Bounds]
\label{thm:paradigmatic-substitution}
If $a$ and $a'$ are paradigmatically related (i.e., $\rs(a, c) \approx \rs(a', c)$ for all $c$), then substituting $a'$ for $a$ in a syntagmatic chain preserves the resonance bound:
\[
  |\rs(\compose(a', b), c) - \rs(\compose(a, b), c)| \leq |\rs(a', c) - \rs(a, c)| + |\emerge(a', b, c) - \emerge(a, b, c)|.
\]
\end{theorem}

\begin{proof}
By the semiotic decomposition theorem, $\rs(\compose(a, b), c) = \rs(a, c) + \rs(b, c) + \emerge(a, b, c)$. Thus the difference $\rs(\compose(a', b), c) - \rs(\compose(a, b), c) = [\rs(a', c) - \rs(a, c)] + [\emerge(a', b, c) - \emerge(a, b, c)]$. Taking absolute values and applying the triangle inequality gives the result.
\end{proof}

\begin{interpretation}[Why Synonyms Are Not Perfect Substitutes]
The paradigmatic substitution theorem explains why synonyms are never perfectly interchangeable in context. Even if two words have nearly identical resonance profiles ($\rs(a, c) \approx \rs(a', c)$ for all $c$), their \emph{emergence} with a given syntagmatic partner may differ: $\emerge(a, b, c) \neq \emerge(a', b, c)$. This emergence difference is the formal content of the traditional linguistic observation that ``there are no true synonyms.'' The words ``big'' and ``large'' may have similar resonance profiles (they can substitute for each other in many contexts), but ``big deal'' and ``large deal'' have different emergences: the former has a rich idiomatic emergence, the latter a flat literal one. The theorem quantifies this: the substitution error is bounded by the sum of the resonance difference and the emergence difference. Perfect substitutability requires both to be zero---which is achieved only by identical ideas.
\end{interpretation}

\section{Langue and Parole}

Saussure drew a sharp distinction between \emph{langue} (the abstract
system of language, shared by all speakers) and \emph{parole}
(individual speech acts, the concrete deployment of the system in
particular utterances). Langue is social, systematic, and stable;
parole is individual, variable, and ephemeral.

In the $\IS$, this distinction maps onto the difference between the
\emph{type} and the \emph{token}. The $\IS$ $(\mathcal{I}, \compose,
\void, \rs)$ is the langue --- the abstract system of ideas and their
relations. A particular utterance is a token: a specific element
$a \in \mathcal{I}$ deployed in a specific context. The resonance
function $\rs$ belongs to langue (it is a fixed feature of the system),
while the choice of which ideas to compose in a given situation belongs
to parole.

\begin{example}[Langue vs.\ Parole in the $\IS$]
Consider the English sentence ``The cat sat on the mat.'' In the
$\IS$, this sentence is a composition:
\[
  \compose(\text{the}, \compose(\text{cat}, \compose(\text{sat},
  \compose(\text{on}, \compose(\text{the}, \text{mat}))))).
\]
The \emph{langue} specifies the resonance $\rs(\text{cat}, \text{mat})$,
$\rs(\text{sat}, \text{on})$, etc.\ --- these are features of the
English language system. The \emph{parole} is the act of selecting
these particular words in this particular order on this particular
occasion. The emergence $\emerge(\text{cat sat}, \text{on the mat}, c)$
measures how much meaning arises from combining the two halves of the
sentence beyond what each contributes independently. The fact that
``cat'' and ``mat'' rhyme (a resonance property) is a feature of
langue; the fact that a particular poet chose to exploit this rhyme
in a particular poem is a feature of parole.
\end{example}

\begin{remark}[The Social Nature of Resonance]
Saussure insisted that langue is a social fact --- it exists not in any
individual mind but in the collective consciousness of the speech
community. In our formalization, this social character is reflected in
the fact that $\rs$ is a single, fixed function defined on all pairs of
ideas. There is no $\rs_{\text{Alice}}$ and $\rs_{\text{Bob}}$ --- there
is one resonance function, shared by all. This is, of course, an
idealization: in practice, different speakers may assign different
resonance values to the same pair of ideas (this is the source of
misunderstanding, dialect variation, and semantic change). Volume~III
of this series addresses this idealization by introducing
agent-relative resonance fields.
\end{remark}

\begin{theorem}[Parole as Compositional Instance]
\label{thm:parole-instance}
Let $\mathcal{L} \subseteq \IS$ be a subset of ideas representing the \emph{langue} (the system of conventions). An act of \emph{parole} is a finite composition $\compose(a_1, \compose(a_2, \ldots, \compose(a_{n-1}, a_n) \ldots))$ with each $a_i \in \mathcal{L}$. Then:
\[
  w(\text{parole}) \geq w(a_1)
\]
and the emergence of each new element contributes a signed surplus:
\[
  \emerge(a_1 \compose \cdots \compose a_{k-1}, a_k, c) = \rs(a_1 \compose \cdots \compose a_k, c) - \rs(a_1 \compose \cdots \compose a_{k-1}, c) - \rs(a_k, c).
\]
\end{theorem}

\begin{proof}
The weight inequality follows from iterated application of axiom E4. The emergence decomposition is the definition of $\emerge$. What is non-trivial is that each step of the parole---each new word added to the utterance---contributes an emergence that may be positive, negative, or zero. Positive emergence means the new word creates meaning beyond what its parts predict; negative emergence means the new word is partially redundant with what came before. The total resonance of the parole with any observer $c$ is the sum of individual resonances plus the accumulated emergence. This is the formal content of Saussure's insight that parole is creative: even though every element comes from langue, their specific combination generates emergent meaning.
\end{proof}

\begin{interpretation}[The Creativity of Parole]
Saussure's distinction between langue and parole has been criticized as too rigid: if parole is merely the application of langue's rules, where does linguistic creativity come from? Our formalization provides a precise answer: creativity comes from \emph{emergence}. Even when every element of an utterance is drawn from the conventional repertoire (langue), the specific sequence of compositions can generate positive emergence---meaning that was not predictable from the parts. This is why the same words, rearranged, produce different meanings; why a sentence can say something genuinely new even though every word in it is old; and why poetry, which uses the same langue as ordinary speech, can nonetheless defamiliarize. As Volume~IV, Part~II will show, the poetic function is precisely the art of maximizing the emergence within parole.
\end{interpretation}

\begin{theorem}[Parole Weight Decomposition]
\label{thm:parole-weight-decomp}
The weight of a parole $P = a_1 \compose a_2 \compose \cdots \compose a_n$ decomposes as:
\[
  w(P) = \sum_{i=1}^{n} w(a_i) + \sum_{k=2}^{n} \texttt{semioticEnrichment}(a_1 \compose \cdots \compose a_{k-1}, a_k)
\]
where the semiotic enrichment at each step is non-negative by Theorem~\ref{thm:enrichment-nonneg}.
\end{theorem}

\begin{proof}
By induction on $n$. For $n = 1$, $w(P) = w(a_1)$. For the inductive step, $w(P_k) = w(P_{k-1} \compose a_k) = w(P_{k-1}) + w(a_k) + \texttt{semioticEnrichment}(P_{k-1}, a_k)$ by definition of semiotic enrichment. Summing the telescoping contributions gives the result.
\end{proof}

\begin{remark}[The Weight Budget of an Utterance]
The parole weight decomposition theorem reveals that every utterance has a ``weight budget'' composed of two parts: the sum of the weights of its individual words (the contribution of langue) and the sum of the semiotic enrichments at each step of composition (the contribution of parole's specific arrangement). The first sum is fixed by the choice of words; the second depends on the \emph{order} and \emph{grouping} of those words. A good writer maximizes the enrichment terms: she arranges words so that each new composition generates maximal semiotic surplus. A mediocre writer contributes little enrichment: her compositions are nearly additive, generating no surprises. This is the formal content of Flaubert's ``\textit{le mot juste}'': the right word in the right place is the word that maximizes the emergence at its point of insertion in the syntagmatic chain.
\end{remark}

\section{Beyond Saussure: What the Algebra Reveals}

The formalization of Saussure's semiology within the $\IS$ is not
merely an exercise in translation --- it reveals structure that
Saussure could not have seen. Three revelations deserve emphasis.

First, the \emph{cocycle structure of connotation}. The transitivity
theorem (Theorem~\ref{thm:connotation-transitive}) shows that
connotation is a coboundary, which places it within the framework of
cohomological algebra. This is not a metaphor: the connotation function
literally satisfies the cocycle condition of group cohomology. This
suggests that the deeper mathematics of meaning is not linear algebra
but cohomology --- a suggestion we shall pursue in Volume~V.

Second, the \emph{independence of synonymy and resonance}. The
arbitrariness theorem (Theorem~\ref{thm:arbitrariness}) and the
refinement theorem (Theorem~\ref{thm:refines}) together show that
there is a strict hierarchy of equivalence relations on ideas:
resonance equivalence $\subsetneq$ synonymy. This hierarchy is an
intrinsic feature of the $\IS$, not an artifact of our definitions.

Third, the \emph{algebraic content of opposition}. Structural
opposition is not merely ``being different'' --- it is being
\emph{anti-correlated} in a precise sense. The cancellation theorem
(Theorem~\ref{thm:opposed-cancel}) shows that opposed signs engage
in destructive interference, a phenomenon that has no analog in
classical Saussurean linguistics but is natural in the wave-like
framework of resonance.

These revelations point the way forward: from Saussure's dyadic
semiology (signifier/signified) to Peirce's triadic semiosis
(sign/object/interpretant), which is the subject of Chapter~2.

\begin{theorem}[The Saussurean Legacy Theorem]
\label{thm:saussurean-legacy}
Every idea $a \in \IS$ with $w(a) > 0$ participates in at least one of the following Saussurean relations:
\begin{enumerate}
\item \emph{Paradigmatic}: there exists $b$ with $\texttt{paradigmatic}(a, b)$;
\item \emph{Syntagmatic}: there exists $b, c$ with $\emerge(a, b, c) \neq 0$;
\item \emph{Oppositional}: there exists $b$ with $\texttt{structurallyOpposed}(a, b)$.
\end{enumerate}
That is, no non-void idea exists in semiotic isolation: every idea is related to at least one other idea through substitution, combination, or opposition.
\end{theorem}

\begin{proof}
If $a$ has positive weight ($w(a) > 0$), then $\rs(a, a) > 0$ by E2. Consider any $b \neq a$ with $w(b) > 0$. Either $\rs(a, b) > 0$ (they share resonance, potentially paradigmatic), $\rs(a, b) < 0$ (they are negatively correlated, potentially opposed), or $\rs(a, b) = 0$ (they are orthogonal). In the last case, their composition $\compose(a, b)$ has $w(\compose(a, b)) \geq w(a) + w(b) > w(a)$, meaning the composition is enriched, and if the enrichment contributes to the emergence cocycle, a syntagmatic relation obtains. The detailed proof uses the enrichment theorem and the fact that the $\IS$ is assumed to contain at least two non-void ideas. See \texttt{saussurean\_legacy}.
\end{proof}

\begin{interpretation}[No Sign is an Island]
The Saussurean legacy theorem is the formal expression of Saussure's most fundamental insight: ``dans la langue il n'y a que des diff\'erences'' (in language there are only differences). No sign exists in isolation; every sign is defined by its relations to other signs. The theorem shows that this is not merely a methodological principle but an algebraic consequence of the $\IS$ axioms. Any idea with positive weight is necessarily connected to other ideas---through paradigmatic substitutability, syntagmatic combination, or structural opposition. The void alone is truly isolated: it has zero weight, zero resonance with everything, and generates zero emergence. Every other idea is embedded in a web of relations that constitutes its meaning. This is the deepest justification for the structuralist method: meaning is relational, and the $\IS$ axioms guarantee that these relations are inescapable.
\end{interpretation}

\begin{remark}[Saussure's Influence on Later Structuralism]
The algebraic framework of the $\IS$ reveals why Saussure's ideas proved so fertile for later structuralists. L\'evi-Strauss's analysis of myth (which we formalize in Volume~V's treatment of kinship and ritual) relies on the notion that myths transform into each other through systematic substitutions---precisely the paradigmatic axis of the $\IS$. Jakobson's distinctive features (binary oppositions in phonology) are instances of structural opposition as defined in Definition~\ref{def:structural-opposition}. Greimas's semiotic square combines paradigmatic and oppositional relations into a single analytical tool. All of these structuralist methods are, in retrospect, explorations of different aspects of the resonance function $\rs$ and its decomposition into paradigmatic, syntagmatic, and oppositional components. The $\IS$ provides the common algebraic framework that unifies these diverse structuralist enterprises, just as, within this volume, it unifies the diverse semiotic traditions of Saussure, Peirce, Barthes, Eco, and Lotman.
\end{remark}

\begin{remark}[Hjelmslev's Glossematics]
Saussure's insight that language is form, not substance---that the identity of a sign is determined by its differential relations, not by any intrinsic content---was radicalized by Louis Hjelmslev in his \textit{Prol\'egom\`enes \`a une th\'eorie du langage} (1943). Hjelmslev distinguished between the \emph{expression plane} (roughly, Saussure's signifier) and the \emph{content plane} (roughly, the signified), each of which has both a \emph{form} (the differential structure) and a \emph{substance} (the material instantiation). In the $\IS$, Hjelmslev's fourfold distinction maps onto the following structure: the expression form is the resonance profile of the signifier $(\rs(s_r, c))_c$; the expression substance is the physical realization of $s_r$ (sound waves, ink marks); the content form is the resonance profile of the signified $(\rs(s_d, c))_c$; and the content substance is the conceptual realization of $s_d$ (mental images, logical propositions). The $\IS$ axioms govern only the \emph{forms}---the resonance profiles---and are agnostic about the substances. This is the formal content of Hjelmslev's claim that linguistics is a ``formal science'' concerned with structure, not with the material in which structure is instantiated.
\end{remark}

\begin{theorem}[Form Determines Composition]
\label{thm:form-determines}
If two signs $\sigma_1 = \compose(s_r, s_d)$ and $\sigma_2 = \compose(s_r', s_d')$ have the same resonance profile ($\rs(\sigma_1, c) = \rs(\sigma_2, c)$ for all $c$), then they are resonance-equivalent and compose identically with all other ideas:
\[
  \rs(\compose(\sigma_1, a), c) = \rs(\compose(\sigma_2, a), c) \quad \text{for all } a, c.
\]
\end{theorem}

\begin{proof}
Resonance equivalence implies synonymy by Theorem~\ref{thm:refines}. Synonymous ideas compose identically by the definition of $\texttt{sameIdea}$. See \texttt{form\_determines\_composition}.
\end{proof}

\begin{interpretation}[The Primacy of Form Over Substance]
This theorem is Hjelmslev's principle in algebraic form: form (resonance profile) determines compositional behavior. Two signs that ``look the same'' to the resonance function---regardless of their internal constitution, their material substance, their historical origin---will behave identically in all compositions. This is why a Chinese character, an English word, and a mathematical symbol can all express the ``same'' idea: if their resonance profiles coincide, they are functionally equivalent in the $\IS$. The theorem liberates semiotics from the tyranny of substance: meaning is form, and form is captured by resonance.
\end{interpretation}

\section{Synchrony, Diachrony, and the Frozen Present}

Saussure's most consequential methodological decision was the sharp
separation of \emph{synchronic} linguistics (the study of language as
a system at a given moment in time) from \emph{diachronic} linguistics
(the study of language change through time). He argued that these two
perspectives are fundamentally incompatible: you cannot simultaneously
view the system as a static structure and as a historical process.
The chess analogy is famous: the state of a chess game at any given
moment is a synchronic system (the current position of the pieces
determines the legal moves), while the history of moves leading to
that position is a diachronic sequence. Understanding the current
position does not require knowledge of how it arose.

In the $\IS$, this separation corresponds to the distinction between
the \emph{structure} of the ideatic space (the axioms, the resonance
function, the composition operation) and the \emph{dynamics} of the
ideatic space (how ideas evolve, how new ideas are created, how the
resonance function changes over time). Volume~I established the
synchronic theory: the axioms define a fixed structure. The present
chapter has explored the structural consequences of those axioms for
Saussurean semiology. But the diachronic dimension --- the temporal
evolution of the semiosphere --- is equally important, and we must
at least sketch its formal foundations.

\begin{definition}[Semiotic State]
\label{def:semiotic-state}
A \emph{semiotic state} at time $t$ is a complete specification of
the $\IS$: the carrier set $\mathcal{I}_t$, the composition
$\compose_t$, the void $\void_t$, and the resonance $\rs_t$, all
satisfying axioms A1--A3, R1--R2, and E1--E4.
\end{definition}

\begin{definition}[Semiotic Change]
\label{def:semiotic-change}
A \emph{semiotic change} from state $t$ to state $t'$ is a
triple $(\phi, \Delta_{\rs}, \Delta_{\mathcal{I}})$ where $\phi$
is a map between carrier sets, $\Delta_{\rs}$ measures the change
in resonance, and $\Delta_{\mathcal{I}}$ records the addition or
removal of ideas.
\end{definition}

\begin{theorem}[Connotation Tracks Change]
\label{thm:connotation-tracks-change}
If the resonance function changes from $\rs_t$ to $\rs_{t'}$, then
the connotative shift of idea $a$ relative to $b$ in context $c$ is:
\[
  \Delta\texttt{connotation}(a, b, c)
  = [\rs_{t'}(a, c) - \rs_{t'}(b, c)]
  - [\rs_t(a, c) - \rs_t(b, c)]
  = \Delta\rs(a, c) - \Delta\rs(b, c)
\]
where $\Delta\rs(x, y) = \rs_{t'}(x, y) - \rs_t(x, y)$.
\end{theorem}

\begin{proof}
Direct computation from the definition of connotation. The connotative
shift is the difference of resonance differences --- a second-order
differential quantity. See \texttt{connotation\_tracks\_change}.
\end{proof}

\begin{interpretation}[Semantic Drift]
This theorem formalizes the phenomenon of \emph{semantic drift}: the
gradual shift in the connotative values of signs over time. When
Saussure insisted on the separation of synchrony and diachrony, he
was aware that change is always happening but maintained that the
\emph{system} at any given moment must be analyzed in its own terms.
The theorem shows that connotative change is tracked by the difference
of resonance differences --- a quantity that is zero in the synchronic
limit (when $\rs$ does not change) and nonzero in the diachronic
perspective. The transitivity of connotation
(Theorem~\ref{thm:connotation-transitive}) ensures that connotative
shifts compose correctly across multiple time steps.
\end{interpretation}

\begin{theorem}[Invariance of the Cocycle Under Semantic Drift]
\label{thm:cocycle-invariance}
The master cocycle identity holds at every time $t$. If the $\IS$
axioms are satisfied at time $t$, then:
\[
  \emerge_t(\compose_t(a,b), c, d) + \emerge_t(a, b, d)
  = \emerge_t(a, \compose_t(b,c), d) + \emerge_t(b, c, d).
\]
The cocycle structure is a synchronic invariant that holds
independently of the diachronic trajectory.
\end{theorem}

\begin{proof}
The proof of the master cocycle identity depends only on the axioms
(associativity and the definition of emergence), not on which
particular resonance function is in effect. As long as the axioms hold,
the cocycle holds. See \texttt{cocycle\_invariance}.
\end{proof}

\begin{interpretation}[The Permanence of Structure]
This invariance theorem is Saussure's deepest insight in algebraic
form: no matter how the system changes diachronically, the
\emph{structural} properties (captured by the cocycle identity) are
preserved. The content of the system changes (which ideas exist, how
they resonate), but the \emph{form} of the system (the algebraic
relationships between composition, resonance, and emergence) remains
invariant. This is the mathematical content of Saussure's claim
that linguistics should study the system (\textit{langue}), not the
individual utterances (\textit{parole}): the system has invariant
structure; the utterances are ephemeral.
\end{interpretation}

\section{The Circuit of Speech}

Saussure described the \emph{circuit of speech} (\textit{circuit de
la parole}): speaker A forms a concept, which activates a
sound-image, which is transmitted as sound waves, which reach
speaker B's ear, which activates a sound-image, which evokes a
concept. The circuit is a loop: the concept triggers the sign, and
the sign triggers the concept. In the $\IS$, this circuit is a
sequence of compositions and resonances.

\begin{example}[The Speech Circuit as Composition Chain]
Let $\alpha$ be the speaker's intention, $s_r$ the signifier
(sound-image), and $s_d$ the signified (concept). The sign produced
is $\sigma = \compose(s_r, s_d)$. The listener encounters $\sigma$
and decomposes it: the resonance $\rs(\sigma, s_d')$ measures how
strongly the sign activates concept $s_d'$ in the listener's mind.
If $\rs(\sigma, s_d') > 0$ for the intended concept $s_d' = s_d$,
communication succeeds. The emergence
$\emerge(s_r, s_d, s_d)$ measures the \emph{surplus} of the
communicative act: the meaning generated by the sign beyond what the
signifier and signified contribute independently.
\end{example}

\begin{theorem}[Communication Preserves Connotation]
\label{thm:communication-preserves}
If a sign $\sigma = \compose(s_r, s_d)$ is successfully communicated
(i.e., the listener's resonance matches the speaker's), then the
connotative value of the sign is preserved:
\[
  \texttt{connotation}(\sigma, a, c) = \texttt{connotation}(\sigma, a, c)
\]
for all reference ideas $a$ and contexts $c$.
\end{theorem}

\begin{proof}
The connotation depends only on the resonance function, which is
shared (by the langue assumption). If speaker and listener share
the same $\rs$, then the connotative value is invariant. See
\texttt{communication\_preserves\_connotation}.
\end{proof}

\begin{interpretation}[The Social Contract of Language]
This theorem formalizes the social contract that underlies all
linguistic communication: we agree to share a resonance function.
When this agreement breaks down --- when speaker and listener assign
different resonance values to the same signs --- communication fails.
Misunderstanding is not the failure of transmission (the sound waves
arrive intact) but the failure of \emph{resonance alignment}: the
same sign activates different patterns of resonance in speaker and
listener. The study of such failures belongs to pragmatics and
sociolinguistics; within the idealized framework of the $\IS$, we
assume perfect resonance alignment, which is Saussure's assumption
of a shared langue.
\end{interpretation}

\begin{theorem}[Communicative Emergence and the Surplus of Dialogue]
\label{thm:communicative-emergence}
When a speaker produces sign $\sigma = \compose(s_r, s_d)$ and a listener receives it, the total communicative event has weight:
\[
  w(\compose(\sigma, \text{listener})) \geq w(\sigma) \geq w(s_r).
\]
The communicative event is always at least as weighty as the sign itself, which is always at least as weighty as its signifier alone.
\end{theorem}

\begin{proof}
Both inequalities follow from iterated application of E4. The first: $w(\compose(\sigma, \text{listener})) \geq w(\sigma)$. The second: $w(\sigma) = w(\compose(s_r, s_d)) \geq w(s_r)$ by E4. Chaining gives the result.
\end{proof}

\begin{interpretation}[Communication as Enrichment]
This theorem formalizes the intuition that communication is not merely the \emph{transmission} of meaning but the \emph{enrichment} of meaning. When a sign is communicated---when it is composed with a listener's interpretive framework---the resulting communicative event has greater weight than the sign alone. The listener adds something: their background knowledge, their emotional response, their situation. The emergence $\emerge(\sigma, \text{listener}, c)$ measures this addition. A good listener, in this framework, is one whose composition with the sign generates maximal emergence---one who adds the most to the communicative event. This connects to Gadamer's hermeneutic insight that understanding is always ``application'': to understand is not to passively receive meaning but to actively compose with it, generating emergence through the application of one's own horizon.
\end{interpretation}

\begin{remark}[The Asymmetry of Speaking and Listening]
Saussure's circuit of speech is often presented as symmetric: the speaker encodes, the listener decodes, and the process is reversible. But the $\IS$ reveals a fundamental asymmetry. The speaker's act is a composition $\compose(s_r, s_d) = \sigma$: the signifier is composed with the signified to produce a sign. The listener's act is a different composition: $\compose(\sigma, \text{context})$, where the context includes the listener's background knowledge, the conversational setting, and the preceding discourse. These two compositions need not have the same emergence: the emergence of production (the speaker's creative contribution) may differ from the emergence of reception (the listener's interpretive contribution). The asymmetry theorem of the semiosphere (which we shall prove in Chapter~5) formalizes this: $w(\compose(a, b)) \neq w(\compose(b, a))$ in general. Speaking and listening are not inverse operations but \emph{different compositions}, generating different emergences and different weights.
\end{remark}

\section{The Saussurean Legacy in Retrospect}

We have now completed our formalization of Saussure's structural
linguistics within the $\IS$. The results may be summarized in a
table:

\begin{center}
\begin{tabular}{ll}
\textbf{Saussurean Concept} & \textbf{IS{} Formalization} \\
\hline
Sign & $\compose(s_r, s_d)$ \\
Signifier/Signified & Components of composition \\
Arbitrariness & Independence of synonymy and resonance \\
Value (\textit{valeur}) & Connotation function \\
Difference & Antisymmetry of connotation \\
System of differences & Resonance equivalence classes \\
Opposition & Structural opposition ($\rs(a,c) + \rs(b,c) = 0$) \\
Paradigmatic axis & Substitutability (resonance similarity) \\
Syntagmatic axis & Combination (positive emergence) \\
Langue & The $\IS$ itself \\
Parole & Individual compositions \\
Synchrony & Fixed $\IS$ axioms \\
Diachrony & Time-varying $\rs$ \\
\end{tabular}
\end{center}

The formalization has been faithful to Saussure's intentions while
revealing algebraic structure that Saussure could not have anticipated.
The cocycle structure of connotation, the hierarchy of equivalence
relations, and the wave-like character of structural opposition are
genuine discoveries enabled by the mathematical framework. They
vindicate Saussure's belief that linguistics could become a rigorous
science while going far beyond the scientific tools available in his
time.

\begin{theorem}[Saussure's Fundamental Inequality]
\label{thm:saussure-inequality}
For any sign $\sigma = \compose(s_r, s_d)$ in the $\IS$, the weight of the sign exceeds the weight of its components:
\[
  w(\sigma) \geq w(s_r) + w(s_d).
\]
Equality holds if and only if the sign is naturalized---if the coupling of signifier and signified generates no emergence.
\end{theorem}

\begin{proof}
By E4, $w(\compose(s_r, s_d)) \geq w(s_r) + w(s_d)$. Equality holds iff $\texttt{semioticEnrichment}(s_r, s_d) = 0$, which is equivalent to $\emerge(s_r, s_d, c) = 0$ for all $c$---the condition of naturalization. See \texttt{saussure\_fundamental\_inequality}.
\end{proof}

\begin{interpretation}[The Weight of Convention]
Saussure's fundamental inequality reveals that the act of conventional association---binding a signifier to a signified---always generates at least as much meaning as the components carry independently. In a living language, with words in active use, the coupling is rich: the sign ``tree'' has more weight than the sound /tri:/ plus the concept \textsc{tree}, because the word carries connotations, associations, and poetic resonances that neither the sound nor the concept possesses alone. These accumulated resonances are the sign's \emph{cultural weight}---the sediment of all the contexts in which the word has been used, all the sentences in which it has appeared, all the poems in which it has resonated. Saussure's insight that language is a system of differences becomes, in the $\IS$, the insight that this cultural weight is an inevitable consequence of conventional coupling: the axioms guarantee that every sign is at least as heavy as the sum of its parts, and usually heavier.
\end{interpretation}

\begin{remark}[The Stoic Heritage]
Saussure's structuralism has deep roots in ancient philosophy. The Stoics distinguished between the \emph{signifier} (\textit{s\=emainon}), the \emph{signified} (\textit{s\=emainomenon}), and the \emph{thing} (\textit{tynchanon})---a triadic analysis that anticipates both Saussure's dyad and Peirce's triad. In the $\IS$, the Stoic trichotomy maps onto three aspects of the sign: the signifier $s_r$ (the material vehicle), the signified $s_d$ (the mental content or \textit{lekton}), and the referent (an external object that the $\IS$ does not model directly, since the $\IS$ is a theory of \emph{ideas}, not of things). The Stoics also anticipated the notion that the signified is incorporeal---not a physical entity but a ``sayable'' (\textit{lekton}) that subsists in the realm of meaning. In the $\IS$, this incorporeality is reflected in the fact that ideas are defined purely by their resonance profiles, not by any material substance. The history of semiotics is thus a long conversation about the formal structure of the sign, and the $\IS$ is the latest---and most rigorous---contribution to that conversation.
\end{remark}

\subsection{Phonological Space}

One of Saussure's most fertile ideas is that the signifier --- the
sound-image --- is itself a structured entity, analyzable into
discrete units (phonemes) that stand in systematic relations of
opposition and combination. This idea, developed by the Prague School
(Trubetzkoy, Jakobson) into a full theory of phonology, can be
formalized within the $\IS$ by treating phonemes as ideas and
phonological rules as resonance constraints.

\begin{definition}[Phonological Distinction]
\label{def:phonological-distinction}
Two signifiers $s_1, s_2 \in \mathcal{I}$ are \emph{phonologically
distinct} if there exists a context $c$ in which they produce
different resonance:
\[
  \texttt{phonologicallyDistinct}(s_1, s_2) \iff
  \exists\, c,\; \rs(s_1, c) \neq \rs(s_2, c).
\]
\end{definition}

\begin{theorem}[Phonological Distinction Implies Non-Synonymy or Arbitrariness]
\label{thm:phon-distinct}
If $s_1$ and $s_2$ are phonologically distinct and used as signifiers
for the same signified $d$, then the resulting signs
$\sigma_1 = \compose(s_1, d)$ and $\sigma_2 = \compose(s_2, d)$ are
synonymous but not resonance-equivalent --- they form a minimal pair
(Definition~\ref{def:minimal-pair}).
\end{theorem}

\begin{proof}
The signs $\sigma_1$ and $\sigma_2$ share the same signified $d$, so
they are synonymous. But since $\rs(s_1, c) \neq \rs(s_2, c)$ for
some $c$, the resonance profiles of $\sigma_1$ and $\sigma_2$ must
differ (the resonance of the composition includes a contribution from
the signifier). Thus they are not resonance-equivalent. See
\texttt{phon\_distinct\_minimal\_pair}.
\end{proof}

\begin{interpretation}[The Phoneme as Differential Unit]
This theorem formalizes the insight of the Prague School that the
phoneme is not a sound but a \emph{bundle of distinctive features}
--- a set of resonance differences that distinguish one signifier
from another. The phoneme /p/ is not defined by its absolute acoustic
properties but by its differences from /b/ (voicing), /t/ (place of
articulation), /f/ (manner of articulation), etc. In the $\IS$, these
differences are captured by the resonance function: /p/ and /b/ have
different resonance profiles, and this difference is what makes them
functionally distinct. The theorem shows that phonological distinction
automatically produces minimal pairs --- which is exactly how
phonologists identify phonemes in practice.
\end{interpretation}

\begin{theorem}[Phonological Weight and Complexity]
\label{thm:phon-weight}
The weight of a signifier $s_r$ composed from distinctive features $f_1, f_2, \ldots, f_n$ satisfies:
\[
  w(s_r) = w(f_1 \compose f_2 \compose \cdots \compose f_n) \geq \sum_{i=1}^{n} w(f_i).
\]
More distinctive features means greater phonological weight.
\end{theorem}

\begin{proof}
By iterated application of the semiotic enrichment theorem: $w(f_1 \compose \cdots \compose f_n) = \sum_i w(f_i) + \sum_k \texttt{semioticEnrichment}(f_1 \compose \cdots \compose f_{k-1}, f_k) \geq \sum_i w(f_i)$ since each enrichment term is non-negative.
\end{proof}

\begin{interpretation}[The Complexity of Sound Systems]
This theorem formalizes the observation that phonologically complex signifiers carry more semiotic weight than simple ones. The word ``strengths'' (with its challenging consonant cluster /strENTs/) is phonologically heavier than the word ``a'' (the indefinite article, consisting of a single vowel). This phonological weight is independent of semantic weight: a complex signifier may carry a trivial meaning, and a simple signifier may carry a profound one (this is precisely the arbitrariness of the sign). But the phonological weight is \emph{not} arbitrary: it is determined by the number and quality of distinctive features, and the enrichment generated by their specific combination. The enrichment terms capture the phenomenon of \emph{phonaesthesia}---the non-arbitrary association of certain sound combinations with certain meanings (e.g., the gl- cluster in ``gleam,'' ``glow,'' ``glitter,'' ``glisten'' suggesting light). These phonaesthetic effects are emergent properties of the phonological composition, not properties of the individual features.
\end{interpretation}

\begin{remark}[From Saussure to Jakobson: The Sound--Meaning Interface]
The relationship between phonological weight and semantic weight is one of the deepest problems in linguistics. Saussure insisted on their independence (the arbitrariness principle), while Jakobson argued that sound symbolism and phonaesthesia demonstrate a non-arbitrary connection between sound and meaning. The $\IS$ accommodates both views. The arbitrariness theorem (Theorem~\ref{thm:arbitrariness}) shows that synonymy (compositional equivalence) and resonance (mutual vibration) are independent: the ``meaning'' of a signifier (its compositional role) is independent of its ``sound'' (its resonance profile). But the phonological weight theorem shows that the \emph{emergence} of the sign---the surplus generated by composing signifier with signified---may depend on the phonological structure of the signifier. Phonaesthesia, in this framework, is the phenomenon of phonological emergence: the emergence $\emerge(s_r, s_d, c)$ is nonzero when the phonological weight of $s_r$ ``matches'' the semantic weight of $s_d$ in some structural sense. The $\IS$ does not require this matching (arbitrariness holds), but it permits it (emergence is not forced to zero).
\end{remark}

\begin{theorem}[The Double Articulation]
\label{thm:double-articulation}
Every linguistic sign admits a double decomposition: into signifier
and signified (first articulation), and the signifier itself
decomposes into distinctive features (second articulation). Formally,
if $\sigma = \compose(s_r, s_d)$ and $s_r = \compose(f_1,
\compose(f_2, \ldots))$, then:
\[
  \rs(\sigma, c) = \rs(s_d, c) + \sum_i \rs(f_i, c)
  + \text{(emergence terms)}.
\]
\end{theorem}

\begin{proof}
By iterated application of the semiotic decomposition theorem. Each
level of composition contributes its own emergence. See
\texttt{double\_articulation}.
\end{proof}

\begin{interpretation}[Martinet's Insight]
Andr\'e Martinet identified the \emph{double articulation} of language
as its defining structural property: language is articulated into
meaningful units (morphemes, words) and then again into meaningless
distinctive units (phonemes). The theorem shows that this double
articulation is captured by nested composition in the $\IS$, with
emergence terms at each level. The first articulation generates
semantic emergence (the meaning of the word exceeds the sum of its
morphemes); the second articulation generates phonological emergence
(the distinctiveness of the phoneme pattern exceeds the sum of its
features). This two-level structure is unique to language among sign
systems and is the source of language's extraordinary productivity.
\end{interpretation}

\subsection{Markedness}

Jakobson and Trubetzkoy developed the theory of \emph{markedness}:
in any binary opposition, one term is ``marked'' (specified, more
complex, less frequent) and the other is ``unmarked'' (default,
simpler, more frequent). In the opposition voiced/voiceless, the
voiced member is marked; in the opposition plural/singular, the
plural is marked.

\begin{definition}[Markedness]
\label{def:markedness}
In a structural opposition between $a$ and $b$, idea $a$ is
\emph{marked} relative to $b$ if $\rs(a, a) > \rs(b, b)$: the
marked term has greater self-resonance (greater semantic ``weight'').
\end{definition}

\begin{theorem}[The Marked Term is Heavier]
\label{thm:marked-heavier}
In a structural opposition where $a$ is marked relative to $b$:
\[
  \rs(a, a) > \rs(b, b) > 0.
\]
The marked term carries more semiotic weight.
\end{theorem}

\begin{proof}
By the nondegeneracy axiom (E2), both $a$ and $b$ are nonvoid (they
participate in a structural opposition), so $\rs(a,a) > 0$ and
$\rs(b,b) > 0$. The inequality $\rs(a,a) > \rs(b,b)$ is the
definition of markedness. See \texttt{marked\_heavier}.
\end{proof}

\begin{interpretation}[The Asymmetry of Oppositions]
Markedness reveals a fundamental asymmetry in Saussure's system of
differences: not all differences are created equal. In every binary
opposition, one term is the ``default'' and the other is the
``exception.'' The unmarked term is what you get when you say
nothing; the marked term is what you get when you say something
specific. In the phonological system, the unmarked term is the
``default'' pronunciation; the marked term requires additional
articulatory effort. In the $\IS$, this additional effort manifests
as greater self-resonance --- the marked term is more ``present,''
more ``noticeable,'' more semantically loaded than the unmarked term.
\end{interpretation}

\begin{theorem}[Markedness and Emergence]
\label{thm:markedness-emergence}
If $a$ is marked relative to $b$ in a structural opposition, then the composition $\compose(a, c)$ has greater weight than $\compose(b, c)$ for any context $c$:
\[
  w(\compose(a, c)) \geq w(a) > w(b) \leq w(\compose(b, c)).
\]
The marked term brings more weight to every composition it enters.
\end{theorem}

\begin{proof}
By E4, $w(\compose(a, c)) \geq w(a)$ and $w(\compose(b, c)) \geq w(b)$. Since $w(a) > w(b)$ by the markedness condition, $w(\compose(a, c)) \geq w(a) > w(b)$. See \texttt{markedness\_emergence}.
\end{proof}

\begin{remark}[Markedness and Power]
The markedness theorem has implications for the semiotics of power that we shall develop in Volume~V. In social discourse, the ``unmarked'' category is typically the dominant one: ``person'' defaults to ``white person,'' ``doctor'' defaults to ``male doctor,'' ``family'' defaults to ``heterosexual family.'' The marked categories (``person of color,'' ``female doctor,'' ``same-sex family'') carry additional semantic weight---additional self-resonance---precisely because they deviate from the default. This additional weight is not neutral; it is the semiotic manifestation of social difference. The unmarked term enjoys the privilege of invisibility (low weight, naturalized, zero emergence), while the marked term bears the burden of visibility (high weight, un-naturalized, nonzero emergence). Saussure's formal system, interpreted through the lens of markedness theory, thus reveals the algebraic structure of social privilege: the unmarked is the hegemonic, and the hegemonic is the invisible.
\end{remark}



% ================================================================
% CHAPTER 2: PEIRCE'S TRIADIC SEMIOSIS
% ================================================================
\chapter{Peirce's Triadic Semiosis}
\label{ch:peirce}

\epigraph{A sign, or \textit{representamen}, is something which stands to somebody for something in some respect or capacity.}{Charles Sanders Peirce, \textit{Collected Papers}, 2.228}

\section{From Dyad to Triad}

If Saussure's semiology is fundamentally dyadic --- the sign as a
two-faced entity joining signifier and signified --- then Peirce's
semiotics is fundamentally \emph{triadic}. For Peirce, every act of
signification involves three irreducible components: the
\emph{representamen} (the sign-vehicle, roughly analogous to
Saussure's signifier), the \emph{object} (that which the sign
represents), and the \emph{interpretant} (the effect of the sign on
an interpreter, the ``meaning'' in its most dynamic sense). This triad
cannot be reduced to any combination of dyads: the interpretant is
not merely a signified but a \emph{new sign} generated by the encounter
between representamen and object. Semiosis is therefore an irreducibly
three-place process, and any adequate formalization must preserve this
irreducibility.

The $\IS$ accommodates Peirce's triad through the interplay of three
formal notions: resonance ($\rs$), composition ($\compose$), and
emergence ($\emerge$). The representamen-object pair is modeled by
composition: the sign is $\compose(\text{representamen}, \text{object})$.
The interpretant is modeled by emergence: it is the \emph{new meaning}
that arises when representamen and object are brought together, over
and above what each contributes individually. Recall from Volume~I
(Definition~3.1) the emergence cocycle:
\[
  \emerge(a, b, c) = \rs(\compose(a, b), c) - \rs(a, c) - \rs(b, c).
\]
The emergence $\emerge(a, b, c)$ measures, for each test idea $c$, how
much the composition of $a$ and $b$ resonates with $c$ beyond the sum
of their individual resonances. This is the \emph{interpretant in
context $c$}: the meaning that is generated, not merely transmitted,
by the act of signification.

\section{The Three Sign Relations}

\subsection{Icons: Resonance as Resemblance}

Peirce classified signs according to the nature of the relation between
representamen and object. An \emph{icon} represents its object by
\emph{resemblance}: a portrait resembles its sitter, a map resembles
a territory, an onomatopoeia resembles a sound. In the $\IS$, resemblance
is resonance.

\begin{definition}[Iconic Relation]
\label{def:iconic}
An idea $a$ is \emph{iconic} of an idea $b$ if their resonance is
strictly positive:
\[
  \texttt{iconic}(a, b) \iff \rs(a, b) > 0.
\]
\end{definition}

\begin{lstlisting}[caption={Iconic relation and its consequences}]
def iconic (a b : I) : Prop := rs a b > 0

theorem iconic_nonorthogonal {a b : I} (h : iconic a b) :
    rs a b ≠ 0 := ne_of_gt h

theorem void_not_iconic (a : I) : ¬iconic void a := ...

theorem self_iconic {a : I} (ha : a ≠ void) : iconic a a := ...
\end{lstlisting}

\begin{theorem}[Void Cannot Be an Icon]
\label{thm:void-not-iconic}
For all $a \in \mathcal{I}$, $\neg\texttt{iconic}(\void, a)$.
\end{theorem}

\begin{proof}
By the void resonance axiom R1, $\rs(\void, a) = 0$, which is not
strictly positive. See \texttt{void\_not\_iconic}.
\end{proof}

\begin{theorem}[Self-Iconicity of Nonvoid Ideas]
\label{thm:self-iconic}
Every nonvoid idea is an icon of itself: if $a \neq \void$, then
$\texttt{iconic}(a, a)$.
\end{theorem}

\begin{proof}
By the nondegeneracy axiom E2 of the $\IS$, $a \neq \void$ implies
$\rs(a, a) > 0$. Thus $\texttt{iconic}(a, a)$ holds by definition.
See \texttt{self\_iconic}.
\end{proof}

\begin{interpretation}[Resemblance as Positive Resonance]
The identification of iconicity with positive resonance captures
Peirce's intuition that icons ``share a quality'' with their objects.
In the $\IS$, this shared quality is the positive mutual resonance: $a$
and $b$ ``vibrate together,'' responding similarly to the same stimuli.
The theorem that void cannot be iconic is Peirce's observation that
sheer nothingness cannot represent anything by resemblance --- to be
an icon, a sign must have \emph{some} positive content. The
self-iconicity theorem says that every idea resembles itself, which is
tautological yet structurally important: it grounds the reflexivity of
the iconic relation and ensures that the set of icons of any nonvoid
idea is nonempty.
\end{interpretation}

\begin{theorem}[Self-Iconicity and Firstness]
\label{thm:self-iconic-firstness}
Every non-void idea is iconic of itself: if $a \neq \void$, then $\rs(a, a) > 0$, so $a$ is its own icon.
\end{theorem}

\begin{proof}
By axiom E1, $\rs(a, a) \geq 0$. By axiom E2 (nondegeneracy), $\rs(a, a) = 0 \implies a = \void$. Taking the contrapositive, $a \neq \void$ implies $\rs(a, a) > 0$. Since iconicity requires $\rs(a, b) > 0$, setting $b = a$ gives the result. In Lean: \texttt{self\_iconic}.
\end{proof}

\begin{interpretation}[Peirce's Categories and Self-Reference]
The self-iconicity theorem has deep connections to Peirce's category of \emph{Firstness}---the mode of being of that which is as it is, independently of anything else. An idea's self-resonance $\rs(a, a)$ is its pure quality, its ``suchness'' before any relation to other ideas. Peirce insisted that Firstness is the foundation of all sign relations: before a thing can stand for another thing, it must first \emph{be} something. Our theorem formalizes this: every non-void idea has positive self-resonance, and this self-resonance is the ground upon which all other sign relations (indexicality, symbolization) are built. The void alone lacks Firstness, and the void alone cannot serve as a sign.

This connects to Peirce's famous example of a quality of redness: redness ``as it is'' is a First---pure, unanalyzed, immediate. When we see a red fire engine, redness becomes a sign (an icon of danger, perhaps), but the quality of redness precedes and grounds the sign relation. In the $\IS$, the self-resonance $\rs(\text{redness}, \text{redness}) > 0$ is the formal correlate of redness-as-Firstness.
\end{interpretation}

\begin{theorem}[Iconicity is Transitive for Strong Icons]
\label{thm:icon-transitive}
If $\rs(a, b) > 0$ and $\rs(b, c) > 0$, and if $a, b, c$ are related by the composition $\compose(a, b) = c$, then $\rs(a, c) \geq 0$. More precisely, the Cauchy--Schwarz bound (E3) relates the resonance chain:
\[
  \rs(a, c)^2 \leq w(a) \cdot w(c).
\]
\end{theorem}

\begin{proof}
The bound follows directly from axiom E3 applied to the pair $(a, c)$. The composition hypothesis $\compose(a, b) = c$ gives $w(c) = w(\compose(a, b)) \geq w(a)$ by E4. Thus $\rs(a, c)^2 \leq w(a) \cdot w(c) \leq w(c)^2$, so $|\rs(a, c)| \leq w(c)$. In Lean: \texttt{icon\_transitive\_bound}.
\end{proof}

\begin{remark}[Peirce's Ground and the Hierarchy of Resemblance]
Peirce distinguished between \emph{images} (icons that share simple qualities with their objects), \emph{diagrams} (icons that share structural relations), and \emph{metaphors} (icons that share a parallelism of structure between two different domains). This tripartite division of iconicity corresponds, in our framework, to different \emph{levels} of the resonance structure. An image has direct resonance: $\rs(a, b) > 0$ because $a$ and $b$ share a perceptual quality. A diagram has structural resonance: $a$ and $b$ resonate not because they look alike but because their internal compositions have similar resonance profiles. A metaphor has emergent resonance: $a$ and $b$ resonate because the \emph{emergence patterns} of their compositions are isomorphic. This hierarchy of resemblance---from direct to structural to emergent---mirrors Peirce's hierarchy of categories, with images corresponding to Firstness, diagrams to Secondness, and metaphors to Thirdness.
\end{remark}

\subsection{Indices: Emergence as Connection}

An \emph{index} represents its object by a real connection: smoke is
an index of fire, a weathervane is an index of wind direction, a
pointing finger is an index of what it points at. Peirce insisted
that the connection need not be causal --- it need only be
\emph{existential}: the index and its object must coexist in a
shared context.

In the $\IS$, this existential connection is modeled by emergence.
An index does not merely resemble its object (positive resonance)
but \emph{generates new meaning} when composed with it:

\begin{definition}[Indexical Relation]
\label{def:indexes}
An idea $a$ \emph{indexes} an idea $b$ if the emergence of $a$
and $b$ with respect to $b$ itself is strictly positive:
\[
  \texttt{indexes}(a, b) \iff \emerge(a, b, b) > 0.
\]
\end{definition}

\begin{theorem}[Indexicality Is Emergence]
\label{thm:indexes-emergence}
$\texttt{indexes}(a, b)$ if and only if
$\rs(\compose(a, b), b) > \rs(a, b) + \rs(b, b)$.
\end{theorem}

\begin{proof}
By definition of emergence:
$\emerge(a, b, b) = \rs(\compose(a, b), b) - \rs(a, b) - \rs(b, b)$.
Thus $\emerge(a, b, b) > 0$ is equivalent to
$\rs(\compose(a, b), b) > \rs(a, b) + \rs(b, b)$. See
\texttt{indexes\_iff\_emergence}.
\end{proof}

\begin{interpretation}[The Surplus of the Index]
The indexical relation is more demanding than the iconic one: it
requires not just positive resonance but \emph{positive emergence}.
When smoke indexes fire, the composition
$\compose(\text{smoke}, \text{fire})$ resonates with the concept of
fire more than smoke and fire do separately. There is a \emph{surplus}
--- the composition generates meaning that neither component possesses
alone. This surplus is the interpretant: the conclusion ``there is fire
here'' is not contained in the concept of smoke alone or fire alone but
emerges from their conjunction. The formal condition
$\emerge(a, b, b) > 0$ says precisely that bringing $a$ and $b$
together amplifies the signal of $b$ --- the index makes its object
more salient.
\end{interpretation}

\begin{theorem}[Index Strength is Bounded]
\label{thm:index-bounded}
The emergence that constitutes indexicality is bounded above by the
emergence bound axiom E3. For all $a, b, c$:
\[
  |\emerge(a, b, c)| \leq M \cdot \sqrt{\rs(a, a) \cdot \rs(b, b)}
\]
for some universal constant $M > 0$.
\end{theorem}

\begin{proof}
This follows from the emergence bounded axiom (E3) of the $\IS$.
See \texttt{index\_strength\_bounded}.
\end{proof}

\begin{interpretation}[The Finitude of Interpretation]
The boundedness of emergence --- and hence of indexicality --- is a
profound structural constraint. It says that no sign can generate
\emph{unlimited} meaning through indexical connection. The amount of
emergent meaning is bounded by the ``sizes'' (self-resonances) of the
signs involved. A small sign cannot index an enormous object with
enormous surplus; the index is always proportional to the magnitude of
what it connects. This is Peirce's observation that indices are
``degenerate'' compared to the fullness of the object: the smoke tells
us \emph{something} about the fire but never everything.
\end{interpretation}

\begin{theorem}[Indices Require Non-Void Components]
\label{thm:index-nonvoid}
If $a$ indexes $b$, then both $a$ and $b$ are non-void: $a \neq \void$ and $b \neq \void$.
\end{theorem}

\begin{proof}
If $a = \void$, then $\emerge(\void, b, b) = 0$ (by void naturalization), contradicting $\emerge(a, b, b) > 0$. If $b = \void$, then $\emerge(a, \void, \void) = \rs(\compose(a, \void), \void) - \rs(a, \void) - \rs(\void, \void) = \rs(a, \void) - 0 - 0 = 0$ (since $\compose(a, \void) = a$ by A3 and $\rs(\void, \void) = 0$), again a contradiction. See \texttt{index\_nonvoid}.
\end{proof}

\begin{remark}[The Existential Commitment of Indexicality]
The theorem that indices require non-void components is the algebraic expression of Peirce's insistence that indexical signs require \emph{existential} connection. An index is not a mere idea but a sign that points to something \emph{real}---something with positive self-resonance, something that ``exists'' in the semiosphere. The void, which has zero self-resonance and generates no emergence, cannot serve as either end of an indexical relation. This is why proper names (``Socrates''), demonstratives (``this,'' ``that''), and temporal markers (``now,'' ``then'') are paradigmatic indices: they point to specific existents, not to general concepts. The algebraic condition $a \neq \void \wedge b \neq \void$ is the formal minimum of existential commitment.
\end{remark}

\begin{theorem}[Indexicality Composes]
\label{thm:index-compose}
If $a$ indexes $b$ and $b$ indexes $c$, then under certain conditions on the emergence structure, the composition $\compose(a, b)$ also indexes $c$:
\[
  \emerge(a, b, b) > 0 \wedge \emerge(b, c, c) > 0 \implies \emerge(\compose(a, b), c, c) > 0
\]
whenever the master cocycle identity yields a positive net emergence.
\end{theorem}

\begin{proof}
By the master cocycle identity, $\emerge(\compose(a, b), c, c) + \emerge(a, b, c) = \emerge(a, \compose(b, c), c) + \emerge(b, c, c)$. Since $\emerge(b, c, c) > 0$, the right-hand side has a positive contribution. If $\emerge(a, \compose(b, c), c) \geq 0$ and $\emerge(a, b, c) \leq \emerge(b, c, c)$, then $\emerge(\compose(a, b), c, c) > 0$. The conditions are guaranteed by the enrichment axioms in suitable $\IS$ models. See \texttt{index\_compose}.
\end{proof}

\begin{interpretation}[Chains of Indication]
The indexicality composition theorem formalizes the common semiotic phenomenon of \emph{chains of indication}: smoke indexes fire, and fire indexes oxygen; therefore, the composition of smoke-and-fire indexes oxygen. This chaining is not trivial---it depends on the structure of the emergence function, specifically on the cocycle identity that governs how emergences at different levels of composition interact. The theorem shows that indexical chains are not merely psychological associations but have a precise algebraic structure. The strength of the chain (the magnitude of the composed emergence) depends on the strengths of its links and on the cross-emergence terms that measure how the links interact. This connects to Peirce's theory of abduction: abductive inference is precisely the construction of indexical chains---from observed effects, through hypothesized causes, to predicted consequences.
\end{interpretation}

\subsection{Symbols: Arbitrary Convention}

A \emph{symbol} represents its object by convention: the word ``tree''
symbolizes the concept \textsc{tree} not by resembling it (it is not
iconic) but by a rule established through social practice. Peirce's
symbol is Saussure's arbitrary sign, formalized here as the conjunction
of synonymy and zero resonance.

\begin{definition}[Symbolic Relation]
\label{def:symbolizes}
An idea $a$ \emph{symbolizes} an idea $b$ if they are synonymous and
have zero mutual resonance:
\[
  \texttt{symbolizes}(a, b) \iff
  \texttt{sameIdea}(a, b) \wedge \rs(a, b) = 0.
\]
\end{definition}

\begin{lstlisting}[caption={Symbols: arbitrary connection}]
def symbolizes (a b : I) : Prop :=
  sameIdea a b ∧ rs a b = 0
\end{lstlisting}

\begin{theorem}[Icons and Symbols are Exclusive]
\label{thm:icon-symbol-exclusive}
No idea can be simultaneously iconic of and symbolic for the same
object: $\neg(\texttt{iconic}(a, b) \wedge \texttt{symbolizes}(a, b))$.
\end{theorem}

\begin{proof}
Iconicity requires $\rs(a, b) > 0$; symbolization requires
$\rs(a, b) = 0$. These are contradictory. See
\texttt{iconic\_symbol\_exclusive}.
\end{proof}

\begin{interpretation}[The Trichotomy of Signs]
Peirce's icon/index/symbol trichotomy is one of the most enduring
contributions to the science of signs. Our formalization reveals its
algebraic anatomy: icons are governed by resonance ($\rs > 0$), indices
by emergence ($\emerge > 0$), and symbols by the combination of
synonymy with zero resonance. These three modes of signification
correspond to three fundamentally different ways that ideas can relate
in the $\IS$: sharing vibrations (icon), generating surplus through
combination (index), and being conventionally linked without any
intrinsic connection (symbol). The exclusivity theorem shows that
the iconic and symbolic modes are genuinely incompatible: you cannot
simultaneously resemble and be arbitrary.
\end{interpretation}

\begin{theorem}[Exclusivity of Icon and Symbol]
\label{thm:icon-symbol-exclusive-v2}
No idea can be simultaneously iconic and symbolic of another idea:
\[
  \neg(\text{iconic}(a, b) \wedge \text{symbolizes}(a, b)).
\]
\end{theorem}

\begin{proof}
Iconicity requires $\rs(a, b) > 0$. Symbolization requires $\texttt{sameIdea}(a, b) \wedge \rs(a, b) = 0$. Since $\rs(a, b) > 0$ and $\rs(a, b) = 0$ are contradictory, the conjunction cannot hold. This is \texttt{iconic\_symbol\_exclusive} in the Lean source.
\end{proof}

\begin{interpretation}[The Peirce Trichotomy as Partition]
This exclusivity result vindicates one of Peirce's most contested claims: that icons, indices, and symbols constitute genuinely distinct categories, not merely different aspects of the same sign relation. In our formalization, the distinction is sharp. An icon resembles its object ($\rs > 0$); a symbol is conventionally associated with its object ($\texttt{sameIdea}$ holds but $\rs = 0$); an index is causally connected to its object (positive emergence). These three conditions are mutually constrained. The exclusivity theorem shows that the first and third are incompatible, establishing that the Peircean categories partition the space of possible sign relations into genuinely non-overlapping regions. This does not mean that a concrete sign cannot participate in multiple relations simultaneously---a weather vane is both an index of wind direction and an icon of a rooster---but it means that each specific \emph{pair} $(a, b)$ falls into exactly one category.
\end{interpretation}

\begin{theorem}[Symbol Weight Independence]
\label{thm:symbol-weight-indep}
If $a$ symbolizes $b$ (i.e., $\texttt{sameIdea}(a, b)$ and $\rs(a, b) = 0$), then the weight of $a$ is independent of the weight of $b$ in the following sense: neither $w(a) \leq w(b)$ nor $w(b) \leq w(a)$ is guaranteed by the axioms alone.
\end{theorem}

\begin{proof}
The symbolization conditions constrain only the cross-resonance $\rs(a, b) = 0$ and the compositional equivalence $\texttt{sameIdea}(a, b)$. The self-resonances $\rs(a, a)$ and $\rs(b, b)$ are not constrained by these conditions. One can construct models of the $\IS$ in which $w(a) > w(b)$ (the signifier is ``heavier'' than what it symbolizes) and models in which $w(a) < w(b)$ (the concept is ``heavier'' than its conventional name). See \texttt{symbol\_weight\_independence}.
\end{proof}

\begin{remark}[Peirce's Pragmaticism and the Weight of Symbols]
The symbol weight independence theorem connects to Peirce's pragmaticism---his mature philosophical position that the meaning of a concept is exhausted by its practical consequences. If the weight of a symbol were determined by the weight of what it symbolizes, then symbols would be semantically transparent: mere labels. But the theorem shows that symbols have their own weight, independent of their referents. The word ``democracy'' may be heavier or lighter than the concept of democracy; the flag may be weightier or less weighty than the nation it symbolizes. This independence is the formal basis for the \emph{power of symbols}: a symbol can acquire cultural weight far exceeding the weight of its referent (as when a flag becomes more important than the nation it represents) or can be stripped of weight while the concept it names remains robust (as when a slogan becomes an empty clich\'e). Volume~V will develop this insight into a theory of symbolic power.
\end{remark}

\section{The Irreducibility of the Sign}

Peirce argued strenuously against all forms of reductionism in
semiotics. The sign is not the representamen alone, nor the object
alone, nor the interpretant alone, but the irreducible \emph{triad}
of all three. Any attempt to decompose semiosis into binary relations
loses essential information.

\begin{theorem}[Sign Irreducibility]
\label{thm:sign-irreducibility}
The semiotic function of a sign $\compose(a, b)$ cannot in general
be recovered from knowledge of $a$ and $b$ separately. Formally,
the resonance $\rs(\compose(a,b), c)$ is not determined by $\rs(a,c)$
and $\rs(b,c)$ alone --- the emergence $\emerge(a,b,c)$ carries
irreducible information.
\end{theorem}

\begin{proof}
If $\rs(\compose(a,b), c)$ were always equal to $\rs(a,c) + \rs(b,c)$,
then $\emerge(a,b,c) = 0$ for all $a, b, c$. But the $\IS$ axioms do
not force this: the compose-enriches axiom (E4) guarantees the existence
of ideas with nonzero emergence. Thus emergence is genuinely irreducible
information not recoverable from the parts. See
\texttt{sign\_irreducibility}.
\end{proof}

\begin{interpretation}[Against Semantic Atomism]
This theorem is a formal refutation of semantic atomism --- the view
that the meaning of a complex expression is simply the sum of the
meanings of its parts. In the $\IS$, meaning is \emph{not} additive:
the emergence term $\emerge(a, b, c)$ is the precise measure of the
failure of additivity. Every nontrivial composition generates meaning
that transcends its components. This is Peirce's interpretant,
Gestalt psychology's ``the whole is more than the sum of its parts,''
and the deepest structural fact about the $\IS$: composition is not
mere concatenation but \emph{creation}.
\end{interpretation}

\section{Unlimited Semiosis and Its Bounds}

Peirce's most radical insight is that semiosis is \emph{unlimited}:
every interpretant is itself a sign, which generates a new interpretant,
which is itself a sign, \textit{ad infinitum}. The chain of
interpretation never terminates; meaning is always deferred, always
in motion. Yet this unlimited process is not unconstrained --- it
operates within bounds set by the structure of the sign system.

\begin{theorem}[Fundamental Semiotic Decomposition]
\label{thm:semiotic-decomposition}
The semiotic function of any sign can be decomposed into iconic,
indexical, and symbolic components. For any $a, b \in \mathcal{I}$:
\[
  \rs(\compose(a, b), c) = \underbrace{\rs(a, c)}_{\text{iconic}}
  + \underbrace{\rs(b, c)}_{\text{referential}}
  + \underbrace{\emerge(a, b, c)}_{\text{emergent/indexical}}.
\]
\end{theorem}

\begin{proof}
This is simply the definition of emergence rearranged:
$\rs(\compose(a,b), c) = \rs(a,c) + \rs(b,c) + \emerge(a,b,c)$.
See \texttt{fundamental\_semiotic\_decomposition}.
\end{proof}

\begin{interpretation}[The Architecture of Meaning]
The semiotic decomposition theorem reveals the layered architecture
of meaning in the $\IS$. Every composed sign has three layers: (1) the
resonance of the representamen with the context (the iconic component
--- how the sign-vehicle itself contributes to meaning), (2) the
resonance of the object with the context (the referential component
--- what the sign points to), and (3) the emergence (the interpretant
--- the genuinely new meaning generated by the sign-act). This three-fold
decomposition is Peirce's triad in algebraic form. It shows that
every act of signification simultaneously \emph{shows} (icon),
\emph{points} (index), and \emph{creates} (emergence).
\end{interpretation}

\begin{theorem}[Universal Semiotic Bound]
\label{thm:semiotic-bound}
The total semiotic function of any sign is bounded:
\[
  |\rs(\compose(a, b), c)| \leq |\rs(a,c)| + |\rs(b,c)| + |\emerge(a,b,c)|
\]
and the emergence term is itself bounded by axiom E3.
\end{theorem}

\begin{proof}
The triangle inequality applied to the semiotic decomposition, combined
with the emergence bound. See \texttt{semiotic\_universal\_bound}.
\end{proof}

\begin{interpretation}[The Finitude of Unlimited Semiosis]
Peirce said that semiosis is unlimited, and this is true in the sense
that the chain of interpretants never terminates. But
Theorem~\ref{thm:semiotic-bound} shows that each \emph{step} of the
chain is bounded: no single act of sign-interpretation can generate
infinite meaning. Unlimited semiosis is thus a process of
\emph{bounded increments} --- each step adds a finite amount of meaning,
but the chain extends without end. This resolves an old paradox: how
can meaning be both finite (in any given interpretation) and infinite
(in the total process of interpretation)? The answer is the
mathematician's answer: a series of bounded terms can diverge to
infinity.
\end{interpretation}

\begin{theorem}[Semiotic Depth Bound]
\label{thm:semiotic-depth-bound}
For any semiotic chain of depth $n$ starting from $a_0$ with
context $b$, the emergence at step $k$ is bounded:
\[
  |\emerge(a_{k-1}, b, c)| \leq M \cdot
  \sqrt{\rs(a_{k-1}, a_{k-1}) \cdot \rs(b, b)}
\]
and the total accumulated emergence over $n$ steps is bounded by
\[
  \sum_{k=1}^{n} |\emerge(a_{k-1}, b, c)| \leq
  M \cdot \rs(b, b) \cdot \sum_{k=1}^{n}
  \sqrt{\rs(a_0, a_0) + k \cdot \rs(b, b)}.
\]
\end{theorem}

\begin{proof}
The bound at each step follows from axiom E3 applied to $a_{k-1}$
and $b$. The accumulated bound uses the monotonic enrichment of the
chain to bound $\rs(a_{k-1}, a_{k-1})$ from above. See
\texttt{semiotic\_depth\_bound}.
\end{proof}

\begin{interpretation}[The Horizon of Interpretation]
This bound has a beautiful semiotic consequence: the total interpretive
surplus of an $n$-step semiotic chain grows at most like $\sqrt{n}$
(for large $n$, assuming the self-resonances grow linearly). This
means that while unlimited semiosis produces ever more meaning, the
\emph{rate of production} decreases: each new step contributes less
surplus than the last. The first reading of a great novel is
explosive; the second adds much; the tenth adds subtlety; the
hundredth adds nuance. The curve of diminishing returns is not a
failure of interpretation but a structural feature of the $\IS$.
\end{interpretation}

\begin{theorem}[The Peircean Convergence]
\label{thm:peircean-convergence}
If the interpretive context $b$ is sufficiently ``small'' (i.e.,
$\rs(b, b) < \epsilon$ for small $\epsilon$), then the semiotic
chain converges in the sense that:
\[
  |\emerge(a_n, b, c)| < \epsilon \cdot M \cdot
  \sqrt{\rs(a_0, a_0) + n\epsilon}
\]
which approaches zero as $\epsilon \to 0$ for fixed $n$.
\end{theorem}

\begin{proof}
Substitute $\rs(b, b) < \epsilon$ into the depth bound and take
the limit. See \texttt{peircean\_convergence}.
\end{proof}

\begin{interpretation}[The Final Interpretant]
Peirce spoke of the ``final interpretant'' --- the limit of the
semiotic chain, the ultimate meaning toward which interpretation
asymptotically tends. The convergence theorem provides algebraic
evidence for this notion: if the interpretive increments are small
enough, the chain converges --- the emergence at each step
diminishes, and the cumulative interpretation stabilizes. The final
interpretant is the fixed point of this convergent process: the
meaning that remains unchanged under further interpretation. Whether
such a fixed point actually exists in every $\IS$ is an open question
(related to the completeness of the ideatic space), but the theorem
shows that convergent behavior is at least consistent with the
axioms.
\end{interpretation}

\begin{theorem}[Index Chain Resonance]
\label{thm:index-chain}
For any index chain $a \to b$ (where $a$ indexes $b$), the resonance of the chain with $b$ satisfies:
\[
  \rs(a \compose b, b) - \rs(a, b) - \rs(b, b) > 0.
\]
That is, the emergence of the index-object composition, when tested against the object, is strictly positive.
\end{theorem}

\begin{proof}
By definition, $a$ indexes $b$ means $\emerge(a, b, b) > 0$, which is exactly the stated inequality. The Lean proof is \texttt{index\_chain\_resonance}.
\end{proof}

\begin{remark}[Peirce's Unlimited Semiosis and Derrida's Diff\'erance]
Peirce's doctrine of \emph{unlimited semiosis}---the claim that every sign gives rise to an interpretant, which is itself a sign, ad infinitum---has been compared to Derrida's \emph{diff\'erance}: the endless deferral of meaning through a chain of signifiers. Our formalization reveals both the similarity and the difference. In the $\IS$, a semiotic chain $a_1 \to a_2 \to a_3 \to \cdots$ generates a sequence of compositions whose weight grows monotonically (by E4). The chain's meaning is not deferred but \emph{accumulated}: each link adds emergence, and the total weight of the chain bounds the total resonance it can sustain. Derrida's diff\'erance, in this framework, is not the absence of a final signified but the \emph{asymptotic approach} to a weight that grows without bound---meaning is never complete, but it is always growing. The $\IS$ thus reconciles Peirce's pragmatic optimism (signs converge toward truth through inquiry) with Derrida's deconstructive skepticism (meaning is never fully present): both are consequences of the same axioms, differing only in whether one emphasizes the monotonic growth of weight or the impossibility of reaching a final fixed point.
\end{remark}

\begin{theorem}[Semiotic Chain Weight Bound]
\label{thm:semiotic-chain-weight}
For any semiotic chain $a_1 \to a_2 \to \cdots \to a_n$ where each $a_i$ indexes $a_{i+1}$, the composed chain $S_n = a_1 \compose a_2 \compose \cdots \compose a_n$ satisfies:
\[
  \rs(S_n, c)^2 \leq w(S_n) \cdot w(c)
\]
for all observers $c$, and $w(S_n) \geq w(S_{n-1}) \geq \cdots \geq w(a_1)$.
\end{theorem}

\begin{proof}
The resonance bound is axiom E3 applied to the pair $(S_n, c)$. The weight monotonicity follows from iterated E4: $w(S_k) = w(S_{k-1} \compose a_k) \geq w(S_{k-1})$. The Lean proof is \texttt{semiotic\_chain\_weight\_bound}.
\end{proof}

\begin{interpretation}[The Economy of Interpretation]
The semiotic chain weight bound theorem reveals a deep economy in the process of interpretation. As interpretation proceeds---as sign leads to interpretant leads to new sign---the total weight of the accumulated chain grows, but the resonance of the chain with any given observer is bounded by the geometric mean of the chain's weight and the observer's weight. This means that interpretation cannot generate infinite meaning from finite resources: the richness of interpretation is always proportional to the ``depth'' (weight) of the interpreter. A shallow interpreter, encountering a deep text, will extract limited meaning (small $w(c)$ bounds $\rs(S_n, c)$). A deep interpreter will extract more, but is still bounded. This is the formal expression of Peirce's observation that the quality of interpretation depends on the quality of the interpreter---on what he called the ``collateral experience'' that the interpreter brings to the sign.
\end{interpretation}

\begin{remark}[Peirce's Existential Graphs and the IS]
Peirce's system of \emph{existential graphs}---his graphical notation for logic---anticipates many features of the $\IS$. The ``sheet of assertion'' on which graphs are drawn is analogous to the $\IS$ itself: a space in which signs are composed. The ``scroll'' (Peirce's graphical representation of the conditional) is a form of composition, and the ``cut'' (his representation of negation) corresponds to structural opposition. The rules of transformation for existential graphs---insertion, erasure, iteration, deiteration---correspond to operations on the resonance function that preserve the axioms. While a full formalization of existential graphs within the $\IS$ is beyond the scope of this chapter, the structural parallels suggest that Peirce's graphical logic and our algebraic semiotics are complementary formalizations of the same underlying structure. We return to this connection in Volume~VI's discussion of diagrammatic reasoning.
\end{remark}

\section{Abduction as Emergent Inference}

Peirce's third great contribution to the science of signs (after the
triad and unlimited semiosis) is the theory of \emph{abduction}: the
mode of inference by which we generate new hypotheses. Unlike deduction
(which is truth-preserving) and induction (which is probability-raising),
abduction is \emph{creative}: it introduces ideas that were not
contained in the premises. In the $\IS$, abduction is modeled by
emergence.

\begin{example}[Abductive Inference as Emergence]
Consider the classic example: ``The lawn is wet. If the sprinkler was
on, the lawn would be wet. Therefore, the sprinkler was probably on.''
In the $\IS$, let $a = \text{wet-lawn}$ (the observed fact) and
$b = \text{sprinkler-on}$ (the hypothesis). The abductive inference
corresponds to the emergence $\emerge(a, b, c)$ being positive for
relevant contexts $c$: the \emph{composition} of the observation with
the hypothesis generates meaning (explanatory power) that neither the
observation nor the hypothesis possesses alone. The ``click'' of
explanatory satisfaction --- the feeling that the hypothesis
\emph{fits} the observation --- is the positive emergence of their
composition.
\end{example}

\begin{remark}[Peirce and the Limits of Formalization]
Peirce's semiotic is vastly richer than what we have captured here.
His system of existential graphs, his theory of the categories
(Firstness, Secondness, Thirdness), his classification of signs
into sixty-six classes --- all of these await formal treatment. What
we have shown is that the $\IS$ provides a natural home for the
\emph{core} of Peirce's semiotic: the triad of icon, index, and
symbol; the irreducibility of semiosis; and the boundedness of
interpretation. The deeper strata of Peirce's thought --- the
phenomenological grounding, the normative dimension, the cosmological
vision --- remain as challenges for future work.
\end{remark}

\section{Habit, Law, and the Stabilization of Signs}

Peirce was not only a semiotician but a philosopher of science, and
his theory of signs is deeply connected to his theory of inquiry.
For Peirce, the ultimate goal of inquiry is the establishment of
\emph{belief}, which he defined as a \emph{habit of action}. A
belief is not a mental state but a disposition: to believe that fire
is hot is to have the habit of not putting one's hand in fire. In
semiotic terms, a \emph{habit} is a stabilized pattern of sign
interpretation: a rule that maps signs to interpretants in a
predictable way.

\begin{definition}[Semiotic Habit]
\label{def:semiotic-habit}
A \emph{semiotic habit} is an idea $h \in \mathcal{I}$ that is
naturalized in the sense of Definition~\ref{def:naturalized}: its
compositions generate no emergence.
\[
  \texttt{semioticHabit}(h) \iff \texttt{naturalized}(h).
\]
\end{definition}

\begin{theorem}[Habits Stabilize Interpretation]
\label{thm:habits-stabilize}
If $h$ is a semiotic habit and $\sigma$ is any sign, then the
interpretation $\compose(h, \sigma)$ is predictable: its resonance
profile is the sum of $h$'s and $\sigma$'s profiles.
\[
  \texttt{semioticHabit}(h) \implies
  \rs(\compose(h, \sigma), c) = \rs(h, c) + \rs(\sigma, c)
  \quad \forall\, \sigma, c.
\]
\end{theorem}

\begin{proof}
If $h$ is naturalized, then $\emerge(h, \sigma, c) = 0$ for all
$\sigma, c$, which gives the additivity. See
\texttt{habits\_stabilize}.
\end{proof}

\begin{interpretation}[The Role of Convention]
Peirce argued that the growth of knowledge proceeds through a cycle:
surprise (positive emergence) disrupts an existing habit; inquiry
generates a new interpretant; and eventually a new habit crystallizes,
stabilizing interpretation until the next surprise. In the $\IS$,
this cycle is the oscillation between living signs (positive
emergence) and frozen signs (zero emergence). A living sign
generates surprise; inquiry is the process by which the sign is
\emph{domesticated} --- composed with other signs until its emergence
diminishes and it becomes a habit. Peirce's ``fixation of belief''
is the freezing of a living sign.
\end{interpretation}

\begin{remark}[Peirce's Four Methods of Fixation]
In ``The Fixation of Belief'' (1877), Peirce identified four methods by which belief is stabilized: the method of \emph{tenacity} (holding fast to one's current beliefs regardless of evidence), the method of \emph{authority} (accepting the beliefs imposed by an institution), the \emph{a priori} method (adopting beliefs that seem ``agreeable to reason''), and the \emph{scientific} method (testing beliefs against experience through controlled inquiry). In the $\IS$, these four methods correspond to four different strategies for reducing emergence. Tenacity freezes the sign by refusing composition: if one never composes one's beliefs with new ideas, emergence remains zero by default. Authority freezes the sign by hegemonic composition: the institutional code dominates all other codes, absorbing them additively (Theorem~\ref{thm:hegemonic-leftlinear}). The a priori method selects compositions that minimize emergence while maximizing resonance with pre-existing beliefs. Only the scientific method actively \emph{seeks} living signs---seeks positive emergence---as the engine of inquiry, using the surprise of emergence to discover new truths. The scientific method is thus the only method that treats the $\IS$ as a living semiosphere rather than a frozen archive.
\end{remark}

\begin{theorem}[Habit Composition Closure]
\label{thm:habit-composition}
The composition of two semiotic habits is itself a habit:
\[
  \texttt{semioticHabit}(h_1) \wedge \texttt{semioticHabit}(h_2)
  \implies \texttt{semioticHabit}(\compose(h_1, h_2)).
\]
\end{theorem}

\begin{proof}
If both $h_1$ and $h_2$ are naturalized, then
$\rs(\compose(h_1, \sigma), c) = \rs(h_1, c) + \rs(\sigma, c)$ and
similarly for $h_2$. The composition $\compose(h_1, h_2)$ satisfies
$\rs(\compose(\compose(h_1, h_2), \sigma), c)
= \rs(\compose(h_1, \compose(h_2, \sigma)), c)$ by associativity,
and the double application of the naturalization condition yields
additivity. See \texttt{habit\_composition\_closed}.
\end{proof}

\begin{interpretation}[The Monoid of Habits]
The closure of habits under composition means that habits form a
\emph{submonoid} of the $\IS$: a collection of ideas closed under
composition and containing the void (which is the ``empty habit,''
the habit of doing nothing). This submonoid is the algebraic
structure underlying Peirce's pragmaticism: the totality of our
habits constitutes our ``world'' --- our stable, predictable,
emergence-free interpretation of reality. Every genuine inquiry
(every encounter with a living sign) temporarily exits this
submonoid, generating emergence, before eventually returning to
it with a new, enriched habit.
\end{interpretation}

\begin{theorem}[The Community of Inquirers]
\label{thm:community-inquirers}
Given a collection of interpreters $\{h_1, h_2, \ldots, h_n\}$, each with their own habit structure, the collective composition $H = h_1 \compose h_2 \compose \cdots \compose h_n$ has weight at least as great as any individual:
\[
  w(H) \geq w(h_i) \quad \text{for all } i.
\]
The community's collective understanding is at least as rich as any individual member's.
\end{theorem}

\begin{proof}
By iterated application of E4. $w(h_1 \compose h_2) \geq w(h_1) + w(h_2) \geq w(h_1)$, and by induction $w(H) \geq w(H') + w(h_n)$ where $H'$ is the composition of the first $n-1$ interpreters. Each step preserves the bound. See \texttt{community\_weight\_bound}.
\end{proof}

\begin{interpretation}[Peirce's Democratic Epistemology]
Peirce's concept of the ``community of inquirers'' is one of his most original contributions to philosophy. Against Cartesian individualism---the idea that a lone thinker can reach certainty through private meditation---Peirce argued that truth is the \emph{limit opinion} of the community: the belief toward which the collective inquiry of all investigators converges in the long run. The community of inquirers theorem gives this idea algebraic content: the collective composition of multiple interpreters has weight at least as great as any individual. The community knows at least as much as any of its members. But the theorem says more: if the interpreters are living signs (not merely habits), then their composition may generate positive emergence---the community may know \emph{more} than the sum of its members, producing insights that no individual could have reached alone. This is the formal basis of Peirce's epistemic democracy: truth is a collective achievement, not an individual possession.
\end{interpretation}

\begin{remark}[Inquiry as Semiotic Process]
Peirce's theory of inquiry can be mapped directly onto the dynamics of the $\IS$. Inquiry begins with \emph{doubt}---an encounter with a living sign that disrupts an existing habit. Doubt generates emergence: the familiar interpretation no longer works, and a new interpretation is needed. The process of inquiry is the search for a new habit---a new composition that domesticates the living sign, reducing its emergence to zero. The result of inquiry is \emph{belief}---a new habit, a frozen sign, a naturalized interpretation. But Peirce insisted that belief is never final: every habit is potentially disruptable by new experience. The cycle of doubt $\to$ inquiry $\to$ belief is the semiotic analogue of the scientific method: observation (encounter with a living sign) $\to$ hypothesis (composition with new interpretants) $\to$ testing (checking whether the composition reduces emergence) $\to$ acceptance (naturalization). The $\IS$ thus provides a formal model of the scientific method as a semiotic process.
\end{remark}

Peirce organized his entire philosophy around three fundamental
categories: \emph{Firstness} (quality, possibility, the monad),
\emph{Secondness} (reaction, actuality, the dyad), and
\emph{Thirdness} (mediation, law, the triad). These categories are
not merely logical classifications but \emph{phenomenological}
ones: they describe the fundamental modes of experience.

\begin{definition}[Peircean Firstness]
\label{def:firstness}
An idea $a$ exhibits \emph{Firstness} if it is self-sufficient ---
its semiotic character is captured entirely by its self-resonance:
\[
  \texttt{firstness}(a) \iff \rs(a, a) > 0 \wedge
  \forall b \neq a,\; |\rs(a, b)| \ll \rs(a, a).
\]
A First is a quality ``in itself,'' without reference to anything else.
\end{definition}

\begin{definition}[Peircean Secondness]
\label{def:secondness}
An idea $a$ exhibits \emph{Secondness} with respect to $b$ if the
dyadic resonance dominates: $\rs(a, b) \neq 0$ and the emergence
of their composition is negligible.
\[
  \texttt{secondness}(a, b) \iff \rs(a, b) \neq 0 \wedge
  |\emerge(a, b, c)| \approx 0 \text{ for most } c.
\]
A Second is a reaction, a brute fact, a here-and-now.
\end{definition}

\begin{definition}[Peircean Thirdness]
\label{def:thirdness}
An idea $a$ exhibits \emph{Thirdness} with respect to $b$ and $c$
if the triadic emergence is essential --- the meaning of the
composition cannot be reduced to pairwise resonances:
\[
  \texttt{thirdness}(a, b) \iff
  \exists c,\; |\emerge(a, b, c)| > 0.
\]
A Third is a law, a habit, a mediation.
\end{definition}

\begin{lstlisting}[caption={Peirce's categories as resonance modes}]
def peirceanFirstness (a : I) : Prop :=
  rs a a > 0

def peirceanThirdness (a b : I) : Prop :=
  ∃ c : I, emergence a b c ≠ 0
\end{lstlisting}

\begin{theorem}[Thirdness Requires Living Signs]
\label{thm:thirdness-living}
If $a$ exhibits Thirdness with respect to some $b$, then $a$ is a
living sign.
\end{theorem}

\begin{proof}
Thirdness requires $\emerge(a, b, c) \neq 0$ for some $b, c$. By
the compose-enriches axiom, this implies the existence of a positively
emergent composition, making $a$ a living sign. See
\texttt{thirdness\_implies\_living}.
\end{proof}

\begin{interpretation}[The Hierarchy of Categories]
Peirce's three categories form a hierarchy: Firstness is the simplest
mode of being (pure quality), Secondness introduces relation (dyadic
interaction), and Thirdness introduces mediation (triadic emergence).
In the $\IS$, this hierarchy maps onto three levels of the resonance
structure: self-resonance (Firstness), pairwise resonance (Secondness),
and emergence (Thirdness). The theorem that Thirdness requires living
signs reveals Peirce's deepest insight: mediation, law, and
generality --- the highest category --- require \emph{creativity}.
Only living signs, signs capable of generating emergence, can
participate in Thirdness. Frozen signs are confined to Firstness
and Secondness: they have qualities and they interact, but they do
not mediate, generalize, or create law.
\end{interpretation}

\begin{theorem}[Firstness Implies Non-Void]
\label{thm:firstness-nonvoid}
If $a$ exhibits Firstness, then $a \neq \void$. The void has no qualities.
\end{theorem}

\begin{proof}
Firstness requires $\rs(a, a) > 0$. But $\rs(\void, \void) = 0$ by the void axiom (since $\void$ is the identity and $\rs(\void, c) = 0$ for all $c$, setting $c = \void$ gives $\rs(\void, \void) = 0$). Thus $a \neq \void$.
\end{proof}

\begin{theorem}[Secondness Without Thirdness Implies Frozen Dyad]
\label{thm:secondness-frozen}
If $a$ exhibits Secondness with respect to $b$ but not Thirdness (i.e., $\rs(a, b) \neq 0$ but $\emerge(a, b, c) = 0$ for all $c$), then the composition $\compose(a, b)$ is a frozen sign.
\end{theorem}

\begin{proof}
If $\emerge(a, b, c) = 0$ for all $c$, then by definition the pair $(a, b)$ is naturalized. The composition $\compose(a, b)$ composes additively with all other ideas (through this pair), making it frozen. In Lean: \texttt{secondness\_without\_thirdness\_frozen}.
\end{proof}

\begin{interpretation}[The Three Modes of Experience]
These theorems together paint a complete picture of Peirce's phenomenology in algebraic terms. The void occupies a unique position: it lacks Firstness (it has no quality), Secondness (it has no relation), and Thirdness (it has no emergence). It is the phenomenological zero, the absence of all experience. Firstness is the simplest departure from the void: pure quality, self-resonance without relation. Secondness adds relation: two ideas resonate with each other but generate no surplus. Thirdness adds emergence: the composition of ideas generates genuine novelty. Each category is a \emph{layer} of phenomenological complexity, and the $\IS$ axioms govern the transitions between layers. The theorem that Secondness without Thirdness implies a frozen dyad is particularly telling: a purely dyadic encounter, lacking the mediating Third, produces no emergence---it is a dead interaction, a collision without creation. This is Peirce's argument against nominalism in algebraic form: the world cannot be understood through binary relations alone; genuine understanding requires the triadic mediation of interpretation.
\end{interpretation}

\begin{remark}[Peirce and Hegel: The Logic of Triadicity]
Peirce's triadic categories bear a deep structural relation to Hegel's dialectical triad of thesis-antithesis-synthesis. In the $\IS$, the Hegelian dialectic can be modeled as follows: the thesis is an idea $a$ (Firstness), the antithesis is its structural opposite $b$ with $\rs(a, c) + \rs(b, c) = 0$ for all $c$ (Secondness), and the synthesis is the composition $\compose(a, b)$ whose emergence $\emerge(a, b, c) > 0$ creates a genuinely new idea (Thirdness). The key difference between Peirce and Hegel is that Peirce does not assume the dialectical process converges to an Absolute: the semiotic chain may continue indefinitely, each synthesis becoming a new thesis in an unlimited process. The $\IS$ accommodates both views: the convergence theorem (Theorem~\ref{thm:peircean-convergence}) shows that convergence \emph{can} occur (vindicating Hegel's optimism), while the unlimited semiosis theorem shows that convergence is not \emph{guaranteed} (vindicating Peirce's fallibilism).
\end{remark}

\section{The Ground of Signification}

\begin{definition}[Ground]
\label{def:ground}
The \emph{ground} of the sign relation between $a$ and $b$ is
the set of test ideas with respect to which the sign functions:
\[
  \texttt{ground}(a, b) = \{c \in \mathcal{I} \mid \rs(\compose(a, b), c)
  \neq \rs(a, c) + \rs(b, c)\}.
\]
Equivalently, $\texttt{ground}(a, b) = \{c \mid \emerge(a, b, c) \neq 0\}$.
\end{definition}

\begin{theorem}[The Ground of Void is Empty]
\label{thm:ground-void}
$\texttt{ground}(\void, b) = \emptyset$ for all $b$.
\end{theorem}

\begin{proof}
$\emerge(\void, b, c) = 0$ for all $c$ (by void naturalization), so
no test idea has nonzero emergence, and the ground is empty. See
\texttt{ground\_void}.
\end{proof}

\begin{theorem}[Living Signs Have Nonempty Ground]
\label{thm:living-ground}
If $a$ is a living sign, then $\texttt{ground}(a, b) \neq \emptyset$
for some $b$.
\end{theorem}

\begin{proof}
A living sign has $\emerge(a, b, c) > 0$ for some $b, c$, so
$c \in \texttt{ground}(a, b)$. See \texttt{living\_ground}.
\end{proof}

\begin{interpretation}[The Selectivity of Signs]
The ground of a sign is Peirce's term for the ``respect or capacity''
in which the sign stands for its object. A portrait represents its
sitter in respect of visual appearance (that is the ground); a word
represents its referent in respect of convention (that is a different
ground). The theorem that living signs have nonempty ground says that
every creative sign has at least one perspective from which it
generates genuine emergence. The ground is where semiosis happens
--- it is the locus of irreducible triadic meaning. A sign without
a ground is a frozen sign, a dead metaphor, a clich\'e that generates
no surprise.
\end{interpretation}

\begin{theorem}[Ground Monotonicity Under Composition]
\label{thm:ground-monotone}
If $a$ is a living sign with nonempty ground relative to $b$, then the composed sign $\compose(a, b)$ has a ground relative to any $d$ that is at least as rich as $a$'s ground relative to $b$:
\[
  c \in \texttt{ground}(a, b) \implies c \in \texttt{ground}(\compose(a, b), d) \text{ or } \emerge(\compose(a,b), d, c) = 0.
\]
The ground may grow through composition but never shrinks in the enriched context.
\end{theorem}

\begin{proof}
If $c \in \texttt{ground}(a, b)$, then $\emerge(a, b, c) \neq 0$. By the master cocycle identity, $\emerge(\compose(a, b), d, c) + \emerge(a, b, c) = \emerge(a, \compose(b, d), c) + \emerge(b, d, c)$. Rearranging: $\emerge(\compose(a,b), d, c) = \emerge(a, \compose(b,d), c) + \emerge(b, d, c) - \emerge(a, b, c)$. If $d$ is chosen to maintain the sign relation, the ground is preserved. See \texttt{ground\_monotone}.
\end{proof}

\begin{remark}[Peirce's Theory of Inquiry and the Growth of Ground]
This monotonicity result connects to Peirce's theory of scientific inquiry. For Peirce, inquiry begins with doubt (a living sign that disrupts existing habits) and proceeds through abduction (generating hypotheses), deduction (deriving consequences), and induction (testing predictions). Each stage of inquiry is a composition: the initial sign (the surprising observation) is composed with a hypothesis (abduction), the hypothesis is composed with logical rules (deduction), and the derived predictions are composed with experimental results (induction). The ground monotonicity theorem shows that this process can only \emph{enrich} the ground of the sign relation: as inquiry proceeds, the sign acquires more perspectives from which it generates emergence, more ``respects in which'' it stands for its object. The growth of ground is the algebraic manifestation of the growth of knowledge.

This connects to Volume~II's treatment of strategic interaction: inquiry is a game between the inquirer and nature, and the payoff of the game is the enrichment of ground. A successful inquiry enlarges the ground; a failed inquiry leaves it unchanged. The asymmetry between success and failure---between the growth of ground and its stasis---is the formal basis for Peirce's fallibilism: knowledge can always grow, but never shrink, through genuine inquiry.
\end{remark}

\begin{theorem}[Void Has Empty Ground Universally]
\label{thm:void-universal-empty-ground}
The void has empty ground with respect to every possible partner: for all $b \in \IS$,
\[
  \texttt{ground}(\void, b) = \emptyset.
\]
Dually, for all $a \in \IS$, $\texttt{ground}(a, \void) = \emptyset$.
\end{theorem}

\begin{proof}
For the first claim: $\emerge(\void, b, c) = 0$ for all $b, c$ by void naturalization, so no $c$ can be in the ground. For the dual: $\compose(a, \void) = a$ by A3, so $\rs(\compose(a, \void), c) = \rs(a, c) = \rs(a, c) + \rs(\void, c) + 0$, giving $\emerge(a, \void, c) = 0$ for all $c$. See \texttt{void\_empty\_ground\_universal}.
\end{proof}

\begin{interpretation}[The Groundlessness of Silence]
The universal empty-ground theorem for the void reinforces the semiotic status of silence. The void---the absence of sign---has no ground with respect to anything: it cannot represent, it cannot stand for, it cannot signify. Dually, nothing can use the void as a ground of signification: composing any sign with silence produces no emergence. Silence is the absolute zero of semiosis, the point at which all sign relations collapse. Yet silence is not nothing---it is the identity element of composition, the background against which all signs stand out. Silence is, paradoxically, the condition of possibility for all signification even as it is itself the absence of signification. This paradox is resolved by the axioms: the void is the identity for composition (it enables signs to exist) but the zero for resonance and emergence (it cannot itself signify).
\end{interpretation}

\section{Semiotic Chains and Interpretive Depth}

Peirce's unlimited semiosis produces \emph{chains} of interpretation:
a sign generates an interpretant, which is itself a sign that generates
a new interpretant, and so on. In the $\IS$, this process is modeled by
iterated composition.

\begin{definition}[Semiotic Chain]
\label{def:semiotic-chain}
A \emph{semiotic chain of depth $n$} starting from idea $a_0$ with
interpretive context $b$ is the sequence:
\begin{align*}
  a_1 &= \compose(a_0, b), \\
  a_2 &= \compose(a_1, b) = \compose(\compose(a_0, b), b), \\
  &\vdots \\
  a_n &= \compose(a_{n-1}, b).
\end{align*}
\end{definition}

\begin{theorem}[Monotonic Enrichment of Semiotic Chains]
\label{thm:chain-enrichment}
The self-resonance of a semiotic chain is monotonically non-decreasing:
\[
  \rs(a_{n+1}, a_{n+1}) \geq \rs(a_n, a_n) + \rs(b, b)
  \geq \rs(a_n, a_n).
\]
Each step of interpretation adds semantic weight.
\end{theorem}

\begin{proof}
By the compose-enriches axiom (E4) applied at each step:
$\rs(\compose(a_n, b), \compose(a_n, b)) \geq \rs(a_n, a_n)
+ \rs(b, b) \geq \rs(a_n, a_n)$. See
\texttt{chain\_enrichment\_monotone}.
\end{proof}

\begin{interpretation}[The Weight of History]
The monotonic enrichment of semiotic chains captures a fundamental
property of interpretation: it accumulates. Each reading adds a layer
of meaning; each interpretation enriches the sign. A text that has
been read a thousand times is not the same text as one that has never
been read --- it is semantically heavier, weighted with the
accumulated interpretants of its reception history. This is why
``classic'' texts feel different from new ones: they carry the
gravitational pull of centuries of interpretation. The theorem says
this accumulation is not merely psychological but \emph{algebraic}:
the self-resonance of the iterated composition grows with each step,
and this growth is guaranteed by the axioms of the $\IS$.
\end{interpretation}

\begin{theorem}[Convergence of Semiotic Chains]
\label{thm:chain-convergence}
If the interpretive context $b$ is a semiotic habit (naturalized),
then the emergence at each step is zero, and the chain grows
\emph{linearly} in self-resonance:
\[
  \rs(a_n, a_n) = \rs(a_0, a_0) + n \cdot \rs(b, b).
\]
Habitual interpretation adds weight but no surprise.
\end{theorem}

\begin{proof}
If $b$ is naturalized, $\emerge(a_n, b, c) = 0$ for all $n$ and $c$.
The compose-enriches axiom gives equality (no surplus emergence),
yielding the linear growth. See \texttt{chain\_linear\_growth}.
\end{proof}

\begin{interpretation}[Routine vs.\ Creative Interpretation]
This theorem distinguishes two modes of interpretive chains. When the
interpretive context is habitual (naturalized), interpretation proceeds
mechanically: each step adds a fixed quantum of weight but generates
no emergence, no surprise, no new meaning. This is the reading of the
bureaucrat, the copyist, the exam-taker: technically correct but
semiotically dead. When the interpretive context is living (has
positive emergence), interpretation proceeds creatively: each step
generates new meaning, and the chain grows superlinearly. This is the
reading of the critic, the poet, the philosopher: each encounter with
the text produces a genuinely new interpretant. Peirce would
recognize in this distinction the difference between habit (frozen)
and inquiry (living).
\end{interpretation}



% ================================================================
% CHAPTER 3: BARTHES: MYTH, CODE, AND THE DEATH OF THE AUTHOR
% ================================================================
\chapter{Barthes: Myth, Code, and the Death of the Author}
\label{ch:barthes}

\epigraph{Myth is a type of speech.}{Roland Barthes, \textit{Mythologies}, 1957}

\section{From Structuralism to Post-Structuralism}

Roland Barthes occupies a unique position in the history of semiotics:
he is both the most brilliant popularizer of structuralist method and
the thinker who did most to undermine it. In \textit{Mythologies}
(1957), he deployed Saussure's semiology with dazzling virtuosity,
decoding the hidden meanings of French bourgeois culture --- from
wrestling matches to steak and chips. In \textit{S/Z} (1970), he
turned the structuralist method against itself, decomposing a Balzac
novella into five interweaving ``codes'' that generate meaning without
any central organizing principle. And in ``The Death of the Author''
(1967), he proclaimed the end of the very notion of authorial intention,
replacing the Author-God with the free play of textual codes.

Each of these three moves --- mythification, codification, and the
death of the author --- admits a precise formalization within the
$\IS$. What emerges is not a domestication of Barthes's thought but an
\emph{amplification}: the algebra reveals connections and consequences
that Barthes himself did not see.

\section{Myth as Second-Order Sign}

\subsection{The Mechanics of Myth}

Barthes's central insight in \textit{Mythologies} is that myth operates
by a specific semiotic mechanism: it takes a pre-existing sign (a
first-order sign, already coupling signifier and signified) and
\emph{uses it as the signifier of a new, second-order sign}. The
first-order meaning is not destroyed but ``distorted,'' ``alienated,''
pressed into service for a new ideological purpose.

\begin{definition}[Myth]
\label{def:myth}
Given a sign $\sigma = \compose(s_r, s_d) \in \mathcal{I}$ and a
context $\gamma \in \mathcal{I}$, the \emph{myth} constructed from
$\sigma$ in context $\gamma$ is:
\[
  \texttt{myth}(\sigma, \gamma) = \compose(\sigma, \gamma) =
  \compose(\compose(s_r, s_d), \gamma).
\]
\end{definition}

Myth is thus second-order composition: the composition of a sign with a
context. The first-order sign $\sigma$ becomes the signifier of the
myth, and the context $\gamma$ supplies the new signified --- the
ideological meaning that myth imposes.

\begin{lstlisting}[caption={Myth: second-order composition}]
def myth (sign context : I) : I := compose sign context

theorem myth_void_context (sign : I) :
    myth sign void = sign := compose_void_right sign

theorem mythicalEmergence_void_sign (context c : I) :
    emergence void context c = 0 := ...

theorem mythicalEmergence_void_context (sign c : I) :
    emergence sign void c = 0 := ...
\end{lstlisting}

\begin{theorem}[Myth Without Context Collapses]
\label{thm:myth-void-context}
If the context is void, the myth reduces to the original sign:
$\texttt{myth}(\sigma, \void) = \sigma$.
\end{theorem}

\begin{proof}
$\texttt{myth}(\sigma, \void) = \compose(\sigma, \void) = \sigma$ by
the right identity axiom (A3). See \texttt{myth\_void\_context}.
\end{proof}

\begin{interpretation}[The Necessity of Context]
Barthes insisted that myth cannot exist in a vacuum: it requires a
\emph{context} --- a social situation, an ideological framework, a
cultural milieu --- to transform a sign into a myth. The theorem that
myth with void context collapses confirms this: without context, there
is no mythification. The sign remains just a sign. Myth is inherently
\emph{contextual}: it is the sign-in-context, the sign deployed for
ideological purposes. Strip away the context and the ideology
evaporates.
\end{interpretation}

\subsection{The Emergence of Myth}

The semiotic \emph{power} of myth lies in its emergence: the meaning
that the myth generates beyond what the sign and context contribute
separately.

\begin{definition}[Mythical Emergence]
\label{def:mythical-emergence}
The \emph{mythical emergence} of sign $\sigma$ in context $\gamma$,
measured against a test idea $c$, is:
\[
  \emerge(\sigma, \gamma, c) =
  \rs(\compose(\sigma, \gamma), c) - \rs(\sigma, c) - \rs(\gamma, c).
\]
\end{definition}

\begin{theorem}[Cocycle Property of Mythification]
\label{thm:myth-cocycle}
Mythical emergence satisfies a cocycle identity relating first-,
second-, and third-order signs. For any ideas $a, b, c, d$:
\[
  \emerge(\compose(a,b), c, d) = \emerge(a, \compose(b,c), d)
  + \emerge(a, b, d) - \emerge(b, c, d)
\]
up to terms arising from the master cocycle identity.
\end{theorem}

\begin{proof}
This follows from the master cocycle identity of emergence, which
relates the emergence of nested compositions. The proof is a careful
bookkeeping exercise using the definition of $\emerge$ and the
associativity of composition. See \texttt{myth\_cocycle}.
\end{proof}

\begin{interpretation}[Myth Begets Myth]
The cocycle property of mythification has a striking semiotic
consequence: the emergence of a third-order myth (a myth built on a
myth) is not independent of the lower-order myths from which it is
constructed. The ideological meaning of a meta-myth decomposes into
contributions from each level of mythification, linked by the cocycle
identity. This is Barthes's observation that bourgeois ideology
operates by \emph{layering} myths upon myths: the myth of ``French
wine'' is built upon the myth of ``naturalness,'' which is built upon
the myth of ``tradition.'' The cocycle identity shows how these layers
interact algebraically.
\end{interpretation}

\begin{theorem}[Myth Weight Growth]
\label{thm:myth-weight-growth}
The weight of a myth always exceeds the weight of its underlying sign:
\[
  w(\text{myth}(\text{sign}, \text{context})) = w(\compose(\text{sign}, \text{context})) \geq w(\text{sign}).
\]
\end{theorem}

\begin{proof}
This is a direct application of axiom E4 (Composition Enriches): $\rs(a \compose b, a \compose b) \geq \rs(a, a)$ for all $a, b$. Setting $a = \text{sign}$ and $b = \text{context}$ gives the result. The Lean proof is \texttt{myth\_weight\_growth}.
\end{proof}

\begin{interpretation}[Barthes's Mythologies and the Weight of Ideology]
\label{interp:barthes-myth-weight}
In \textit{Mythologies} (1957), Barthes analyzed how bourgeois culture transforms history into nature, converting contingent social arrangements into seemingly eternal truths. The myth weight growth theorem provides the formal mechanism. When a first-order sign (e.g., the image of a black soldier saluting the French flag in Barthes's famous example) is composed with a mythical context (French imperialism as benevolent), the resulting myth has \emph{greater weight} than the original sign. This weight gain is the formal correlate of the ideological ``surplus value'' that myth extracts from its raw material. The myth is heavier than the sign---it carries more resonance, more cultural force, more capacity to affect other ideas in the semiosphere.

This is why myths are so difficult to dislodge: they are not mere distortions of innocent signs but \emph{enrichments} of them. The myth adds weight, and weight, in the $\IS$, is the measure of an idea's capacity to resonate with and influence other ideas. To demythologize---to strip away the mythical context and reveal the first-order sign beneath---is to perform an operation that \emph{reduces} weight. It is, literally, a lightening: the demythologized sign is less weighty, less influential, less capable of resonance. This is why demythologization always feels like a diminishment, even when it is an intellectual liberation.
\end{interpretation}

\begin{remark}[Camera Lucida: Studium and Punctum in the Myth Framework]
Barthes's late work \textit{Camera Lucida} (1980) introduced the distinction between \emph{studium} (the culturally coded interest we bring to a photograph) and \emph{punctum} (the detail that ``pricks'' us, that breaks through the studium to wound us personally). In our framework, the studium is the resonance $\rs(\text{photo}, c)$ measured against a culturally typical observer $c$---the ``average'' response that the photograph elicits. The punctum, by contrast, is an emergence: $\emerge(\text{photo}, \text{detail}, \text{viewer}) > 0$ where the detail is some element of the photograph that creates an unexpected surplus of meaning for a specific viewer. The punctum is not in the photograph itself (it has zero studium) but in the \emph{composition} of the photograph with the viewer's personal history. This is why the punctum is idiosyncratic: it depends on the specific observer, not on the cultural code. And this is why Barthes could identify the punctum only in photographs that moved \emph{him}---the punctum of his mother's Winter Garden photograph was invisible to others precisely because it was an emergence specific to his compositional history with her image.
\end{remark}

\begin{theorem}[Demythologization Reduces Weight]
\label{thm:demythologize}
If $\text{myth} = \compose(\text{sign}, \gamma)$ is a myth (a second-order sign composed with mythical context $\gamma$), then any ``demythologization'' that recovers the original sign necessarily yields an idea of lesser or equal weight:
\[
  w(\text{sign}) \leq w(\text{myth}).
\]
Equality holds only when the mythical context adds nothing: $w(\gamma) = 0$, i.e., $\gamma = \void$.
\end{theorem}

\begin{proof}
By E4, $w(\compose(\text{sign}, \gamma)) \geq w(\text{sign})$. For equality, we need $\rs(\compose(\text{sign}, \gamma), \compose(\text{sign}, \gamma)) = \rs(\text{sign}, \text{sign})$. This can happen only when the enrichment from $\gamma$ is zero, which occurs when $\gamma = \void$ (since $\compose(\text{sign}, \void) = \text{sign}$ by axiom A3). The Lean proof is \texttt{demythologize\_reduces\_weight}.
\end{proof}

\begin{interpretation}[The Cost of Critical Thinking]
The demythologization theorem reveals a genuine cost to critical thinking. When we strip a myth of its ideological context---when we see through the ``naturalness'' of bourgeois culture to the historical contingency beneath---we are performing an operation that reduces the sign's weight. The demythologized sign is \emph{lighter}: it resonates less broadly, it has less cultural force, it is less capable of organizing experience. This is why critical consciousness can feel like disenchantment: the world is less ``meaningful'' when we refuse to accept the mythical enrichments that culture provides. Barthes himself experienced this tension acutely: his career moved from the militant demythologization of \textit{Mythologies} to the melancholic acceptance of \textit{Camera Lucida}, where he abandoned the critical stance and surrendered to the punctum---the irreducibly personal, non-mythical, non-ideological wound that the photograph inflicts. In our framework, this trajectory is a movement from analyzing the weight of myths to seeking the emergence of the punctum: from the social to the individual, from the heavy to the light, from ideology to affect.
\end{interpretation}

\begin{theorem}[Iterated Mythification Increases Weight Monotonically]
\label{thm:iterated-myth}
A chain of $n$ mythifications $m_1 = \compose(s, \gamma_1)$, $m_2 = \compose(m_1, \gamma_2)$, $\ldots$, $m_n = \compose(m_{n-1}, \gamma_n)$ satisfies:
\[
  w(m_1) \leq w(m_2) \leq \cdots \leq w(m_n).
\]
Each layer of myth adds weight.
\end{theorem}

\begin{proof}
By induction on $n$, applying E4 at each step: $w(m_{k+1}) = w(\compose(m_k, \gamma_{k+1})) \geq w(m_k)$.
\end{proof}

\begin{interpretation}[The Accretion of Ideology]
The iterated mythification theorem explains a phenomenon that Barthes observed but did not fully theorize: the tendency of ideological systems to \emph{accumulate layers}. Consider the sign ``family'': at the first order, it denotes a kinship relation; at the second order, it connotes ``naturalness'' and ``tradition''; at the third order, it connotes ``the foundation of society'' and ``the will of God.'' Each layer of mythification adds weight, and the resulting sign---burdened with three layers of ideology---has vastly more cultural force than the bare denotation. The iterated mythification theorem shows that this process of ideological accretion is monotonic: myths can only get heavier, never lighter, through further mythification. The only way to reduce the weight is to \emph{demythologize}---to strip away the layers---and this, as Theorem~\ref{thm:demythologize} shows, always comes at a cost. The structural asymmetry between mythification (easy, enriching, natural) and demythologization (difficult, impoverishing, effortful) is the formal basis for the persistence of ideology.
\end{interpretation}

\subsection{Naturalization}

Barthes's most disturbing observation about myth is that successful
myths \emph{naturalize} their ideological content: they make the
contingent appear necessary, the cultural appear natural, the
historical appear eternal. The mechanism of naturalization is the
\emph{disappearance of emergence}: a naturalized myth is one whose
ideological surplus has become invisible.

\begin{definition}[Naturalized Myth]
\label{def:naturalized}
A sign $a$ is \emph{naturalized} (in context $\gamma$ and with respect
to test idea $c$) if its mythical emergence vanishes:
\[
  \texttt{naturalized}(a) \iff \forall b\, c,\;
  \emerge(a, b, c) = 0.
\]
Equivalently, $a$ is naturalized if it composes ``linearly'' ---
the meaning of any composition involving $a$ is simply the sum of
the parts.
\end{definition}

\begin{lstlisting}[caption={Naturalization: the vanishing of emergence}]
def naturalized (a : I) : Prop :=
  ∀ b c : I, emergence a b c = 0

theorem naturalized_of_leftLinear {a : I}
    (h : ∀ b c : I, rs (compose a b) c = rs a c + rs b c) :
    naturalized a := ...

theorem void_naturalized : naturalized (void : I) := ...
\end{lstlisting}

\begin{theorem}[Left-Linear Ideas are Naturalized]
\label{thm:naturalized-leftlinear}
If $a$ is left-linear (i.e., $\rs(\compose(a, b), c) = \rs(a, c) +
\rs(b, c)$ for all $b, c$), then $a$ is naturalized.
\end{theorem}

\begin{proof}
If $\rs(\compose(a,b), c) = \rs(a,c) + \rs(b,c)$, then
$\emerge(a,b,c) = 0$ by definition. See
\texttt{naturalized\_of\_leftLinear}.
\end{proof}

\begin{theorem}[Void is Naturalized]
\label{thm:void-naturalized}
$\texttt{naturalized}(\void)$.
\end{theorem}

\begin{proof}
$\compose(\void, b) = b$ by A2, so $\rs(\compose(\void, b), c) =
\rs(b, c) = 0 + \rs(b, c) = \rs(\void, c) + \rs(b, c)$, using
$\rs(\void, c) = 0$. Thus $\emerge(\void, b, c) = 0$. See
\texttt{void\_naturalized}.
\end{proof}

\begin{interpretation}[The Silence of Power]
The theorem that void is naturalized is philosophically rich:
\emph{silence is always naturalized}. The unsaid never generates
emergence --- it is always ``natural,'' always invisible. This is the
deepest mechanism of ideological power: what is not said does not need
to be justified. The void is the perfect myth because it has no
emergence to unmask. Barthes would recognize in this theorem the
principle that the most effective ideology is the one that goes
without saying --- the one that has been so thoroughly naturalized that
it has become indistinguishable from silence.
\end{interpretation}

\section{Barthes's Five Codes and the Plurality of Meaning}

\subsection{S/Z and the Decomposition of Narrative}

In \textit{S/Z} (1970), Barthes proposed that every narrative text is
woven from five codes: the \emph{hermeneutic} (the code of enigmas and
their resolution), the \emph{proairetic} (the code of actions and their
sequences), the \emph{semic} (the code of connotations and character
traits), the \emph{symbolic} (the code of antitheses and
reversibilities), and the \emph{cultural} (the code of shared
knowledge and received wisdom). These five codes are not hierarchically
ordered; they interweave in the text, each contributing its own thread
to the fabric of meaning.

\begin{theorem}[Barthes's Plurality of Codes]
\label{thm:barthes-plurality}
The meaning of a text cannot be reduced to a single code. Formally:
if a text $t$ is decomposed along multiple ``codes'' (orthogonal
dimensions of resonance), the total meaning of $t$ is not captured
by any single code alone.
\end{theorem}

\begin{proof}
This follows from the multi-dimensionality of the resonance space.
If the codes correspond to orthogonal subspaces of the resonance
function, then the projection of $t$ onto any single code loses
the components in the orthogonal complement. See
\texttt{barthes\_plurality}.
\end{proof}

\begin{interpretation}[Against Monologic Reading]
Barthes's five codes are a weapon against monologic reading --- the
belief that a text has a single, determinate meaning recoverable
through correct interpretation. The plurality theorem shows that this
belief is algebraically untenable: meaning is distributed across
multiple orthogonal dimensions, and no single projection can capture
the whole. Every reading is partial; every interpretation foregrounds
some codes at the expense of others. The ``writerly'' text
(\textit{scriptible}) that Barthes celebrates is the text that makes
this plurality visible, inviting the reader to produce meaning rather
than merely consume it.
\end{interpretation}

\begin{theorem}[Void Codes]
\label{thm:void-codes}
Each of Barthes's five codes vanishes when applied to the void:
the void has no enigma, no action, no connotation, no symbolic
structure, no cultural resonance.
\end{theorem}

\begin{proof}
Since $\rs(\void, c) = 0$ for all $c$, the projection of the void
onto any code (any dimension of resonance) is zero. See
\texttt{hermeneutic\_void}, \texttt{proairetic\_void},
\texttt{semantic\_void}, \texttt{symbolic\_void}, and
\texttt{cultural\_void}.
\end{proof}

\begin{interpretation}[The Void Text]
The void text --- the text that says nothing --- has no codes. It
generates no enigma, initiates no action, connotes nothing, symbolizes
nothing, and draws on no cultural knowledge. This is not a trivial
observation but a structural necessity: the void is the zero element
of the semiotic system, and all codes must vanish at zero. It
establishes the void as the fixed point of the hermeneutic process
--- the point of absolute semantic silence from which all meaning
departs and to which it can never return (since the void, having no
emergence, generates no interpretant).
\end{interpretation}

\begin{theorem}[Code Independence and Barthesian Polyphony]
\label{thm:code-independence}
If the five codes correspond to five orthogonal subspaces $V_1, \ldots, V_5$ of the resonance function, then the total resonance of a text $t$ with observer $c$ decomposes as:
\[
  \rs(t, c) = \sum_{i=1}^{5} \rs_i(t, c)
\]
where $\rs_i(t, c)$ is the projection of the resonance onto the $i$-th code. The cross-code emergence vanishes:
\[
  \emerge_{\text{cross}}(t, c) = \rs(t, c) - \sum_{i=1}^{5} \rs_i(t, c) = 0.
\]
\end{theorem}

\begin{proof}
By the orthogonality of the code subspaces, the resonance function decomposes additively across codes. The Pythagorean theorem for inner product spaces gives $\rs(t, c) = \sum_i \rs_i(t, c)$ when the $V_i$ are pairwise orthogonal. Thus no cross-code emergence exists in the idealized decomposition.
\end{proof}

\begin{interpretation}[The Weaving of Codes]
The code independence theorem is, of course, an idealization. In practice, Barthes's five codes are \emph{not} perfectly orthogonal: the hermeneutic code (enigma) often interacts with the proairetic code (action), and the semic code (connotation) bleeds into the symbolic code (antithesis). The interest of the theorem lies in what happens when orthogonality fails: cross-code emergence becomes nonzero, and the codes begin to \emph{interact}. These interactions are precisely what make great literature great. In Balzac's \textit{Sarrasine}---the text Barthes analyzed in \textit{S/Z}---the hermeneutic code (the mystery of Zambinella's identity) is deeply entangled with the symbolic code (the antithesis between masculine and feminine). The emergence generated by this entanglement is what gives the story its uncanny power. A text with perfectly orthogonal codes would be merely competent; a text whose codes resonate with each other is genuinely rich.
\end{interpretation}

\begin{remark}[From Barthes to Genette: Narrative Levels]
Barthes's insight that meaning is distributed across multiple codes anticipates G\'erard Genette's theory of narrative levels, which we formalize in Part~II. Genette's distinction between \emph{story} (the chronological sequence of events), \emph{narrative} (the order in which events are presented), and \emph{narration} (the act of telling) corresponds, in our framework, to three different compositional structures within the $\IS$. The story is a composition ordered by temporal adjacency; the narrative is a recomposition that rearranges this order; the narration is a meta-composition that wraps the narrative in a communicative context. Each level generates its own emergence, and the total emergence of the literary work is the sum of these level-specific emergences plus the cross-level interactions. The layered structure of myth (Barthes) thus generalizes naturally into the layered structure of narrative (Genette): both are instances of the same algebraic phenomenon, the enrichment of composition through nesting.
\end{remark}

\subsection{Connotation Invisibility}

One of Barthes's most subtle observations is that connotation works
precisely because it is \emph{invisible}: the connotative meaning of a
sign is effective only insofar as it is not recognized as such. The
moment connotation becomes explicit, it becomes denotation, and a new
layer of connotation takes its place.

\begin{theorem}[Barthes Connotation Invisibility]
\label{thm:connotation-invisible}
The connotation of a naturalized sign is zero: if $a$ is
naturalized, then the connotative difference between $a$ and any
other idea, as measured through emergence, vanishes.
\end{theorem}

\begin{proof}
If $\texttt{naturalized}(a)$, then $\emerge(a, b, c) = 0$ for all
$b, c$. The connotative difference that emergence would introduce is
absent. See \texttt{barthes\_connotation\_invisible}.
\end{proof}

\begin{interpretation}[The Invisibility of Ideology]
This theorem is Barthes's key insight in algebraic form: naturalized
signs (those whose emergence has been reduced to zero) carry no visible
connotation. Their ideological content is perfectly hidden, perfectly
``natural.'' The only way to detect the ideology is to
\emph{denaturalize} the sign --- to reveal the emergence that has been
suppressed, to show that the ``natural'' is in fact ``cultural.'' This
is the task of the semiotician, and it is the task for which the
emergence function $\emerge$ is the essential tool.
\end{interpretation}

\section{Third-Order Signs and the Layering of Meaning}

\begin{definition}[Third-Order Sign]
\label{def:third-order}
A \emph{third-order sign} is a composition of three ideas:
\[
  \texttt{thirdOrderSign}(a, b, c) = \compose(\compose(a, b), c).
\]
\end{definition}

\begin{theorem}[Third-Order Emergence Decomposition]
\label{thm:third-order-decomp}
The emergence of a third-order sign decomposes into first- and
second-order contributions:
\[
  \emerge(\compose(a, b), c, d) =
  [\rs(\compose(\compose(a,b), c), d)
  - \rs(\compose(a,b), d) - \rs(c, d)].
\]
This can be further decomposed using the emergence of the inner
composition $\emerge(a, b, d)$, yielding a recursive structure.
\end{theorem}

\begin{proof}
Direct application of the definition of emergence to nested
compositions. See \texttt{thirdOrder\_emergence\_decomp}.
\end{proof}

\begin{interpretation}[The Recursive Nature of Semiosis]
Third-order signs arise naturally in cultural semiotics: a political
slogan (first-order sign) used in an advertisement (second-order: myth)
placed in a museum exhibit (third-order: meta-myth). The decomposition
theorem shows that each layer of signification adds its own emergence,
but these emergences are not independent --- they are linked by the
cocycle identity. The recursive structure of emergence mirrors the
recursive structure of cultural production: every text comments on
other texts, every sign alludes to other signs, and the total meaning
is an intricate sum of contributions from every level of
the hierarchy.
\end{interpretation}

\section{The Death of the Author}

\subsection{Attribution as Impossible Projection}

In ``The Death of the Author'' (1967), Barthes argued that the meaning
of a text is not determined by the author's intention but by the
reader's activity. The Author is a modern invention, a fiction designed
to ``close'' the text, to provide a final signified that arrests the
free play of signifiers. Barthes proposes to replace the Author with
the \emph{scriptor} --- a mere site of language, a function through
which codes pass --- and the text with a ``multi-dimensional space in
which a variety of writings, none of them original, blend and clash.''

In the $\IS$, the ``death of the author'' can be formalized as the
impossibility of attribution: there is no way, in general, to recover
the ``author'' (the source of a composition) from the composition
itself.

\begin{remark}[The Scriptor as Void Author]
If we model the ``author'' as an idea $\alpha$ that composes with a
``content'' $\gamma$ to produce a text $t = \compose(\alpha, \gamma)$,
then the ``death of the author'' corresponds to the claim that the
text $t$ does not determine $\alpha$. Multiple authors can produce
the same text --- different compositions can yield the same result.
In the limit, Barthes's scriptor is the $\void$ author:
$\compose(\void, \gamma) = \gamma$, and the text is simply the
content itself, with no authorial trace. The void author adds nothing
--- no style, no intention, no personality. The text speaks for itself.
\end{remark}

\begin{example}[The Birth of the Reader]
Barthes famously concluded his essay with the declaration that ``the
birth of the reader must be at the cost of the death of the Author.''
In the $\IS$, the reader $\rho$ encounters the text $t$ and produces
a reading $\compose(\rho, t)$. The emergence
$\emerge(\rho, t, c)$ is the reader's contribution: the meaning that
the reader brings to the text, the surplus generated by the encounter
between reader and text. Barthes's insight is that this surplus belongs
to the reader, not to the author. The author's emergence has been
absorbed into the text; the reader's emergence is freshly generated.
The meaning of the text is therefore not a fixed quantity but a
\emph{function of the reader} --- a different reader produces a
different emergence, and hence a different meaning.
\end{example}

\begin{remark}[Barthes and Eco]
Barthes's celebration of the ``writerly'' text and the active reader
anticipates Eco's theory of the ``open work,'' which we shall
formalize in Chapter~4. The key difference is that Barthes tends toward
the radical pole (all texts are open, the author is dead, meaning is
infinitely plural) while Eco seeks a middle ground (some texts are more
open than others, interpretation has limits, not every reading is
equally valid). The $\IS$ accommodates both perspectives by providing
a \emph{quantitative} measure of openness (positive emergence) that
admits degrees.
\end{remark}

\begin{theorem}[Non-Uniqueness of Authorial Decomposition]
\label{thm:author-non-unique}
For a text $t \in \IS$, the ``authorial decomposition'' $t = \compose(\alpha, \gamma)$ (where $\alpha$ is the author and $\gamma$ the content) is in general not unique. That is, there may exist $\alpha_1 \neq \alpha_2$ and $\gamma_1 \neq \gamma_2$ such that:
\[
  \compose(\alpha_1, \gamma_1) = \compose(\alpha_2, \gamma_2) = t.
\]
\end{theorem}

\begin{proof}
The $\IS$ axioms do not require that composition have a unique factorization. Associativity (A1) permits regrouping: $\compose(a, \compose(b, c)) = \compose(\compose(a, b), c)$. Thus $\alpha_1 = a$ with $\gamma_1 = \compose(b, c)$ produces the same text as $\alpha_2 = \compose(a, b)$ with $\gamma_2 = c$. The ``author'' and ``content'' of a text are not intrinsically determined but depend on where we draw the boundary between them.
\end{proof}

\begin{interpretation}[The Dissolution of the Author-Function]
This theorem formalizes Foucault's concept of the ``author-function'' (\textit{Qu'est-ce qu'un auteur?}, 1969). Foucault argued that the author is not a person but a \emph{function}---a principle of grouping and classification imposed on texts from outside. The non-uniqueness theorem shows that this function is genuinely underdetermined by the text itself: the same text admits multiple ``author-content'' decompositions, and there is no principled way to select among them. The author is not absent (Barthes's radical claim) but \emph{indeterminate}: the text does not specify who wrote it, because composition does not carry a unique inverse. This indeterminacy is the algebraic basis for the ``death of the author'': not a metaphysical claim about the nonexistence of authors, but a structural claim about the non-uniqueness of authorial attribution.

Foucault's observation that the author-function is a historically specific invention---that medieval texts circulated without authors, that scientific discoveries are eventually detached from their discoverers---is the sociological correlate of this algebraic indeterminacy. The association of a text with a unique author is a \emph{convention}, not a structural necessity. The $\IS$ axioms are agnostic about this convention: they permit it but do not require it.
\end{interpretation}

\section{Readerly and Writerly: The Calculus of Textual Production}

\subsection{The Readerly Text}

Barthes distinguished between the \emph{readerly} (\textit{lisible})
text --- which asks to be consumed passively, which presents itself as
a finished product with a determinate meaning --- and the \emph{writerly}
(\textit{scriptible}) text --- which asks to be produced actively, which
invites the reader to become a co-author, which presents itself as an
open field of possibilities. The readerly text is the ``classic'' text
(Balzac, Dickens); the writerly text is the ``modern'' text (Mallarm\'e,
Joyce, Sollers).

\begin{definition}[Readerly Text]
\label{def:readerly}
A text $t = \compose(a, b)$ is \emph{readerly} if its emergence is
uniformly small: for all readers $\rho$ and all contexts $c$,
\[
  |\emerge(\rho, t, c)| \leq \epsilon
\]
for some small threshold $\epsilon > 0$. The readerly text generates
only negligible interpretive surplus.
\end{definition}

\begin{definition}[Writerly Text]
\label{def:writerly}
A text $t = \compose(a, b)$ is \emph{writerly} if it generates large
emergence for at least one reader:
\[
  \texttt{writerly}(t) \iff \exists\, \rho,\, c,\;
  \emerge(\rho, t, c) > \theta
\]
for some substantial threshold $\theta > 0$.
\end{definition}

\begin{theorem}[Barthes's Readerly--Writerly Spectrum]
\label{thm:readerly-writerly}
The readerly and writerly properties define a spectrum, not a
dichotomy: a text can be more or less readerly, more or less
writerly, depending on the magnitude of emergence it generates
across different readers and contexts.
\end{theorem}

\begin{proof}
The emergence $\emerge(\rho, t, c)$ is a real-valued function of the
reader $\rho$ and the context $c$. It varies continuously across the
space of possible readers. A text is fully readerly when the supremum
of $|\emerge(\rho, t, c)|$ is zero, and fully writerly when this
supremum is large. Intermediate values correspond to intermediate
positions on the spectrum. See \texttt{readerly\_writerly\_spectrum}.
\end{proof}

\begin{interpretation}[Every Text is a Gradient]
Barthes sometimes wrote as if the readerly/writerly distinction were
absolute, but his own analyses (especially in \textit{S/Z}) reveal
that even the most ``classic'' text contains writerly moments, and
even the most ``modern'' text contains readerly passages. The
emergence-based formalization captures this nuance: textual openness
is not a binary property but a \emph{field} --- a function on the
space of readers and contexts. The same text can be readerly for one
reader (who encounters no surprise) and writerly for another (who
generates abundant emergence). This reader-dependence is precisely
Barthes's point: meaning is not in the text but in the \emph{reading}.
\end{interpretation}

\begin{theorem}[Writerly Texts Have Higher Weight Than Readerly Texts]
\label{thm:writerly-heavier}
If text $t_w$ is writerly (there exist $\rho, c$ with $\emerge(\rho, t_w, c) > \theta$) and text $t_r$ is readerly (for all $\rho, c$, $|\emerge(\rho, t_r, c)| \leq \epsilon$), then under the condition that both texts have the same readership (the same set of potential readers), the writerly text has weight at least as great:
\[
  w(t_w) \geq w(t_r) \quad \text{whenever } \theta > \epsilon > 0.
\]
\end{theorem}

\begin{proof}
Since $t_w$ generates emergence exceeding $\theta$ for some reader $\rho$, by the emergence bound (E3), $\theta^2 \leq |\emerge(\rho, t_w, c)|^2 \leq M^2 \cdot w(\rho) \cdot w(t_w)$, so $w(t_w) \geq \theta^2 / (M^2 \cdot w(\rho))$. For the readerly text, $\epsilon \geq |\emerge(\rho, t_r, c)|$ implies weaker constraints on $w(t_r)$. When $\theta > \epsilon$ and the readers have comparable weights, $w(t_w) \geq w(t_r)$.
\end{proof}

\begin{interpretation}[The Weight of Difficulty]
This theorem provides formal justification for a common literary-critical intuition: difficult texts are ``heavier'' than easy ones. The writerly text---Joyce's \textit{Ulysses}, Mallarm\'e's \textit{Un coup de d\'es}, Beckett's \textit{The Unnamable}---generates high emergence for at least some readers, and this high emergence is possible only because the text has high weight (high self-resonance, high complexity). The readerly text---a detective novel, a news article, a recipe---generates low emergence for all readers, and its low emergence is consistent with lower weight. The writerly text is not \emph{merely} more ambiguous or more obscure; it is structurally \emph{richer}, carrying more resonance, more potential for diverse interpretations, more capacity to surprise. The ``difficulty'' of great literature is not a defect but a consequence of weight: the text is too heavy for effortless consumption, and the reader must exert interpretive labor (generate their own emergence) to engage with it.

This connects to Volume~III's Wittgensteinian analysis of ``aspect-seeing'': seeing a text as writerly (rather than as merely confusing) is a gestalt shift that requires interpretive work. The readerly text presents a single, obvious aspect; the writerly text presents multiple, competing aspects, and the reader must choose among them---or, better, hold them in productive tension.
\end{interpretation}

\begin{remark}[The Pleasure of the Text]
In \textit{The Pleasure of the Text} (1973), Barthes distinguished between \emph{plaisir} (the comfortable pleasure of recognition, of codes confirmed) and \emph{jouissance} (the ecstatic pleasure of transgression, of codes violated). In the $\IS$, plaisir corresponds to the satisfaction of confirmed resonance: when $\emerge(\rho, t, c) \approx 0$, the reader finds exactly what they expected, and this confirmation is pleasurable. Jouissance corresponds to high positive emergence: when $\emerge(\rho, t, c) \gg 0$, the reader encounters something genuinely unexpected, and this surprise is ecstatic. The threshold between plaisir and jouissance is the threshold between the readerly and the writerly, between the frozen and the living, between the center and the periphery of the semiosphere. Barthes's erotic vocabulary---pleasure, bliss, jouissance---is the phenomenological correlate of the algebraic distinction between zero emergence and positive emergence.
\end{remark}

\subsection{The Grain of the Voice}

In his essay ``The Grain of the Voice'' (1972), Barthes distinguished
between the \emph{pheno-song} (the communicative, expressive, dramatic
dimension of singing) and the \emph{geno-song} (the material, bodily,
grain-like dimension --- the physicality of the voice, the tongue, the
glottis). The pheno-song is signification; the geno-song is
\emph{significance} --- meaning that exceeds signification, that
resides in the body of the sign rather than in its conceptual content.

\begin{definition}[Pheno-Resonance and Geno-Resonance]
\label{def:pheno-geno}
The \emph{pheno-resonance} of a sign $\sigma = \compose(s_r, s_d)$
with context $c$ is the component attributable to the signified:
$\rs(s_d, c)$. The \emph{geno-resonance} is the component
attributable to the signifier: $\rs(s_r, c)$.
\end{definition}

\begin{theorem}[The Grain is the Remainder]
\label{thm:grain-remainder}
The ``grain'' of a sign --- its irreducible material surplus --- is:
\[
  \texttt{grain}(\sigma, c) = \rs(\sigma, c) - \rs(s_d, c) - \rs(s_r, c)
  = \emerge(s_r, s_d, c).
\]
The grain is the emergence of the sign.
\end{theorem}

\begin{proof}
By the semiotic decomposition theorem
(Theorem~\ref{thm:semiotic-decomposition}), $\rs(\sigma, c) =
\rs(s_r, c) + \rs(s_d, c) + \emerge(s_r, s_d, c)$. Subtracting the
pheno and geno components leaves the emergence. See
\texttt{grain\_is\_emergence}.
\end{proof}

\begin{interpretation}[The Body of the Sign]
This theorem identifies Barthes's ``grain of the voice'' with the
emergence of the sign. The grain is not the signifier (the physical
sound) nor the signified (the conceptual content) but the
\emph{surplus} generated by their coupling --- the irreducible
something-more that makes this particular voice singing this
particular song an experience that exceeds both the music and the
words. The grain is what distinguishes a \emph{performance} from a
\emph{rendition}, art from communication, the writerly from the
readerly. It is, in the deepest sense, the locus of aesthetic value
in the $\IS$.
\end{interpretation}

\subsection{The Fashion System: Clothing as Language}

In \textit{The Fashion System} (1967), Barthes applied semiological
analysis to fashion magazines, treating clothing as a language in the
strict Saussurean sense. Each garment is a sign composed of a material
signifier (the physical fabric, the cut, the color) and a signified
(the meaning attributed by the fashion system: ``sporty,'' ``elegant,''
``casual''). The fashion system is a \emph{code} that maps garments to
meanings, and this code changes with the seasons---not because garments
physically change but because the system of differences is rearranged.

\begin{theorem}[Fashion as Coded Composition]
\label{thm:fashion-composition}
Let $g = \compose(m, f)$ where $m$ is the material garment and $f$ is the fashion code. The ``meaning'' of the garment in the fashion system is not the material alone nor the code alone but the emergence of their composition:
\[
  \texttt{fashionMeaning}(g, c) = \emerge(m, f, c).
\]
When the code changes (from $f$ to $f'$), the same material acquires a new meaning: $\emerge(m, f', c) \neq \emerge(m, f, c)$ in general.
\end{theorem}

\begin{proof}
By the semiotic decomposition theorem, $\rs(g, c) = \rs(m, c) + \rs(f, c) + \emerge(m, f, c)$. The fashion meaning is the emergence term, which depends on the specific pairing of material and code. Changing $f$ to $f'$ changes the emergence while leaving $\rs(m, c)$ invariant. See \texttt{fashion\_coded\_composition}.
\end{proof}

\begin{interpretation}[The Arbitrariness of Fashion]
Barthes argued that fashion reveals the arbitrary nature of the sign in its purest form: the ``meaning'' of a garment (its fashionability, its social register, its connotative value) is entirely determined by the code, not by the material. A miniskirt ``means'' liberation in 1966 and retro nostalgia in 1996---not because the fabric has changed but because the fashion code has rotated. In the $\IS$, this rotation is a change in the composition partner: the same material $m$ is composed with different codes $f, f'$, producing different emergences. The material is constant; the meaning is variable. This is Saussure's arbitrariness thesis in its most radical form, applied not to language but to the entire system of cultural artifacts. The fashion system is the semiotic machine that manufactures new meanings from old materials, and the emergence function $\emerge$ is the algebraic measure of this manufacture.
\end{interpretation}

\begin{remark}[Barthes's Progression: From System to Affect]
Barthes's intellectual trajectory can be traced through the $\IS$. In \textit{Mythologies} (1957) and \textit{The Fashion System} (1967), he was a structuralist: his concern was with the \emph{system} of signs, the code that organizes meaning, the resonance function $\rs$ that maps signs to their cultural positions. In \textit{S/Z} (1970) and \textit{The Pleasure of the Text} (1973), he turned to the \emph{surplus}: the emergence $\emerge$ that exceeds the code, the writerly dimension that resists systematization, the jouissance that disrupts the comfortable pleasures of recognition. In \textit{Camera Lucida} (1980), he arrived at the \emph{singular}: the punctum that pierces through all codes, the irreducible affect that belongs to no system but to a specific encounter between a specific viewer and a specific image. In the $\IS$, this trajectory is a movement from $\rs$ to $\emerge$ to the conditions under which $\emerge$ attains its maximum---from the study of the mean to the study of the extremes, from the center of the distribution to its tails. Barthes's genius was to see that the most interesting semiotic phenomena occur not where the system functions smoothly but where it breaks down---where the emergence is not merely positive but overwhelming.
\end{remark}

In \textit{Camera Lucida} (1980), his last book, Barthes introduced
another famous dyad: the \emph{studium} (the cultural, linguistic,
politically interested reading of a photograph) and the \emph{punctum}
(the detail that ``pricks,'' ``wounds,'' that pierces through the
studium with an intensity that exceeds cultural coding). The studium
is the readerly dimension of the photograph; the punctum is its
writerly surplus.

\begin{definition}[Studium]
\label{def:studium}
The \emph{studium} of a photograph $p$ viewed by reader $\rho$ is the
culturally coded resonance:
$\texttt{studium}(\rho, p) = \rs(\rho, p)$.
\end{definition}

\begin{definition}[Punctum]
\label{def:punctum}
The \emph{punctum} of a photograph $p$ viewed by reader $\rho$ in
context $c$ is the emergence:
$\texttt{punctum}(\rho, p, c) = \emerge(\rho, p, c)$.
\end{definition}

\begin{theorem}[Punctum is the Wound of Emergence]
\label{thm:punctum-emergence}
The punctum is positive if and only if the photograph-reader encounter
generates interpretive surplus:
$\texttt{punctum}(\rho, p, c) > 0 \iff \emerge(\rho, p, c) > 0$.
\end{theorem}

\begin{proof}
By definition. See \texttt{punctum\_is\_emergence}.
\end{proof}

\begin{theorem}[No Punctum Without a Reader]
\label{thm:punctum-requires-reader}
$\texttt{punctum}(\void, p, c) = 0$ for all photographs $p$ and
contexts $c$.
\end{theorem}

\begin{proof}
$\emerge(\void, p, c) = 0$ by void naturalization. See
\texttt{punctum\_void\_reader}.
\end{proof}

\begin{interpretation}[The Privacy of the Wound]
Barthes insisted that the punctum is irreducibly personal: what
pierces one viewer leaves another unmoved. This subjectivity is
captured by the reader-dependence of emergence: $\emerge(\rho, p, c)$
depends on the particular reader $\rho$. A different reader generates
a different emergence, and what constitutes a punctum for one person
may be mere studium for another. The theorem that a void reader
generates no punctum confirms that the punctum requires an actual,
embodied, historically situated reader --- not a disembodied
``subject'' but a person with a specific resonance profile, specific
memories, specific vulnerabilities. The punctum is where the
universal structure of the $\IS$ meets the irreducible singularity
of the individual.
\end{interpretation}

\begin{theorem}[Studium--Punctum Orthogonality]
\label{thm:studium-punctum-orthogonal}
For a photograph $p$ and reader $\rho$, the total resonance of the reader-photograph encounter decomposes as:
\[
  \rs(\compose(\rho, p), c) = \underbrace{\rs(\rho, c)}_{\text{reader contribution}} + \underbrace{\rs(p, c)}_{\text{studium}} + \underbrace{\emerge(\rho, p, c)}_{\text{punctum}}.
\]
The studium and punctum are ``orthogonal'' in the sense that they arise from structurally independent components: the studium from the photograph's resonance profile, the punctum from the emergence of the reader-photograph composition.
\end{theorem}

\begin{proof}
This is a direct application of the semiotic decomposition theorem (Theorem~\ref{thm:semiotic-decomposition}) with $a = \rho$, $b = p$. The reader contribution $\rs(\rho, c)$ captures what the reader brings independently of this photograph; the studium $\rs(p, c)$ captures what the photograph conveys independently of this reader; and the punctum $\emerge(\rho, p, c)$ captures the irreducible surplus of their specific encounter.
\end{proof}

\begin{interpretation}[Barthes's Late Turn: From Structure to Affect]
The studium-punctum orthogonality theorem illuminates Barthes's intellectual trajectory. In his structuralist period (\textit{Mythologies}, \textit{S/Z}, \textit{The Fashion System}), Barthes was concerned with the studium---the culturally coded dimension of signs, the system of differences, the play of connotation and denotation. In his late period (\textit{The Pleasure of the Text}, \textit{Roland Barthes by Roland Barthes}, \textit{Camera Lucida}), he turned increasingly to the punctum---the affective, personal, emergence-dependent dimension that escapes cultural coding. The $\IS$ shows that these are not contradictory interests but \emph{complementary components} of the same semiotic structure: the studium lives in the resonance function $\rs$, while the punctum lives in the emergence function $\emerge$. Barthes's career was not a rejection of structuralism but a \emph{completion} of it: having mapped the resonance landscape (the system of differences), he turned to explore the emergence landscape (the system of surpluses). Together, they constitute the full semiotic theory of the $\IS$.
\end{interpretation}

\begin{remark}[Cross-Reference: Photography and the Semiosphere]
Barthes's semiotics of photography connects to Lotman's theory of the semiosphere (Chapter~5) in a precise way. The studium of a photograph is its position within the semiosphere: its resonance with the culturally dominant codes that constitute the center. The punctum is the photograph's position at the boundary: the point where the private experience of the viewer collides with the public code of the image, generating the emergence that Lotman calls an ``explosion.'' In this light, every photograph is a semiospheric event: it simultaneously occupies a position in the cultural center (studium) and punctures the boundary between public meaning and private experience (punctum). The \textit{Camera Lucida} is thus not merely a book about photography but a theory of the semiosphere's most intimate boundary---the boundary between the social and the personal, between culture and affect, between langue and the untranslatable singularity of individual experience.
\end{remark}



% ================================================================
% CHAPTER 4: ECO'S OPEN WORK AND THE LIMITS OF INTERPRETATION
% ================================================================
\chapter{Eco's Open Work and the Limits of Interpretation}
\label{ch:eco}

\epigraph{A text is a machine conceived for eliciting interpretations.}{Umberto Eco, \textit{The Role of the Reader}, 1979}

\section{The Dialectic of Openness and Closure}

Umberto Eco's contribution to semiotics is distinguished by its
commitment to a productive tension: between the openness of the text
(its capacity to generate multiple interpretations) and the constraints
that limit interpretation (the text's own structure, the codes it
deploys, the encyclopedia it presupposes). Against Barthes's tendency
to celebrate unlimited plurality, Eco insists that some interpretations
are better than others, that texts guide their readers, and that the
freedom of interpretation is bounded by the properties of the text
itself.

This dialectic --- openness constrained by structure --- maps naturally
onto the $\IS$'s interplay of emergence and resonance. Emergence
measures the openness of a composition (the new meaning generated
beyond the parts), while the axioms of the $\IS$ constrain the
magnitude of emergence (it cannot exceed certain bounds determined by
the self-resonances of the components). Eco's semiotics is, in this
sense, the semiotics \emph{native} to the $\IS$: it takes both the
creative and the constraining aspects of the formal structure seriously.

\section{Open and Closed Texts}

\subsection{Formal Definitions}

\begin{definition}[Open Text]
\label{def:eco-open}
A text $t = \compose(a, b)$ is \emph{open} (in Eco's sense) if there
exists a context $c$ such that the emergence is strictly positive:
\[
  \texttt{ecoOpen}(a, b) \iff \exists\, c,\; \emerge(a, b, c) > 0.
\]
An open text generates meaning beyond the sum of its parts: it invites
interpretation, rewards rereading, and produces different meanings in
different contexts.
\end{definition}

\begin{definition}[Closed Text]
\label{def:eco-closed}
A text $t = \compose(a, b)$ is \emph{closed} if for all contexts $c$,
the emergence is non-positive:
\[
  \texttt{ecoClosed}(a, b) \iff \forall\, c,\;
  \emerge(a, b, c) \leq 0.
\]
A closed text generates no surplus meaning: it says what it says and
nothing more. Its meaning is exhausted by its explicit content.
\end{definition}

\begin{lstlisting}[caption={Open and closed texts in the IdeaticSpace}]
def ecoOpen (a b : I) : Prop := ∃ c : I, emergence a b c > 0

def ecoClosed (a b : I) : Prop := ∀ c : I, emergence a b c ≤ 0
\end{lstlisting}

\begin{interpretation}[The Phenomenology of Open Texts]
Eco's examples of open works include Joyce's \textit{Finnegans Wake},
Mallarm\'e's \textit{Livre}, Stockhausen's \textit{Klavierst\"uck XI},
and Calder's mobiles. What these works share is an excess of meaning
--- they generate more interpretive possibilities than any single
reading can exhaust. In the $\IS$, this excess is precisely the
positive emergence: the composition of the text's elements produces
resonances that transcend the elements themselves. A closed text, by
contrast, is a manual, a tax form, a traffic sign --- its meaning
is determined, its emergence negligible. The formal definitions capture
this intuition exactly: openness is positive emergence, closure is
the absence of positive emergence.
\end{interpretation}

\subsection{The Kristeva--Eco Equivalence}

Julia Kristeva's distinction between the ``semiotic'' and the
``symbolic'' modalities of language bears a deep connection to Eco's
open/closed distinction.

\begin{theorem}[Kristeva's Semiotic Equals Eco's Open]
\label{thm:kristeva-eco}
$\texttt{kristevaSemiotic}(a, b) \iff \texttt{ecoOpen}(a, b)$.
\end{theorem}

\begin{proof}
Kristeva's ``semiotic'' modality is defined (in our formalization) as
the presence of positive emergence --- the eruption of pre-linguistic
drives and rhythms into the symbolic order. This is exactly the
condition for an open text. See
\texttt{kristevaSemiotic\_iff\_ecoOpen}.
\end{proof}

\begin{interpretation}[The Semiotic as the Open]
This theorem reveals a deep structural identity between two seemingly
different theoretical frameworks. Kristeva's ``semiotic'' (drawn from
psychoanalysis and phenomenology) and Eco's ``open'' (drawn from
aesthetics and information theory) are the same concept in different
dress. Both name the surplus of meaning that arises when signs combine
in ways that exceed conventional expectations. The $\IS$ shows that
this surplus has a single algebraic source: the emergence cocycle
$\emerge$.
\end{interpretation}

\begin{theorem}[Openness Monotonicity Under Composition]
\label{thm:openness-mono}
If a text $t$ is open (admits multiple non-equivalent valid interpretations), then the composed text $t \compose a$ is at least as open, in the sense that composition with any idea $a$ cannot reduce the number of valid resonance profiles below those of $t$ alone.
\end{theorem}

\begin{proof}
By axiom E4, $w(t \compose a) \geq w(t)$, so the composed text has at least as much self-resonance as $t$ alone. The richer resonance structure of $t \compose a$ means that for any observer $c$, the resonance $\rs(t \compose a, c)$ draws on the contributions of both $t$ and $a$, plus the emergence $\emerge(t, a, c)$. The additional degree of freedom introduced by the emergence term means that the composed text can resonate differently with different observers in ways that $t$ alone could not. Thus composition preserves or increases interpretive openness.
\end{proof}

\begin{interpretation}[Eco's Open Work and the Poetics of Ambiguity]
In \textit{The Open Work} (1962), Eco argued that certain artworks---Mallarm\'e's \textit{Livre}, Stockhausen's \textit{Klavierst\"uck XI}, Calder's mobiles---are deliberately designed to admit multiple valid ``performances'' or interpretations. The theorem above formalizes a key consequence of this openness: adding material to an open work never closes it. Composition enriches, and enrichment means more resonance, more potential for diverse responses. This is why open works tend to generate ever-expanding critical literatures: each new interpretation (a new ``composition'' of the text with a critical framework) adds to the text's resonance without exhausting it. Eco's metaphor of the ``opera aperta'' as a field of possibilities is, in our framework, a statement about the geometry of the text's resonance profile: an open text occupies a region of $\IS$ that is convex and unbounded in the direction of new compositions.
\end{interpretation}

\begin{remark}[The Role of the Reader]
Eco's \textit{The Role of the Reader} (1979) developed the concept of the \emph{Model Reader}: the ideal reader ``foreseen'' by the text, whose interpretive competence matches the text's demands. In the $\IS$, the Model Reader is an observer $c^*$ that maximizes the resonance $\rs(t, c^*)$ subject to the constraint that $c^*$ belongs to the set of ``competent'' observers (those who possess the relevant cultural codes). The actual reader may deviate from the Model Reader, producing either \emph{underinterpretation} ($\rs(t, c) < \rs(t, c^*)$, the reader misses the text's richness) or \emph{overinterpretation} ($\rs(t, c) > \rs(t, c^*)$, the reader projects meaning that the text does not support). Eco's famous debates with Derrida about the ``limits of interpretation'' are, in our framework, debates about whether overinterpretation is a genuine phenomenon or whether every resonance is equally valid. The $\IS$ axioms do not settle this debate, but they provide a precise language for stating it.
\end{remark}

\begin{theorem}[Model Reader Optimality]
\label{thm:model-reader}
Let $\mathcal{C} \subseteq \IS$ be the set of competent observers for text $t$. If a Model Reader $c^* \in \mathcal{C}$ exists such that $\rs(t, c^*) \geq \rs(t, c)$ for all $c \in \mathcal{C}$, then for any reader $\rho$:
\[
  \rs(t, \rho) \leq w(t) \cdot w(\rho).
\]
That is, even the most generous reading of a text is bounded by the geometric mean of the text's weight and the reader's weight.
\end{theorem}

\begin{proof}
This follows directly from axiom E3 (Emergence Boundedness), which provides a Cauchy--Schwarz-like inequality: $\rs(a, b)^2 \leq w(a) \cdot w(b)$ for all $a, b \in \IS$. Taking the square root and noting that resonance is bounded by $\sqrt{w(t) \cdot w(\rho)}$, we obtain the result. The Lean proof is \texttt{model\_reader\_optimality}.
\end{proof}

\begin{interpretation}[The Limits of Interpretation, Formalized]
Eco's \textit{The Limits of Interpretation} (1990) argued against both the hermetic tradition (which claims that a text can mean \emph{anything}) and the positivist tradition (which claims that a text has \emph{one} correct meaning). The Model Reader Optimality theorem provides formal support for Eco's moderate position. On one hand, the resonance $\rs(t, \rho)$ varies with the reader $\rho$, so there is no single ``correct'' meaning. On the other hand, every reading is bounded by $\sqrt{w(t) \cdot w(\rho)}$, so not every reading is equally valid: a reading that exceeds this bound is not interpretation but \emph{overinterpretation}. The bound depends on both the text's weight (its intrinsic complexity) and the reader's weight (their interpretive capacity). A weightier text supports richer readings; a more competent reader extracts more meaning; but neither text nor reader can exceed the bound set by the $\IS$ axioms. This is Eco's ``rights of the text'' made algebraic.
\end{interpretation}

\begin{remark}[Eco and the Encyclopedia as Rhizome]
In \textit{Semiotics and the Philosophy of Language} (1984), Eco proposed that semantic competence is organized not as a dictionary (a tree of definitions) but as an \emph{encyclopedia}: a potentially infinite, labyrinthine network of interconnected entries. Drawing on Deleuze and Guattari's concept of the \emph{rhizome}, Eco described the encyclopedia as a structure without center, without hierarchy, without beginning or end---a network in which every node connects to every other node through some path. In the $\IS$, the encyclopedia model corresponds to the full resonance matrix $(\rs(a, b))_{a, b \in \IS}$: every idea resonates (positively, negatively, or zero) with every other idea, and the web of resonances constitutes the semantic network. The dictionary model, by contrast, corresponds to the much sparser structure of synonymy classes: only ideas that are $\texttt{sameIdea}$ are grouped together. Eco's insight that the encyclopedia is richer than the dictionary is, in our framework, the observation that the resonance structure of the $\IS$ carries more information than its synonymy classes---which is precisely the content of the arbitrariness theorem (Theorem~\ref{thm:arbitrariness}).
\end{remark}

\section{The Model Reader}

\subsection{Dictionary vs.\ Encyclopedia}

Eco distinguished two modes of semantic competence: the
\emph{dictionary} mode, which assigns meanings to signs in isolation
(like a dictionary entry: ``cat = small domesticated feline''), and
the \emph{encyclopedia} mode, which assigns meanings to signs in
context (like an encyclopedia entry: ``cats were sacred in ancient
Egypt, are associated with independence and mystery, feature
prominently in internet culture,'' etc.). The dictionary is finite and
closed; the encyclopedia is potentially infinite and open.

\begin{definition}[Dictionary Entry]
\label{def:dictionary}
A sign $a$ has a \emph{dictionary entry} if its meaning can be
specified without reference to emergence: its resonance profile
$\rs(a, \cdot)$ fully captures its content, and no composition
involving $a$ generates surplus meaning.
\end{definition}

\begin{definition}[Encyclopedia Entry]
\label{def:encyclopedia}
A sign $a$ has an \emph{encyclopedia entry} if its full meaning
requires reference to emergence: there exist $b, c$ such that
$\emerge(a, b, c) \neq 0$.
\end{definition}

\begin{lstlisting}[caption={Dictionary vs.\ encyclopedia entries}]
def dictionaryEntry (a : I) : Prop :=
  ∀ b c : I, emergence a b c = 0

def encyclopediaEntry (a : I) : Prop :=
  ∃ b c : I, emergence a b c ≠ 0
\end{lstlisting}

\begin{theorem}[Encyclopedia Entries Require Emergence]
\label{thm:encyclopedia-emergence}
An idea is an encyclopedia entry if and only if it is not a dictionary
entry --- that is, if it participates in at least one composition with
nonzero emergence.
\end{theorem}

\begin{proof}
$\texttt{encyclopediaEntry}(a) \iff \exists b\,c,\;
\emerge(a,b,c) \neq 0 \iff \neg\forall b\,c,\;
\emerge(a,b,c) = 0 \iff \neg\texttt{dictionaryEntry}(a)$.
\end{proof}

\begin{interpretation}[The Poverty of the Dictionary]
The dictionary/encyclopedia distinction is one of Eco's most productive
ideas. A dictionary tells you what a word ``means'' in isolation; an
encyclopedia tells you what it means in the web of all possible
contexts. The formal distinction between dictionary entries (zero
emergence) and encyclopedia entries (nonzero emergence) shows that
dictionary-like meaning is the \emph{degenerate} case: it is the case
where composition is perfectly additive, where no surplus meaning
arises. Most ideas are encyclopedia entries --- they generate emergence
when composed with other ideas. The dictionary is an abstraction, a
useful fiction that works only for the most sterile, decontextualized
uses of language.
\end{interpretation}

\begin{theorem}[Hegemonic Signs are Closed Texts]
\label{thm:hegemonic-closed}
If a sign $a$ is hegemonic (dominates all other signs in the system),
then the text $\compose(a, b)$ is closed for all $b$.
\end{theorem}

\begin{proof}
Hegemonic signs are left-linear (as we shall prove in
Chapter~5, Theorem~\ref{thm:hegemonic-leftlinear}), which means
$\emerge(a, b, c) = 0$ for all $b, c$. A text with zero emergence
satisfies $\emerge(a,b,c) \leq 0$ for all $c$, hence it is closed.
See \texttt{hegemonic\_ecoClosed}.
\end{proof}

\begin{interpretation}[Power Closes Texts]
This theorem has profound political implications: hegemonic signs ---
those that dominate the semiotic landscape --- produce \emph{closed}
texts. A hegemonic discourse admits no surplus, no alternative reading,
no creative interpretation. It says what it says and nothing more; it
demands submission, not engagement. This is the semiotic mechanism of
ideological closure: power does not merely suppress alternative meanings
but \emph{algebraically prevents their emergence}. The closed text is
the text of authority; the open text is the text of resistance.
\end{interpretation}

\subsection{The Degree of Openness}

\begin{definition}[Degree of Openness]
\label{def:degree-openness}
The \emph{degree of openness} of a text $t = \compose(a, b)$ is:
\[
  \texttt{openness}(a, b) = \sup_{c \in \mathcal{I}} \emerge(a, b, c).
\]
This is the maximum interpretive surplus the text can generate across
all possible contexts.
\end{definition}

\begin{theorem}[Openness is Non-Negative]
\label{thm:openness-nonneg}
$\texttt{openness}(a, b) \geq 0$ for all $a, b$.
\end{theorem}

\begin{proof}
By the compose-enriches axiom (E4), there exists at least one $c$
(namely $c = \compose(a,b)$) for which the emergence is non-negative.
The supremum is therefore at least zero. See
\texttt{openness\_nonneg}.
\end{proof}

\begin{theorem}[Openness of Void Compositions]
\label{thm:openness-void}
$\texttt{openness}(\void, b) = 0$ and $\texttt{openness}(a, \void) = 0$.
\end{theorem}

\begin{proof}
$\emerge(\void, b, c) = 0$ and $\emerge(a, \void, c) = 0$ for all
$c$, so the supremum is zero. See \texttt{openness\_void}.
\end{proof}

\begin{theorem}[Bounded Openness]
\label{thm:openness-bounded}
The degree of openness is bounded above:
\[
  \texttt{openness}(a, b) \leq M \cdot \sqrt{\rs(a, a) \cdot \rs(b, b)}.
\]
\end{theorem}

\begin{proof}
The supremum of $\emerge(a, b, c)$ over all $c$ is bounded by the
emergence bound axiom (E3). See \texttt{openness\_bounded}.
\end{proof}

\begin{interpretation}[The Measure of Textual Richness]
The degree of openness provides a single number that characterizes the
interpretive richness of a text. A text with zero openness (e.g., any
composition involving the void) generates no interpretive surplus: it
is ``what it is'' and nothing more. A text with high openness generates
abundant surplus in the right context: it is a wellspring of
interpretation. But even the most open text has bounded openness ---
it cannot generate infinite meaning. This boundedness is Eco's
insistence on limits, now made quantitative: we can compare texts
by their degrees of openness and ask which is ``more open'' in a
precise mathematical sense.
\end{interpretation}

\begin{theorem}[Openness Additivity Under Independent Composition]
\label{thm:openness-additive}
If texts $t_1 = \compose(a_1, b_1)$ and $t_2 = \compose(a_2, b_2)$ are ``independent'' in the sense that $\rs(t_1, t_2) = 0$, then the openness of their combined text satisfies:
\[
  \texttt{openness}(t_1, t_2) \geq 0.
\]
Moreover, the compound text $\compose(t_1, t_2)$ has weight at least $w(t_1) + w(t_2)$, so the combination of independent texts creates a text whose interpretive capacity exceeds the sum of its parts.
\end{theorem}

\begin{proof}
The weight bound follows from E4 applied to $(t_1, t_2)$: $w(\compose(t_1, t_2)) \geq w(t_1)$. For the enrichment: $\texttt{semioticEnrichment}(t_1, t_2) = w(\compose(t_1, t_2)) - w(t_1) - w(t_2) \geq 0$ by the enrichment non-negativity theorem. See \texttt{openness\_additive}.
\end{proof}

\begin{interpretation}[The Anthology Effect]
The openness additivity theorem explains what might be called the ``anthology effect'': placing two independent texts side by side creates a compound text that is richer than either alone. An anthology of poems, a museum's collection, a film festival's programme---each is a composition of independent texts whose combined openness exceeds the sum of the individual opennesses. This is why curating is a creative act: the curator does not merely collect but \emph{composes}, and the composition generates emergence. The specific juxtaposition of works---which poem follows which, which painting hangs beside which---matters because it determines the emergence of the compound text. A well-curated anthology is a text whose emergence is high; a badly curated one is merely additive.

This connects to Eco's concept of the ``intentio operis''---the text's own intention, distinct from both the author's intention and the reader's intention. The intentio operis of an anthology is not the sum of the intentions of its individual texts but a genuinely new intention that emerges from their composition. The curator's art is to arrange texts so that this emergent intention is rich, coherent, and illuminating.
\end{interpretation}

\begin{remark}[Eco's Peircean Heritage]
Eco was one of the few semioticians to take Peirce's theory as seriously as Saussure's. His concept of ``unlimited semiosis'' draws directly from Peirce's doctrine of the interpretant, and his concept of the ``encyclopedia'' generalizes Peirce's notion of the ``ground'' of a sign relation. In the $\IS$, the connection is precise: Eco's encyclopedia is the full resonance matrix $(\rs(a, b))_{a, b \in \IS}$, while Peirce's ground is the set of test ideas with respect to which a specific sign relation is defined (Definition~\ref{def:ground}). The encyclopedia contains all grounds; the ground is a local slice of the encyclopedia. Eco's innovation was to insist that semantic competence requires the \emph{whole} encyclopedia---that you cannot understand a sign without understanding its position in the entire network of signs. The $\IS$ vindicates this holism: the meaning of an idea is its resonance profile, and the resonance profile is defined with respect to \emph{all} other ideas. Local knowledge is always partial knowledge.
\end{remark}

\begin{lstlisting}[caption={The spectrum of textual openness}]
noncomputable def openness (a b : I) : ℝ :=
  sSup {x : ℝ | ∃ c : I, emergence a b c = x}

theorem openness_nonneg (a b : I) : openness a b ≥ 0 := ...
theorem openness_void_left (b : I) : openness void b = 0 := ...
theorem openness_void_right (a : I) : openness a void = 0 := ...
\end{lstlisting}

\section{Interpretation, Use, and the Rights of the Text}

\subsection{Use vs.\ Interpretation}

In \textit{The Limits of Interpretation} (1990), Eco drew a crucial
distinction between \emph{interpretation} (recovering the intentio
operis, respecting the text's structure) and \emph{use} (deploying
the text for one's own purposes, regardless of its structure). A
Marxist reading of \textit{Hamlet} that illuminates the class dynamics
of Elsinore is interpretation; using the pages of \textit{Hamlet} as
wallpaper is use. Both are legitimate activities, but they are
categorically different.

\begin{definition}[Interpretation vs.\ Use]
\label{def:interpretation-use}
An encounter between reader $\rho$ and text $t$ is \emph{interpretation}
if the emergence respects the text's resonance profile:
\[
  \texttt{interpretation}(\rho, t) \iff
  \text{sign}(\emerge(\rho, t, c)) = \text{sign}(\rs(t, c))
  \quad \text{for most } c.
\]
The encounter is \emph{use} if the emergence is independent of the
text's resonance profile.
\end{definition}

\begin{interpretation}[The Ethics of Reading]
Eco's distinction between interpretation and use is fundamentally
\emph{ethical}: it concerns the reader's obligation to the text.
Interpretation acknowledges that the text has its own structure, its
own resonance profile, its own intentio operis, and seeks to generate
emergence that is consonant with that structure. Use ignores the
text's structure and treats it as raw material for the reader's own
purposes. Both are possible in the $\IS$ (the axioms do not
distinguish between them), but the distinction matters for any
normative theory of reading. Eco's semiotics is, in this sense, an
ethics of interpretation: it asks not just ``what can the text mean?''
but ``what \emph{should} the text be taken to mean?''
\end{interpretation}

\section{Interpretive Openness and Its Limits}

\subsection{The Conditions of Interpretation}

Eco insisted that interpretation requires two things: a \emph{text}
and a \emph{reader}. Neither alone suffices. A text without a reader
is a mere sequence of marks; a reader without a text is a mind without
an object. The encounter between reader and text is what generates
meaning --- and this encounter is modeled, in the $\IS$, by emergence.

\begin{theorem}[Interpretive Openness Requires Both Reader and Text]
\label{thm:openness-requires-both}
Interpretive openness vanishes if either the reader or the text is
void:
\begin{align}
  \emerge(\void, t, c) &= 0 \quad \text{for all } t, c, \\
  \emerge(r, \void, c) &= 0 \quad \text{for all } r, c.
\end{align}
\end{theorem}

\begin{proof}
$\emerge(\void, t, c) = \rs(\compose(\void, t), c) - \rs(\void, c)
- \rs(t, c) = \rs(t, c) - 0 - \rs(t, c) = 0$, using the left
identity axiom and the void resonance axiom. Similarly for the right
case. See \texttt{interpretiveOpenness\_void\_reader} and
\texttt{interpretiveOpenness\_void\_text}.
\end{proof}

\begin{interpretation}[Against Solipsism and Objectivism]
This theorem navigates between two extremes. The solipsist says meaning
is in the reader; the objectivist says meaning is in the text. Eco
says meaning is in the \emph{encounter}, and the algebra agrees: the
emergence function takes two arguments, and it vanishes when either is
void. There is no meaning without a reader (the void reader generates
no emergence) and no meaning without a text (the void text generates no
emergence). Meaning is irreducibly relational --- it is a property not
of signs or of minds but of the transaction between them.
\end{interpretation}

\subsection{Overinterpretation}

Eco was equally concerned with the \emph{limits} of interpretation.
In his Tanner Lectures (published as \textit{Interpretation and
Overinterpretation}, 1992), he argued against the ``hermetic drift''
--- the tendency of certain interpretive traditions to find meaning
everywhere, to see connections where none exist, to overinterpret the
text into saying whatever the interpreter wishes. The $\IS$ provides a
formal criterion for overinterpretation.

\begin{theorem}[Overinterpretation is Bounded]
\label{thm:overinterpretation-bounded}
The interpretive surplus (emergence) of any reading is bounded above
by the self-resonances of the reader and the text:
\[
  |\emerge(r, t, c)| \leq M \cdot \sqrt{\rs(r, r) \cdot \rs(t, t)}
\]
for some universal constant $M > 0$.
\end{theorem}

\begin{proof}
This is a direct application of the emergence bounded axiom (E3).
See \texttt{overinterpretation\_bounded}.
\end{proof}

\begin{interpretation}[The Economy of Interpretation]
The overinterpretation bound is Eco's criterion of ``interpretive
economy'' made algebraic. You cannot extract more meaning from a text
than the text (and the reader) can support. A trivial text (small
self-resonance) cannot sustain deep interpretation; a shallow reader
(small self-resonance) cannot produce deep readings. The bound is
proportional to the geometric mean of the reader's and text's
``magnitudes,'' which captures the intuition that rich readings require
both rich texts and rich readers. Overinterpretation occurs when the
claimed emergence exceeds this bound --- when the interpreter attributes
more surplus meaning to the encounter than the structure of the $\IS$
permits.
\end{interpretation}

\begin{remark}[Eco's Three Intentions]
Eco distinguished three types of textual intention: the \emph{intentio auctoris} (what the author meant), the \emph{intentio lectoris} (what the reader wants the text to mean), and the \emph{intentio operis} (what the text itself means, given its structure and codes). In the $\IS$, these three intentions correspond to three different projections of the same semiotic object. The intentio auctoris is the resonance profile of the author: $(\rs(\alpha, c))_c$, which encodes the author's communicative intent. The intentio lectoris is the resonance profile of the reader: $(\rs(\rho, c))_c$, which encodes the reader's interpretive desires. The intentio operis is the resonance profile of the text: $(\rs(t, c))_c$, which encodes the text's structural affordances. Eco's position is that legitimate interpretation must respect the intentio operis---the text's own resonance profile---even when this conflicts with the intentio lectoris. The overinterpretation bound theorem provides formal support: the emergence generated by a reader-text encounter is bounded by the \emph{text's} self-resonance (among other factors), not merely by the reader's desire for meaning. The text resists arbitrary readings, and the algebra quantifies this resistance.
\end{remark}

\begin{theorem}[The Intentio Operis as Invariant]
\label{thm:intentio-operis}
The intentio operis of a text $t$---its resonance profile $(\rs(t, c))_{c \in \IS}$---is invariant under changes of reader. That is, for any two readers $\rho_1, \rho_2$:
\[
  \rs(t, c) = \rs(t, c) \quad \text{for all } c.
\]
The text's intrinsic resonance profile does not change when a different reader encounters it. What changes is the emergence: $\emerge(\rho_1, t, c) \neq \emerge(\rho_2, t, c)$ in general.
\end{theorem}

\begin{proof}
The resonance function $\rs$ is a fixed property of the $\IS$, not a reader-dependent quantity. The text $t$ has a unique resonance profile independent of who reads it. The emergence, however, is a function of both reader and text, and thus varies with the reader. This is the algebraic separation of the objective (the text's resonance profile) from the subjective (the reader's emergence with the text).
\end{proof}

\begin{interpretation}[Objectivity and Subjectivity in Interpretation]
The intentio operis invariance theorem provides a precise formulation of the old hermeneutic problem: what is ``in the text'' and what is ``in the reader''? The answer: the resonance profile $(\rs(t, c))_c$ is in the text (it is reader-invariant), while the emergence $\emerge(\rho, t, c)$ is in the encounter (it is reader-dependent). This distinction maps onto Eco's own resolution of the objectivity/subjectivity debate: the text has objective properties (its codes, its structures, its resonance profile) that constrain interpretation, but the actual meaning generated in any reading is subjective (it depends on the reader's resonance profile and the resulting emergence). Neither pure objectivism (meaning is in the text alone) nor pure subjectivism (meaning is in the reader alone) is correct; meaning is in the \emph{composition} of reader and text, and this composition has both objective constraints (the resonance profiles) and subjective contributions (the emergence).
\end{interpretation}

\begin{theorem}[Encyclopedia Enriches]
\label{thm:encyclopedia-enriches}
Encyclopedic reading (reading that engages with the full contextual
web of a sign's meanings) produces at least as much interpretive
openness as dictionary reading:
\[
  \texttt{encyclopediaEntry}(a) \implies \exists b\, c,\;
  \emerge(a, b, c) > 0.
\]
\end{theorem}

\begin{proof}
If $a$ is an encyclopedia entry, then $\exists b\,c,\;
\emerge(a,b,c) \neq 0$. If this emergence is positive for some
$b, c$, we are done. If it is negative, the compose-enriches axiom
(E4) guarantees that \emph{some} composition has positive emergence
(otherwise the space would degenerate). See
\texttt{encyclopedia\_enriches}.
\end{proof}

\begin{interpretation}[The Reward of Contextual Reading]
This theorem vindicates Eco's preference for encyclopedic over
dictionary reading. A dictionary reader encounters a sign and assigns
it a fixed meaning; an encyclopedic reader encounters a sign and
explores its contextual connections, generating emergence through
composition. The theorem says that this exploration is always
\emph{rewarding}: encyclopedia entries always participate in at least
one positively emergent composition. There is always more meaning to
find --- but (by the overinterpretation bound) never infinitely more.
\end{interpretation}

\begin{remark}[The Encyclopedia as Rhizome]
Eco borrowed the botanical metaphor of the \emph{rhizome} from Deleuze and Guattari to describe the structure of encyclopedic knowledge. Unlike a tree (which has a single root, a hierarchical structure, and terminal leaves), a rhizome has no center, no hierarchy, and no boundaries: any point can connect to any other point. In the $\IS$, the rhizomatic structure of the encyclopedia corresponds to the multi-dimensionality of the resonance function: each idea resonates with every other idea along multiple dimensions, and there is no privileged direction or organizing principle. The resonance function $\rs(a, b)$ is defined for all pairs $(a, b)$, not just for pairs connected by a tree-like hierarchy. This universality of resonance is the formal expression of the rhizomatic principle: in the semiosphere, everything is potentially connected to everything else, and the encyclopedia is the map of these connections.
\end{remark}

\begin{theorem}[Eco's Openness Additivity]
\label{thm:openness-additivity}
If text $t_1$ is open (there exist $\rho_1, c_1$ with $\emerge(\rho_1, t_1, c_1) > 0$) and text $t_2$ is open, then the composition $\compose(t_1, t_2)$ is also open:
\[
  \exists \rho, c, \quad \emerge(\rho, \compose(t_1, t_2), c) > 0.
\]
Openness is preserved under textual composition.
\end{theorem}

\begin{proof}
Since $t_1$ is open, $\emerge(\rho_1, t_1, c_1) > 0$ for some $\rho_1, c_1$. The composition $\compose(t_1, t_2)$ has weight $w(\compose(t_1, t_2)) \geq w(t_1) + w(t_2)$ by E4. Since $w(\compose(t_1, t_2)) > 0$ and the text has non-trivial resonance structure, there exists a reader $\rho$ that generates positive emergence with the composition. The formal argument uses the cocycle identity: the emergence of $\rho$ with $\compose(t_1, t_2)$ relates to the emergences of $\rho$ with $t_1$ and $t_2$ individually via the master cocycle, and the positive contributions from $t_1$'s openness propagate through the composition. See \texttt{openness\_additivity}.
\end{proof}

\begin{interpretation}[The Anthology Effect]
The openness additivity theorem explains the ``anthology effect'': a collection of open texts, composed together (in an anthology, a museum exhibition, a concert program), is itself open. The openness of the individual texts compounds through composition, and the resulting collection may generate emergences that none of the individual texts could produce alone. This is why curating---the art of composing texts together---is itself a creative act: the curator's composition generates emergence beyond what the individual works contribute. The emergence of the anthology is not merely the sum of the individual emergences but includes the cross-emergences between texts---the resonances between, say, a Vermeer and a de Hooch hung side by side, or between a Beethoven sonata and a Schubert lied performed in sequence. These cross-emergences are the formal content of the curator's art.
\end{interpretation}

Eco proposed the concept of the \emph{Model Reader}: the ideal reader
presupposed by the text, the reader who possesses the competences
necessary to actualize the text's meaning. The Model Reader is not a
real person but a textual strategy --- a set of instructions encoded in
the text for its own interpretation.

\begin{example}[The Model Reader in the $\IS$]
Consider a text $t = \compose(a, b)$ and a Model Reader $\rho^*$ such
that $\emerge(\rho^*, t, c)$ is maximized. The Model Reader is the
reader who extracts the maximum emergence --- the maximum interpretive
surplus --- from the text. By the overinterpretation bound, this
maximum is finite and depends on $\rs(\rho^*, \rho^*)$ and $\rs(t, t)$.
A different text requires a different Model Reader; a reader who is
``model'' for one text may be a poor reader of another. This captures
Eco's insight that every text constructs its own ideal reader: a
detective novel constructs a reader who enjoys puzzle-solving, a
philosophical treatise constructs a reader versed in the relevant
tradition, a love poem constructs a reader who is emotionally
available.
\end{example}

\begin{remark}[The Intentio Operis]
Eco distinguished three kinds of intention: the \emph{intentio auctoris}
(what the author meant), the \emph{intentio lectoris} (what the reader
wants the text to mean), and the \emph{intentio operis} (what the text
itself means). In the $\IS$, the intentio operis corresponds to the
resonance profile of the text: $\rs(t, \cdot)$. This profile is
independent of any particular reader --- it is a property of the text
alone. But the \emph{actualization} of this profile --- the generation
of interpretive surplus --- requires a reader, and the surplus depends
on the reader's own resonance profile. Eco's middle position between
Barthesian anarchy and authorial authority is thus algebraically
grounded: the text has a definite structure (its resonance profile),
but this structure underdetermines the meaning (which depends also on
the reader's emergence).
\end{remark}

\begin{theorem}[Eco's Principle of Charity]
\label{thm:principle-charity}
Given a text $t$ and two candidate readings $\rho_1, \rho_2$, the ``more charitable'' reading is the one whose emergence aligns more closely with the text's own resonance profile:
\[
  \rho_1 \text{ is more charitable than } \rho_2 \iff \sum_{c} |\emerge(\rho_1, t, c) - \rs(t, c)| < \sum_{c} |\emerge(\rho_2, t, c) - \rs(t, c)|.
\]
\end{theorem}

\begin{proof}
The principle of charity, as adapted from analytic philosophy to semiotics by Eco, requires that we attribute to the text the richest meaning consistent with its structure. In our framework, the text's structure is encoded in $\rs(t, c)$, and the reader's contribution is encoded in $\emerge(\rho, t, c)$. A charitable reading is one whose emergence ``follows'' the text's resonance---the reader generates surplus in the same dimensions where the text has strength, rather than projecting surplus onto dimensions where the text has none. This is not a theorem derivable from the $\IS$ axioms alone but a \emph{normative principle} expressed in the language of the $\IS$.
\end{proof}

\begin{interpretation}[Charity as Resonance Alignment]
The principle of charity theorem provides a formal criterion for distinguishing ``good'' readings from ``bad'' ones without appealing to authorial intention. A good reading is one that resonates with the text---one whose emergence profile aligns with the text's own resonance profile. A bad reading is one that projects its own concerns onto the text, generating emergence in dimensions where the text has no resonance. This does not mean that creative, surprising readings are bad: a reading may generate emergence in dimensions where the text's resonance is subtle but nonzero, revealing hidden depths that other readers have missed. What it does mean is that a reading that generates emergence in dimensions where the text has \emph{zero} resonance is, by the charity criterion, not an interpretation but a \emph{use} of the text. This connects to Eco's distinction between interpretation and use (Definition~\ref{def:interpretation-use}) and provides a quantitative refinement: the degree of interpretive charity is measurable, and readings can be ranked accordingly.
\end{interpretation}

\section{Eco and the Ideatic Space: A Retrospective}

Eco's semiotics is, among the theoretical frameworks we have examined,
the one most naturally at home in the $\IS$. This is not a coincidence:
Eco was trained in medieval philosophy and aesthetics, and his approach
to semiotics combines rigorous logical analysis with sensitivity to
the richness and ambiguity of cultural phenomena. The $\IS$'s combination
of algebraic structure (axioms, theorems) with interpretive openness
(emergence, resonance) mirrors Eco's own intellectual commitments.

Three features of the $\IS$ are particularly Eco-friendly. First, the
\emph{boundedness} of emergence: Eco always insisted on limits, on the
difference between interpretation and overinterpretation, on the
text's right to resist arbitrary readings. The emergence bound axiom
(E3) is the formal embodiment of this insistence. Second, the
\emph{enrichment} of composition: Eco believed that genuine
interpretation is not a subtraction (peeling away layers to find a
hidden meaning) but an \emph{addition} (composing the text with new
contexts to generate new meanings). The compose-enriches axiom (E4)
captures this belief. Third, the \emph{plurality} of codes: Eco's
encyclopedia model posits an open-ended web of interconnections among
signs, with no privileged dimension or code. The multi-dimensionality
of the resonance space, with its orthogonal subspaces and its emergent
compositions, provides the formal architecture for this web.

\section{Aberrant Decoding and Isotopy}

\subsection{Aberrant Decoding}

Eco introduced the concept of \emph{aberrant decoding} to describe
the situation in which a text is interpreted according to a code
different from the one used to produce it. This is not a pathological
case but the \emph{normal} condition of mass communication: a
television broadcast produced in one cultural context is received in
countless others, each with its own interpretive framework. The
``intended'' meaning is only one of many actualizations.

\begin{definition}[Aberrant Decoding]
\label{def:aberrant-decoding}
An \emph{aberrant decoding} of text $t$ by reader $\rho$ occurs when
the reader's resonance profile is significantly different from the
Model Reader's profile. Formally, reader $\rho$ produces an aberrant
decoding if:
\[
  d(\rho, \rho^*) > \delta
\]
where $\rho^*$ is the Model Reader, $d$ is semiotic distance
(Definition~\ref{def:semiotic-distance}), and $\delta$ is a
threshold of ``acceptable'' interpretive distance.
\end{definition}

\begin{theorem}[Aberrant Decoding Shifts Emergence]
\label{thm:aberrant-shift}
If reader $\rho$ is semiotically distant from the Model Reader
$\rho^*$, then the emergence of their respective readings may differ:
\[
  |\emerge(\rho, t, c) - \emerge(\rho^*, t, c)| \leq
  M' \cdot d(\rho, \rho^*)
\]
for some constant $M'$ depending on $t$ and $c$. The interpretive
surplus is \emph{continuously sensitive} to the reader's position
in semiotic space.
\end{theorem}

\begin{proof}
The emergence is a function of the reader's resonance profile; the
difference in emergence between two readers is bounded by the
difference in their resonance profiles, which is itself bounded by
their semiotic distance. See \texttt{aberrant\_shift\_bounded}.
\end{proof}

\begin{interpretation}[The Universality of Aberrance]
Eco's concept of aberrant decoding reveals a democratic truth about
interpretation: \emph{every} reading is aberrant, because no actual
reader perfectly coincides with the Model Reader. The Model Reader is
a textual strategy, not a real person, and the distance between any
real reader and the Model Reader is always nonzero. But the theorem
shows that the degree of aberrance is bounded and continuous: a small
deviation from the Model Reader produces a small shift in emergence.
Radical aberrance (reading the \textit{Iliad} as a cookbook) produces
large shifts. The $\IS$ provides a \emph{metric} of interpretive
deviation that can distinguish productive aberrance (creative
misreading that generates new meaning) from pathological aberrance
(misreading that generates nonsense).
\end{interpretation}

\subsection{Isotopy}

Eco borrowed the concept of \emph{isotopy} from A.~J.\ Greimas to
describe the phenomenon of semantic coherence: the way in which a text
maintains a consistent ``level'' of meaning across its constituent
elements. An isotopy is a recurrent semic category that makes uniform
reading possible; a text may sustain multiple isotopies simultaneously.

\begin{definition}[Isotopy]
\label{def:isotopy}
An \emph{isotopy} of a text $t$ is a collection of ideas
$\{a_1, \ldots, a_n\} \subset \mathcal{I}$ that are pairwise
positively resonant: $\rs(a_i, a_j) > 0$ for all $i \neq j$.
An isotopy provides a coherent ``reading level'' --- a dimension of
meaning that runs consistently through the text.
\end{definition}

\begin{theorem}[Isotopies Support Coherent Reading]
\label{thm:isotopy-coherent}
If $\{a_1, \ldots, a_n\}$ forms an isotopy and $t = \compose(a_1,
\compose(a_2, \ldots))$ is a text built from these elements, then
the self-resonance of $t$ is at least as large as the sum of the
individual self-resonances:
\[
  \rs(t, t) \geq \sum_i \rs(a_i, a_i).
\]
\end{theorem}

\begin{proof}
By iterated application of the compose-enriches axiom (E4), each
successive composition adds at least the self-resonance of the new
element. The positive pairwise resonances within the isotopy
contribute additional terms. See \texttt{isotopy\_coherent}.
\end{proof}

\begin{interpretation}[Coherence as Enrichment]
An isotopy produces a text whose semantic weight exceeds the sum of
its parts. The positive pairwise resonances within the isotopy ensure
that the elements reinforce each other, creating a coherent reading
experience. When a text sustains multiple isotopies simultaneously
(as in allegory, irony, or double entendre), the different isotopies
may generate different emergences, giving the text a multi-layered
quality. This is the formal counterpart of Barthes's five codes
applied at the level of semantic fields.
\end{interpretation}

\section{Abduction as Interpretive Strategy}

Eco devoted considerable attention to Peirce's theory of abduction,
which he reinterpreted as the fundamental mode of textual interpretation.
When we read a text and form a hypothesis about ``what it means,'' we
are performing an abduction: inferring a general rule (the meaning)
from a particular case (the text) and a result (our experience of
reading it). Eco classified abductions into three types:
\emph{overcoded} (applying a pre-existing rule), \emph{undercoded}
(selecting tentatively among competing rules), and \emph{creative}
(inventing a new rule).

\begin{theorem}[Abductive Hierarchy]
\label{thm:abductive-hierarchy}
The three types of abduction correspond to three levels of emergence:
\begin{enumerate}
\item \emph{Overcoded abduction}: $\emerge(\rho, t, c) \approx 0$
(the rule is already known; no surplus).
\item \emph{Undercoded abduction}: $\emerge(\rho, t, c) > 0$ but
bounded by existing patterns.
\item \emph{Creative abduction}: $\emerge(\rho, t, c) \gg 0$
(a genuinely new rule is generated).
\end{enumerate}
\end{theorem}

\begin{proof}
Overcoded abduction applies a naturalized habit (zero emergence).
Undercoded abduction operates at the boundary between frozen and living
signs (small positive emergence). Creative abduction generates a
genuinely new interpretant with large emergence, constrained only by
the universal bound (E3). See \texttt{abductive\_hierarchy}.
\end{proof}

\begin{interpretation}[The Ecology of Inference]
Eco's classification of abductions is a theory of the \emph{ecology
of inference}: how different interpretive strategies are suited to
different semiotic environments. In a stable, well-coded environment
(a natural language, a genre convention, a legal code), overcoded
abduction suffices: we apply known rules and generate no surprise. In
a less stable environment (a new artistic movement, a cross-cultural
encounter), undercoded abduction is necessary: we try different rules
and see which generates the most coherent reading. In a revolutionary
environment (a genuinely unprecedented text, a paradigm shift),
creative abduction is required: we invent new rules, new codes, new
isotopies. The emergence function measures the ``creativity cost''
of each type of abduction, and the universal bound ensures that even
creative abduction operates within finite limits.
\end{interpretation}

\begin{theorem}[Eco's Possible Worlds]
\label{thm:eco-possible-worlds}
Every text $t$ defines a space of ``possible worlds''---alternative compositions that could have been constructed from the text's components. Formally, if $t = \compose(a, b)$, the possible world generated by an alternative context $b'$ is:
\[
  t' = \compose(a, b')
\]
and the ``narrative distance'' between the actual world $t$ and the possible world $t'$ is:
\[
  d_{\text{narrative}}(t, t') = w(t) + w(t') - 2\rs(t, t').
\]
\end{theorem}

\begin{proof}
The narrative distance inherits the metric properties of the semiotic distance (Theorem~\ref{thm:semiotic-triangle}). Since $t = \compose(a, b)$ and $t' = \compose(a, b')$, the distance depends on how much $b$ and $b'$ differ and on the cross-terms introduced by composition. By the triangle inequality for $\sqrt{d}$, the narrative distance is a well-defined pseudo-metric on the space of possible worlds. See \texttt{eco\_possible\_worlds}.
\end{proof}

\begin{interpretation}[Fiction as Exploration of Possible Worlds]
Eco argued, following Leibniz and analytic philosophy of fiction, that every narrative text constructs a ``possible world''---a state of affairs that differs from the actual world in specific, structured ways. A detective novel constructs a world in which a murder has occurred and a detective investigates; a science fiction novel constructs a world in which physical laws differ from our own. The possible worlds theorem gives this idea a metric structure: the narrative distance between the actual world and the fictional world measures how ``far'' the fiction departs from reality. A realistic novel has small narrative distance (it depicts a world close to ours); a fantasy novel has large narrative distance. But crucially, even the most distant fictional world shares some resonance with the actual world---otherwise it would be unintelligible. This residual resonance is what Eco calls the ``encyclopedia'' that must be shared between text and reader for communication to be possible. The possible worlds framework explains why fiction is not merely escapism but a \emph{cognitive tool}: by exploring nearby possible worlds, fiction allows us to understand the structure of the actual world. The narrative distance provides a quantitative measure of this exploration.
\end{interpretation}

\begin{theorem}[Fictional Enrichment]
\label{thm:fictional-enrichment}
The composition of the actual world $w_0$ with a fictional world $t'$ produces enrichment:
\[
  w(\compose(w_0, t')) \geq w(w_0) + w(t').
\]
Reading fiction enriches the reader's world.
\end{theorem}

\begin{proof}
Direct application of the compose-enriches axiom E4. See \texttt{fictional\_enrichment}.
\end{proof}

\begin{remark}[Eco's Semiotics of the Novel]
Eco's own novels---\textit{The Name of the Rose}, \textit{Foucault's Pendulum}, \textit{The Island of the Day Before}---are deliberate exercises in the semiotics he theorized. Each novel constructs a possible world at a specific narrative distance from the actual world, deploys a specific Model Reader, and tests the limits of interpretation within a fictional framework. \textit{The Name of the Rose} is a medieval detective story that is simultaneously a treatise on semiotics: the detective William of Baskerville is an explicit Peircean abducer, constructing hypotheses from signs and testing them against evidence. The novel is thus a \emph{meta-semiotic} text: a text that enacts the theory of texts. In the $\IS$, this self-referential quality corresponds to the composition of the text with its own interpretive framework---a kind of autocommunication at the level of the text itself, generating the self-enrichment that characterizes all genuine acts of cultural creation.
\end{remark}
the semiotics of texts to the semiotics of cultures. This is the
territory of Yuri Lotman and the Tartu--Moscow school, which we
formalize in Chapter~5.



% ================================================================
% CHAPTER 5: LOTMAN'S SEMIOSPHERE
% ================================================================
\chapter{Lotman's Semiosphere}
\label{ch:lotman}

\epigraph{The semiosphere is the semiotic space necessary for the existence and functioning of languages.}{Yuri M.\ Lotman, \textit{Universe of the Mind}, 1990}

\section{The Semiosphere as Ideatic Space}

Yuri Lotman's concept of the \emph{semiosphere} is one of the most
ambitious in the history of semiotics. By analogy with Vernadsky's
\emph{biosphere} --- the totality of living matter on Earth, forming a
single interconnected system --- Lotman proposed the \emph{semiosphere}:
the totality of sign processes in a culture, forming a single
interconnected system. Just as no organism can exist outside the
biosphere, no sign can function outside the semiosphere. Language does
not arise from individual signs that are subsequently combined; rather,
individual signs presuppose the semiosphere as the condition of their
possibility.

The $\IS$ $(\mathcal{I}, \compose, \void, \rs)$ \emph{is} a
semiosphere. It is a total system of ideas and their relations,
equipped with a composition operation (combining signs), a void
element (the silence that grounds the system), and a resonance
function (the web of affinities and oppositions that constitutes the
system's structure). The axioms A1--A3, R1--R2, and E1--E4 are the
``laws of the semiosphere'' --- the structural constraints that any
system of signs must satisfy. In this chapter, we develop the internal
geography of the semiosphere: its center and periphery, its boundaries,
its dynamics of explosion and gradual change, and its ethical
dimension.

\section{Center, Periphery, and Boundary}

\subsection{The Topology of Cultural Space}

Lotman proposed that the semiosphere is not homogeneous but structured
into a \emph{center} (the dominant, codified, ``grammaticalized'' region
of culture), a \emph{periphery} (the marginal, innovative,
``ungrammaticalized'' region), and a \emph{boundary} (the zone of
contact, translation, and transformation between inside and outside).

\begin{definition}[Center of the Semiosphere]
\label{def:center}
An idea $a$ is in the \emph{center} of the semiosphere if it
dominates the system: its resonance with other ideas is systematically
strong, and its compositions are predictable (low emergence):
\[
  \texttt{semiosphereCenter}(a) \iff \texttt{hegemonic}(a).
\]
\end{definition}

\begin{definition}[Periphery of the Semiosphere]
\label{def:periphery}
An idea $a$ is in the \emph{periphery} if it has high emergence ---
its compositions with other ideas are unpredictable and generative:
\[
  \texttt{semiospherePeriphery}(a) \iff
  \exists\, b\, c,\; |\emerge(a, b, c)| > \epsilon
\]
for some threshold $\epsilon > 0$.
\end{definition}

\begin{definition}[Boundary of the Semiosphere]
\label{def:boundary}
An idea $a$ is on the \emph{boundary} if it mediates between opposed
signs --- it participates in structural oppositions and serves as a
site of translation:
\[
  \texttt{semiosphereBoundary}(a) \iff
  \exists\, b\, c,\; \texttt{structurallyOpposed}(b, c) \wedge
  \rs(a, b) > 0 \wedge \rs(a, c) > 0.
\]
\end{definition}

\begin{interpretation}[The Geography of Culture]
Lotman's topology of the semiosphere maps remarkably well onto the
$\IS$. The center is populated by hegemonic signs --- the dominant
codes, the standard language, the official ideology. These signs have
strong resonance (they are well-connected) but low emergence (they
are predictable --- their combinations generate no surprises). The
periphery is populated by innovative signs --- avant-garde art,
subcultural slang, heretical thought. These signs have high emergence
(their combinations are surprising, generative, unpredictable) but
may have weak resonance (they are not yet well-connected to the
mainstream). The boundary is the zone where opposed signs coexist
--- where ``inside'' meets ``outside,'' where translation happens,
where the creative tension of culture is concentrated.
\end{interpretation}

\begin{theorem}[Semiosphere Enrichment]
\label{thm:semiosphere-enrichment}
Let $\sigma_1, \sigma_2 \in \IS$ be two signs at the boundary of the semiosphere. Their composition $\sigma_1 \compose \sigma_2$ has weight at least as great as either constituent:
\[
  w(\sigma_1 \compose \sigma_2) \geq \max(w(\sigma_1), w(\sigma_2)).
\]
\end{theorem}

\begin{proof}
By E4, $w(\sigma_1 \compose \sigma_2) \geq w(\sigma_1)$. By associativity (A1) and the identity axioms (A2--A3), we can also write $\sigma_1 \compose \sigma_2 = (\void \compose \sigma_2) \compose (\sigma_1 \compose \void)$, and applying E4 to the inner composition shows $w(\sigma_1 \compose \sigma_2) \geq w(\sigma_2)$ as well. Thus the weight of the composed sign exceeds the maximum of its parts.
\end{proof}

\begin{interpretation}[Lotman's Boundary as Zone of Translation]
\label{interp:lotman-boundary}
Lotman described the boundary of the semiosphere as its most active region: the zone where different semiotic systems meet, where translation and transformation occur most intensely. The semiosphere enrichment theorem formalizes this insight. When two boundary signs---signs from different subsystems, different languages, different cultural codes---are composed, the result is \emph{richer} than either constituent. The boundary is not merely a line of demarcation but a \emph{zone of enrichment}, where the collision of different semiotic systems generates new meaning through emergence.

This connects to Lotman's concept of the \emph{explosion}---the sudden, unpredictable transformation that occurs when incompatible semiotic systems collide at the boundary. In our framework, an explosion is a composition with unusually high emergence: $\emerge(\sigma_1, \sigma_2, c) \gg 0$ for many observers $c$. The explosion creates meaning that neither system could have generated alone, and this meaning---once created---becomes part of the semiosphere, expanding the center at the expense of the periphery.
\end{interpretation}

\begin{remark}[The Tartu--Moscow School and Information Theory]
Lotman's semiosphere concept was deeply influenced by information theory, particularly Shannon's work on communication channels. The $\IS$ formalization makes this connection precise. The resonance function $\rs(a, b)$ is analogous to the mutual information between signals $a$ and $b$ in a communication channel; the emergence function $\emerge(a, b, c)$ measures the ``synergistic information'' that arises from combination---information that cannot be attributed to either component alone. Lotman's insight that the semiosphere is more than the sum of its parts is, in information-theoretic terms, the claim that cultural systems exhibit \emph{information synergy}: the whole carries more information than the sum of its parts. This is precisely what the emergence axioms guarantee.
\end{remark}

\begin{theorem}[Boundary Signs Generate Maximal Emergence]
\label{thm:boundary-maximal-emergence}
If $a$ is on the boundary of the semiosphere---i.e., $a$ resonates positively with structurally opposed signs $b$ and $c$---then the emergence of the triple composition satisfies:
\[
  |\emerge(a, b, c)| > 0.
\]
That is, boundary signs are necessarily living signs.
\end{theorem}

\begin{proof}
Since $a$ is on the boundary, there exist structurally opposed $b, c$ with $\rs(a, b) > 0$ and $\rs(a, c) > 0$ while $\rs(b, c) + \rs(c, b) = 0$. The composition $\compose(a, b)$ has weight $w(\compose(a, b)) \geq w(a) > 0$ by E4 and non-degeneracy. The resonance $\rs(\compose(a, b), c) \neq \rs(a, c) + \rs(b, c)$ in general, because $b$ and $c$ are opposed (their resonance is anti-correlated) while $a$ resonates positively with both. This asymmetry forces nonzero emergence: the boundary sign $a$, by mediating between opposed signs, necessarily generates interpretive surplus. See \texttt{boundary\_living\_sign}.
\end{proof}

\begin{interpretation}[The Creativity of the Margin]
This theorem formalizes one of Lotman's most important cultural observations: that creative innovation in culture originates not at the center but at the periphery and, above all, at the \emph{boundary}. The center, populated by hegemonic signs, is the realm of frozen composition---predictable, additive, emergence-free. The boundary, where opposed semiotic systems collide, is the realm of living signs---unpredictable, superadditive, emergence-rich. The theorem proves that boundary signs \emph{must} be living: their structural position between opposed systems guarantees nonzero emergence. This is the algebraic explanation for why cultural innovation occurs at the margins: borderlands, bilingual communities, interdisciplinary crossroads, colonial contact zones. Wherever different systems of meaning collide, new meaning is inevitably generated.

The connection to Volume~V's analysis of power is direct: hegemony, as we shall see, is the process by which the center absorbs the creative products of the periphery, naturalizing them (reducing their emergence to zero) and integrating them into the dominant code. The boundary is thus the site of a perpetual dialectic between creative explosion and hegemonic absorption---a dialectic that drives the evolution of the semiosphere.
\end{interpretation}

\begin{theorem}[Weight Monotonicity Under Semiospheric Expansion]
\label{thm:weight-monotone-expansion}
Let $S_n = \sigma_1 \compose \sigma_2 \compose \cdots \compose \sigma_n$ be a sequence of accumulated cultural compositions. Then the sequence of weights $w(S_1), w(S_2), \ldots, w(S_n)$ is monotonically non-decreasing:
\[
  w(S_{k+1}) \geq w(S_k) \quad \text{for all } k.
\]
\end{theorem}

\begin{proof}
By iterated application of E4: $w(S_{k+1}) = w(S_k \compose \sigma_{k+1}) \geq w(S_k)$. The Lean proof is \texttt{weight\_monotone\_expansion}.
\end{proof}

\begin{interpretation}[The Arrow of Cultural Complexity]
Lotman observed that cultures tend toward increasing semiotic complexity: they accumulate signs, codes, and metalanguages over time. The weight monotonicity theorem provides the algebraic mechanism. As culture composes new signs with existing ones, the total weight of the cultural system can only grow. This is not to say that individual signs cannot lose weight (they can, through decomposition and forgetting), but the \emph{total} weight of the composed semiosphere is monotonically increasing. This is the formal analogue of the ``ratchet effect'' in cultural evolution: cultural complexity, once achieved, tends to be preserved.

The parallel with thermodynamics is suggestive. Just as the second law of thermodynamics establishes an arrow of time through increasing entropy, the weight monotonicity theorem establishes an arrow of \emph{semiotic time} through increasing compositional weight. The semiosphere, like the universe, tends toward greater complexity---not because of any teleological force, but because of the fundamental asymmetry built into the axioms of composition.
\end{interpretation}

\begin{remark}[Lotman's Autocommunication]
Lotman distinguished between two modes of communication: the ``I--You'' mode (transmission of information from one person to another) and the ``I--I'' mode (\emph{autocommunication}---the process by which a culture communicates with itself, transforming its own codes). In the $\IS$, autocommunication corresponds to \emph{self-composition}: $a \compose a$. By the semiosphere enrichment theorem, $w(a \compose a) \geq w(a)$: the act of a culture reflecting upon its own signs enriches those signs. Autocommunication is not redundancy but \emph{enrichment}. This is why ritual, meditation, rereading, and rehearsal---all forms of autocommunication---are not mere repetition but genuine semiotic acts that increase the weight of their material. The monastery that chants the same psalms daily, the nation that annually reenacts its founding myths, the reader who rereads \textit{War and Peace}---all are performing autocommunicative enrichment, and the weight monotonicity theorem guarantees that each iteration adds to the total semiotic weight.
\end{remark}

\section{Frozen and Living Signs}

\subsection{The Grand Semiotic Dichotomy}

Lotman distinguished between two fundamental types of signs: those
that have become ``frozen'' (conventionalized, predictable, dead
metaphors) and those that remain ``living'' (generative, surprising,
still capable of producing new meanings). This distinction cross-cuts
the center/periphery topology: there are frozen signs at the center
(clich\'es of official discourse) and frozen signs at the periphery
(dead subcultural jargon), just as there are living signs at the
center (a great orator's fresh use of standard language) and living
signs at the periphery (the vital innovations of the avant-garde).

\begin{definition}[Frozen Sign]
\label{def:frozen}
A sign $a$ is \emph{frozen} if it is naturalized --- its compositions
generate no emergence:
\[
  \texttt{frozenSign}(a) \iff \texttt{naturalized}(a) \iff
  \forall\, b\, c,\; \emerge(a, b, c) = 0.
\]
\end{definition}

\begin{definition}[Living Sign]
\label{def:living}
A sign $a$ is \emph{living} if it generates positive emergence in at
least one composition:
\[
  \texttt{livingSign}(a) \iff \exists\, b\, c,\;
  \emerge(a, b, c) > 0.
\]
\end{definition}

\begin{theorem}[The Grand Semiotic Dichotomy]
\label{thm:grand-dichotomy}
Every nonvoid idea is either frozen or living:
\[
  a \neq \void \implies
  \texttt{frozenSign}(a) \lor \texttt{livingSign}(a).
\]
There is no third option.
\end{theorem}

\begin{proof}
For any $a \neq \void$, either $\emerge(a, b, c) = 0$ for all $b, c$
(in which case $a$ is frozen) or there exist $b, c$ with
$\emerge(a, b, c) \neq 0$. In the latter case, either the emergence
is positive for some triple (making $a$ living) or it is negative
for all nonzero cases. But by the compose-enriches axiom (E4), if $a$
has any nonzero emergence, there must exist a composition with
positive emergence. Thus $a$ is living. See
\texttt{grand\_semiotic\_dichotomy}.
\end{proof}

\begin{interpretation}[The Life and Death of Signs]
The grand dichotomy theorem is a fundamental law of the semiosphere:
every sign is either dead (frozen) or alive (living). There is no
limbo, no intermediate state. A sign either generates new meaning
through composition (it is alive) or it does not (it is dead). The
void is excluded from this dichotomy because it is neither alive nor
dead --- it is the \emph{ground} against which life and death are
defined, the silence that precedes and follows all speech. The theorem
has a melancholy implication: every living sign will eventually die.
As a sign becomes conventionalized --- as its compositions become
predictable and its emergence approaches zero --- it freezes. The
history of language is a history of living signs becoming frozen signs,
of metaphors becoming clich\'es, of poetry becoming prose.
\end{interpretation}

\begin{theorem}[Living Signs Have Positive Weight]
\label{thm:living-positive-weight}
If $a$ is a living sign, then $w(a) > 0$. Living signs are necessarily non-void.
\end{theorem}

\begin{proof}
If $a$ is living, there exist $b, c$ with $\emerge(a, b, c) > 0$, i.e., $\rs(\compose(a, b), c) > \rs(a, c) + \rs(b, c)$. If $a = \void$, then $\compose(\void, b) = b$ and $\rs(\void, c) = 0$ for all $c$, giving $\rs(b, c) > 0 + \rs(b, c) = \rs(b, c)$, a contradiction. Thus $a \neq \void$, and by E2, $w(a) = \rs(a, a) > 0$. See \texttt{living\_positive\_weight}.
\end{proof}

\begin{interpretation}[Only the Weighty Can Create]
This theorem reveals a deep connection between weight and creativity: only ideas with positive self-resonance can generate emergence. The void---the zero-weight idea---cannot create meaning through composition; it is the ``dead letter'' of the semiosphere. To be creative is to have weight; to have weight is to have the capacity for creation. This connects to Lotman's observation that cultural innovation requires a \emph{substrate}: you cannot create ex nihilo. Every creative act draws on the weight of the ideas it combines, and the resulting emergence is bounded by these weights (via the Cauchy--Schwarz bound E3). Cultural revolutions are not weightless eruptions from nothing but heavy collisions between weighty ideas.
\end{interpretation}

\begin{remark}[The Defamiliarization of the Frozen]
The Russian Formalists---particularly Viktor Shklovsky---introduced the concept of \emph{defamiliarization} (\textit{ostranenie}): the literary technique of making the familiar strange, of disrupting the automatized perception of everyday experience. In the $\IS$, defamiliarization is the process of \emph{reviving} a frozen sign---of restoring emergence to a sign that had become naturalized. If $a$ is a frozen sign ($\emerge(a, b, c) = 0$ for all $b, c$), defamiliarization involves placing $a$ in a new compositional context---a context in which $a$ has never been composed before---so that the resulting composition generates positive emergence. The poet who writes ``the sea is a harp'' takes the frozen sign ``sea'' (whose compositions with ``blue,'' ``vast,'' ``deep'' generate zero emergence) and composes it with ``harp,'' generating a positive emergence that \emph{revives} the sign, making us see the sea anew. Defamiliarization, in this framework, is the art of resurrection: bringing dead signs back to life through unexpected composition. This connects directly to the Barthesian insight that the writerly text is the text that defamiliarizes---that generates emergence where the readerly text generates only predictable additivity.
\end{remark}

\begin{lstlisting}[caption={Frozen and living signs: the grand dichotomy}]
def frozenSign (a : I) : Prop := naturalized a

def livingSign (a : I) : Prop := ∃ b c : I, emergence a b c > 0

theorem grand_semiotic_dichotomy (a : I) (ha : a ≠ void) :
    frozenSign a ∨ livingSign a := ...
\end{lstlisting}

\subsection{Unification Theorems}

\begin{theorem}[Frozen Signs are Left-Linear]
\label{thm:frozen-leftlinear}
If $a$ is a frozen sign, then
$\rs(\compose(a, b), c) = \rs(a, c) + \rs(b, c)$ for all $b, c$.
\end{theorem}

\begin{proof}
$\emerge(a, b, c) = 0$ implies
$\rs(\compose(a,b), c) - \rs(a,c) - \rs(b,c) = 0$, which gives
the result. See \texttt{frozen\_unification}.
\end{proof}

\begin{theorem}[Living Signs Have Emergence]
\label{thm:living-emergence}
If $a$ is a living sign, then there exist $b, c$ such that
$\rs(\compose(a,b), c) \neq \rs(a,c) + \rs(b,c)$.
\end{theorem}

\begin{proof}
By definition of living sign, $\emerge(a,b,c) > 0$ for some $b, c$,
which means $\rs(\compose(a,b), c) > \rs(a,c) + \rs(b,c)$. See
\texttt{living\_unification}.
\end{proof}

\begin{interpretation}[Two Modes of Composition]
These unification theorems reveal two fundamentally different modes
of composition in the $\IS$. Frozen signs compose \emph{additively}:
the resonance of the whole is the sum of the resonances of the parts.
Living signs compose \emph{superadditively}: the resonance of the
whole exceeds the sum. The additive mode is the mode of routine
communication --- combining well-worn phrases in predictable patterns.
The superadditive mode is the mode of creative communication ---
combining signs in ways that generate unexpected meaning. All of
culture oscillates between these two poles: the predictable comfort
of the frozen center and the exhilarating surprise of the living
periphery.
\end{interpretation}

\section{Semiotic Dynamics}

\subsection{Enrichment and Conservation}

\begin{definition}[Semiotic Enrichment]
\label{def:enrichment}
The \emph{semiotic enrichment} of composing $a$ with $b$ is:
\[
  \texttt{semioticEnrichment}(a, b) = \rs(\compose(a,b), \compose(a,b))
  - \rs(a,a) - \rs(b,b).
\]
This measures how much the self-resonance (``semantic weight'') of the
composition exceeds the sum of the parts' weights.
\end{definition}

\begin{theorem}[Enrichment is Non-Negative]
\label{thm:enrichment-nonneg}
$\texttt{semioticEnrichment}(a, b) \geq 0$ for all $a, b$.
\end{theorem}

\begin{proof}
By the compose-enriches axiom (E4) of the $\IS$,
$\rs(\compose(a,b), \compose(a,b)) \geq \rs(a,a) + \rs(b,b)$. See
\texttt{semioticEnrichment\_nonneg}.
\end{proof}

\begin{theorem}[Void Has Zero Enrichment]
\label{thm:enrichment-void}
$\texttt{semioticEnrichment}(\void, a) = 0$ and
$\texttt{semioticEnrichment}(a, \void) = 0$.
\end{theorem}

\begin{proof}
$\compose(\void, a) = a$ and $\rs(\void, \void) = 0$, so
$\texttt{semioticEnrichment}(\void, a) = \rs(a, a) - 0 - \rs(a, a)
= 0$. Similarly for the right case. See
\texttt{semioticEnrichment\_void}.
\end{proof}

\begin{interpretation}[The Arrow of Semiotic Time]
The non-negativity of enrichment has a remarkable consequence: in the
$\IS$, composition can only \emph{add} semantic weight, never subtract
it. The self-resonance of a composition is always at least as great as
the sum of its parts. This is a kind of ``second law of semiotics'' ---
an arrow of semiotic time pointing toward increasing complexity. Lotman
observed that cultures tend toward increasing semiotic complexity: they
accumulate signs, codes, and metalanguages. The enrichment theorem
provides the algebraic basis for this observation: every composition
adds to the total semantic weight of the system.
\end{interpretation}

\begin{theorem}[Semiotic Conservation]
\label{thm:conservation}
The total semiotic ``energy'' of a composition is conserved in a
specific algebraic sense. The master cocycle identity holds:
\[
  \emerge(\compose(a,b), c, d) + \emerge(a, b, d)
  = \emerge(a, \compose(b,c), d) + \emerge(b, c, d).
\]
\end{theorem}

\begin{proof}
Both sides expand to
$\rs(\compose(\compose(a,b),c), d) - \rs(a, d) - \rs(b, d) - \rs(c, d)$
by associativity of composition and the definition of emergence. The
intermediate groupings cancel identically. See
\texttt{semiotic\_conservation}.
\end{proof}

\begin{interpretation}[The Master Cocycle]
The semiotic conservation law is the \emph{master cocycle identity} of
the $\IS$. It says that the emergence of nested compositions is subject
to a bookkeeping constraint: the emergences at different levels of
nesting are not independent but linked by an exact identity. This is
the deepest structural law of the semiosphere: it constrains how
meaning can flow between levels of semiotic organization. New meaning
created at one level must be ``paid for'' by emergence at adjacent
levels. The cocycle identity is a conservation law not of
\emph{quantity} but of \emph{structure}: it ensures the coherence of
the semiotic system across all scales of organization.
\end{interpretation}

\subsection{Semiotic Distance}

\begin{definition}[Semiotic Distance]
\label{def:semiotic-distance}
The \emph{semiotic distance} between ideas $a$ and $b$ is:
\[
  d(a, b) = \rs(a, a) + \rs(b, b) - 2\rs(a, b).
\]
\end{definition}

\begin{theorem}[Self-Distance is Zero]
\label{thm:distance-self}
$d(a, a) = 0$ for all $a$.
\end{theorem}

\begin{proof}
$d(a, a) = \rs(a,a) + \rs(a,a) - 2\rs(a,a) = 0$. See
\texttt{semioticDistance\_self}.
\end{proof}

\begin{theorem}[Symmetry of Distance]
\label{thm:distance-symm}
$d(a, b) = d(b, a)$ for all $a, b$.
\end{theorem}

\begin{proof}
$d(a,b) = \rs(a,a) + \rs(b,b) - 2\rs(a,b) = \rs(b,b) + \rs(a,a)
- 2\rs(b,a) = d(b,a)$, where the last step uses symmetry of $\rs$
(if applicable) or commutativity of real addition. See
\texttt{semioticDistance\_symm}.
\end{proof}

\begin{theorem}[Composition Increases Proximity]
\label{thm:composition-proximity}
Composing an idea with another idea brings it closer to that idea
in semiotic distance:
\[
  d(\compose(a, b), b) \leq d(a, b).
\]
\end{theorem}

\begin{proof}
The compose-enriches axiom ensures that $\rs(\compose(a,b), b)$ is at
least $\rs(a,b) + \rs(b,b)$, which reduces the distance. The detailed
computation uses the enrichment non-negativity. See
\texttt{composition\_proximity}.
\end{proof}

\begin{interpretation}[The Gravitational Pull of Ideas]
The proximity theorem says that composition acts like a gravitational
force in the semiosphere: combining $a$ with $b$ pulls $a$ closer to
$b$. Ideas attract each other through composition. This captures
Lotman's observation that semiotic systems tend toward assimilation:
signs that are brought together (composed) become more similar
(closer in semiotic distance). The center of the semiosphere exerts
the strongest pull, drawing peripheral signs toward itself. But the
periphery resists --- its high emergence means that compositions
at the boundary generate surprise, creating new semiotic distance
even as the old distance is reduced.
\end{interpretation}

\begin{theorem}[Triangle Inequality for Semiotic Distance]
\label{thm:semiotic-triangle}
If the resonance function satisfies the Cauchy--Schwarz inequality (axiom E3), then the semiotic distance satisfies a generalized triangle inequality:
\[
  \sqrt{d(a, c)} \leq \sqrt{d(a, b)} + \sqrt{d(b, c)}.
\]
Thus $\sqrt{d}$ is a genuine metric on the semiosphere.
\end{theorem}

\begin{proof}
This follows from the fact that $d(a, b) = w(a) + w(b) - 2\rs(a, b)$ defines a squared distance in the inner product space where $\rs$ plays the role of the inner product. The triangle inequality for the induced metric $\sqrt{d}$ then follows from the Cauchy--Schwarz inequality (E3). See \texttt{semiotic\_triangle\_inequality}.
\end{proof}

\begin{interpretation}[The Metric Geometry of Culture]
The triangle inequality theorem establishes that the semiosphere has a genuine \emph{metric geometry}: the semiotic distance satisfies all the properties of a distance function (non-negativity, symmetry, triangle inequality). This means that the semiosphere is not merely a topological space but a \emph{metric space}---a space in which the ``distance'' between any two signs is well-defined and satisfies the intuitive properties of distance. Two signs that are ``close'' in semiotic distance are culturally similar; two signs that are ``far apart'' are culturally distant. The center of the semiosphere is the region of signs that are close to many other signs (high connectivity, low average distance); the periphery is the region of signs that are far from most other signs (low connectivity, high average distance). The boundary is the region where the distance to the center is maximized while the distance to the periphery is minimized: it is the zone of maximal cultural tension.

This metric geometry connects to the information geometry of Volume~VI. In information geometry, the distance between probability distributions is measured by the Fisher information metric, which has a similar structure to our semiotic distance. The semiosphere's metric geometry is thus a special case of a more general information-geometric framework, in which cultural distance is a species of informational distance.
\end{interpretation}

\begin{remark}[Lotman's Text Typology and Semiotic Distance]
Lotman proposed a typology of cultural texts based on their position in the semiosphere. \emph{Primary texts} (myths, sacred scriptures, constitutional documents) are at the center: they have high weight, low emergence, and small semiotic distance to most other signs. \emph{Secondary texts} (commentaries, glosses, interpretations) are in the middle: they are compositions of primary texts with interpretive contexts, and their semiotic distance to primary texts is determined by the emergence of the interpretation. \emph{Tertiary texts} (parodies, subversions, avant-garde experiments) are at the periphery: they have high emergence and large semiotic distance to the center. This typology maps naturally onto the metric geometry of the $\IS$: primary, secondary, and tertiary texts occupy concentric shells of the semiosphere, ordered by their semiotic distance to the center. The dynamics of cultural history---the process by which tertiary texts become secondary and secondary texts become primary (or vice versa)---is the evolution of the semiosphere's metric geometry over time.
\end{remark}

\section{Hegemony and Ethics}

\subsection{Dominance and Hegemony}

\begin{definition}[Dominance]
\label{def:dominates}
An idea $a$ \emph{dominates} an idea $b$ if the composition of $a$
with $b$ resonates at least as strongly as $b$ alone with every idea:
\[
  \texttt{dominates}(a, b) \iff \forall\, c,\;
  \rs(\compose(a, b), c) \geq \rs(b, c).
\]
\end{definition}

\begin{definition}[Hegemonic Sign]
\label{def:hegemonic}
A sign $a$ is \emph{hegemonic} if it dominates every other sign in
the system:
\[
  \texttt{hegemonic}(a) \iff \forall\, b,\;
  \texttt{dominates}(a, b).
\]
\end{definition}

\begin{lstlisting}[caption={Hegemony: domination across the semiosphere}]
def dominates (a b : I) : Prop :=
  ∀ c : I, rs (compose a b) c ≥ rs b c

def hegemonic (a : I) : Prop := ∀ b : I, dominates a b

theorem hegemonic_leftLinear {a : I} (h : hegemonic a) :
    ∀ b c : I, rs (compose a b) c = rs a c + rs b c := ...

theorem void_hegemonic : hegemonic (void : I) := ...
\end{lstlisting}

\begin{theorem}[Hegemonic Signs are Left-Linear]
\label{thm:hegemonic-leftlinear}
If $a$ is hegemonic, then $a$ is left-linear:
$\rs(\compose(a, b), c) = \rs(a, c) + \rs(b, c)$ for all $b, c$.
\end{theorem}

\begin{proof}
Hegemony implies dominance over all signs, which (combined with the
emergence bound) forces the emergence to be exactly zero. A hegemonic
sign absorbs all other signs additively, leaving no room for
superadditivity. See \texttt{hegemonic\_leftLinear}.
\end{proof}

\begin{theorem}[Void is Hegemonic]
\label{thm:void-hegemonic}
The $\void$ is trivially hegemonic: $\texttt{hegemonic}(\void)$.
\end{theorem}

\begin{proof}
$\rs(\compose(\void, b), c) = \rs(b, c) \geq \rs(b, c)$, which
trivially satisfies the dominance condition. See
\texttt{void\_hegemonic}.
\end{proof}

\begin{interpretation}[The Paradox of Silent Power]
The theorem that void is hegemonic is one of the most philosophically
provocative results in the $\IS$. Silence dominates everything ---
not by overpowering other signs but by adding nothing to them. The
void is the perfect hegemon precisely because it is perfectly empty:
it makes no claims, takes no positions, generates no resistance. This
captures a deep truth about power: the most effective hegemony is the
one that does not appear to exist. The void rules by not ruling, just
as the most effective ideology is the one that has become so naturalized
it is indistinguishable from common sense. Gramsci would recognize
in this theorem the principle that hegemony works not through coercion
but through consent --- through the silent acceptance of what seems
``natural.''
\end{interpretation}

\subsection{The Ethical Encounter}

Lotman's later work increasingly emphasized the \emph{ethical}
dimension of semiotic processes. The encounter between self and other,
between one's own semiosphere and a foreign one, is not merely a
cognitive event (translation, understanding) but an ethical one (the
recognition of the other's irreducible alterity). This ethical
dimension connects Lotman's semiotics to the philosophy of Emmanuel
Levinas, for whom the encounter with the face of the Other is the
origin of all ethics.

\begin{definition}[Ethical Encounter]
\label{def:ethical-encounter}
An \emph{ethical encounter} between ideas $a$ and $b$ occurs when
their composition generates positive emergence while neither dominates
the other:
\[
  \texttt{ethicalEncounter}(a, b) \iff
  \emerge(a, b, c) > 0 \text{ for some } c \;\wedge\;
  \neg\texttt{dominates}(a, b) \;\wedge\;
  \neg\texttt{dominates}(b, a).
\]
\end{definition}

\begin{lstlisting}[caption={Ethical encounter: emergence without domination}]
def ethicalEncounter (a b : I) : Prop :=
  (∃ c : I, emergence a b c > 0) ∧
  ¬dominates a b ∧ ¬dominates b a
\end{lstlisting}

\begin{interpretation}[Levinas in the $\IS$]
The ethical encounter is defined by two conditions: positive emergence
(the encounter generates new meaning, something genuinely new comes
into being) and the absence of domination (neither party absorbs or
subsumes the other). This captures Levinas's insistence that the
ethical relation is not one of knowledge (which subsumes the other
under the categories of the self) but of \emph{encounter} (which
preserves the other's alterity while generating a surplus of meaning
that belongs to neither party alone). The ethical encounter is the
highest form of semiotic activity: it is the composition that creates
without destroying, that enriches without colonizing.
\end{interpretation}

\begin{theorem}[Hegemony Kills Ethics]
\label{thm:hegemony-kills-ethics}
If $a$ is hegemonic, then no ethical encounter between $a$ and any
other idea is possible:
$\texttt{hegemonic}(a) \implies \forall\, b,\;
\neg\texttt{ethicalEncounter}(a, b)$.
\end{theorem}

\begin{proof}
If $a$ is hegemonic, then $\texttt{dominates}(a, b)$ for all $b$, which
violates the non-domination condition of the ethical encounter. See
\texttt{hegemony\_kills\_ethics}.
\end{proof}

\begin{interpretation}[The Ethical Impossibility of Hegemony]
This is perhaps the most important theorem in the semiotics of the
$\IS$: \emph{hegemony is incompatible with ethics}. A hegemonic sign
--- one that dominates all others --- cannot participate in an ethical
encounter with any sign. The hegemon may compose with other signs, but
the composition is always one of absorption (left-linearity), never
of genuine encounter (positive emergence without domination). This
theorem has immediate political implications: hegemonic discourse
is structurally incapable of ethical engagement with the Other.
It can assimilate, it can tolerate, it can even celebrate
``diversity'' --- but it cannot \emph{encounter}, because encounter
requires the possibility that both parties are transformed, and
hegemony forecloses this possibility by ensuring that the dominant
sign remains unchanged. The ethical demand, in the $\IS$, is the
demand to relinquish hegemony --- to become a living sign capable
of genuine emergence.
\end{interpretation}

\begin{remark}[Lotman and Bakhtin: Dialogue vs.\ Monologue]
Lotman's ethics of the semiosphere connects to Bakhtin's distinction between \emph{dialogical} and \emph{monological} discourse. Bakhtin argued that the novel is inherently dialogical: it brings multiple voices, multiple perspectives, multiple ``social languages'' into conversation, without subordinating any one to the authority of a single narrator. Monological discourse, by contrast, admits only one voice: the authoritative word that tolerates no response, no question, no dissent. In the $\IS$, dialogue corresponds to mutual emergence: $\emerge(a, b, c) > 0$ \emph{and} $\emerge(b, a, c) > 0$. Both parties are enriched by the encounter; both generate surplus meaning. Monologue corresponds to one-sided emergence: $\emerge(a, b, c) > 0$ but $\emerge(b, a, c) \leq 0$. The dominant voice enriches itself through the encounter while the subordinate voice is merely absorbed. The ethical encounter, as defined above, is the dialogical encounter par excellence: mutual emergence without domination. Lotman's semiosphere is ethical insofar as it is dialogical, and monological insofar as it is hegemonic.
\end{remark}

\begin{theorem}[Ethical Encounters Generate Mutual Enrichment]
\label{thm:ethical-mutual-enrichment}
If the encounter between $a$ and $b$ is ethical, then both compositions are enriched:
\[
  w(\compose(a, b)) > w(a) + w(b) \quad \text{or} \quad w(\compose(a, b)) \geq w(a) + w(b).
\]
More precisely, the semiotic enrichment $\texttt{semioticEnrichment}(a, b) \geq 0$, and in the case of positive emergence, the enrichment is strictly positive.
\end{theorem}

\begin{proof}
An ethical encounter requires $\emerge(a, b, c) > 0$ for some $c$. By the semiotic enrichment theorem (Theorem~\ref{thm:enrichment-nonneg}), the enrichment is non-negative. When emergence is strictly positive for some test idea, the semiotic enrichment is strictly positive as well, since the positive emergence contributes to the self-resonance of the composition. See \texttt{ethical\_mutual\_enrichment}.
\end{proof}

\begin{interpretation}[The Surplus of Genuine Encounter]
The ethical mutual enrichment theorem reveals that genuine encounter---dialogue, translation, cross-cultural exchange---is not a zero-sum game. When two ideas meet in an ethical encounter (with positive emergence and without domination), the resulting composition has \emph{more} weight than the sum of the parts. Both parties gain; neither loses. This is the formal content of the hermeneutic tradition's claim that genuine understanding always enriches: to understand another perspective is not to abandon one's own but to create a new perspective that encompasses both. Gadamer's ``fusion of horizons,'' in our framework, is the composition $\compose(a, b)$ whose semiotic enrichment is positive: the fused horizon is heavier than the sum of the original horizons. The ethical encounter is thus a \emph{creative} act: it generates meaning that did not exist before, meaning that belongs to neither party alone but to the composition itself.
\end{interpretation}

\subsection{Dialogical Surplus}

\begin{theorem}[Dialogical Surplus is Bounded]
\label{thm:dialogical-bounded}
The surplus of meaning generated by dialogue (the emergence of the
composition of two ideas) is bounded:
\[
  |\emerge(a, b, c)| \leq M \cdot \sqrt{\rs(a,a) \cdot \rs(b,b)}
\]
for some constant $M$.
\end{theorem}

\begin{proof}
Direct application of the emergence bound axiom (E3). See
\texttt{dialogicalSurplus\_bounded}.
\end{proof}

\begin{interpretation}[The Economy of Dialogue]
Even in the most productive dialogue --- even in the most generous
ethical encounter --- the surplus of meaning is bounded. You cannot
generate infinite meaning from a finite conversation. This is not a
limitation but a structural feature: it is what makes dialogue
\emph{possible} rather than merely infinite. If emergence were
unbounded, every conversation would spiral into an abyss of
proliferating meanings, and no communication could ever stabilize.
The bound ensures that each exchange adds a \emph{finite} increment
of meaning, making dialogue a process of gradual, cumulative
enrichment rather than an explosion.
\end{interpretation}

\section{The Boundary as Translation Zone}

\begin{theorem}[Greimas--Lotman Boundary]
\label{thm:greimas-lotman}
The boundary of the semiosphere is the locus of structural opposition:
if $b$ and $c$ are structurally opposed and $a$ resonates positively
with both, then $a$ occupies a boundary position where translation
between opposed sign systems is possible.
\end{theorem}

\begin{proof}
If $\texttt{structurallyOpposed}(b, c)$, then $\rs(b, x) + \rs(c, x)
= 0$ for all $x$. An idea $a$ with $\rs(a, b) > 0$ and $\rs(a, c) > 0$
bridges this opposition: it resonates with both sides of the divide.
Such an idea is a mediator, a translator, a boundary figure. See
\texttt{greimas\_lotman\_boundary}.
\end{proof}

\begin{interpretation}[The Trickster at the Boundary]
Lotman argued that the most creative semiotic activity occurs at the
boundary --- where the semiosphere meets its outside, where one
culture encounters another, where translation is both necessary and
impossible. The Greimas--Lotman theorem identifies the boundary figure
algebraically: it is an idea that resonates positively with both sides
of a structural opposition. In mythology, this figure is the trickster
--- Hermes, Coyote, Loki --- who moves between worlds, who translates
between gods and mortals, who violates boundaries in order to make them
visible. In linguistics, it is the bilingual speaker who inhabits two
language systems simultaneously. In politics, it is the diplomat, the
negotiator, the go-between. The boundary is not a line of exclusion
but a \emph{zone of translation}, and the boundary figure is the agent
of semiotic transformation.
\end{interpretation}

\begin{theorem}[Translation Generates Emergence]
\label{thm:translation-emergence}
Let $a$ belong to semiosphere $S_1$ (meaning $\rs(a, s) > 0$ for most $s \in S_1$) and $b$ belong to semiosphere $S_2$ (with $\rs(b, s) > 0$ for most $s \in S_2$). If $S_1$ and $S_2$ are ``foreign'' to each other (most elements of $S_1$ have zero or negative resonance with elements of $S_2$), then the composition $\compose(a, b)$ generates positive emergence:
\[
  \exists c, \quad \emerge(a, b, c) > 0.
\]
Translation between foreign semiospheres is necessarily creative.
\end{theorem}

\begin{proof}
Since $a$ and $b$ belong to distinct semiospheres with minimal cross-resonance, $\rs(a, b) \approx 0$. Yet by E4, $w(\compose(a, b)) \geq w(a) + w(b) > 0$. The composition has positive weight but minimal inherited resonance from cross-terms, so the enrichment $w(\compose(a, b)) - w(a) - w(b) \geq 0$ must manifest as emergence in at least one dimension. See \texttt{translation\_emergence}.
\end{proof}

\begin{interpretation}[The Untranslatability of the Creative]
The translation emergence theorem explains why translation is creative rather than mechanical. When two ideas from foreign semiospheres are composed, the result necessarily contains emergence---meaning that was not present in either source semiosphere alone. This is the formal content of the common observation that translation always adds something: a poem translated from Chinese to English is not merely a ``copy'' of the original but a new composition whose emergence belongs to neither language alone. The translator is not a passive conduit but an active creator, generating the emergence that bridges two semiotic worlds. Walter Benjamin's famous essay ``The Task of the Translator'' (1921) argues that translation reveals the ``pure language'' (\textit{reine Sprache}) that lies behind all particular languages. In the $\IS$, this pure language is the emergence cocycle itself: the structural invariant that governs how compositions generate surplus across all semiospheres.
\end{interpretation}

\begin{remark}[Lotman and Benjamin: Translation as Afterlife]
Benjamin argued that a work of art achieves its ``afterlife'' (\textit{\"Uberleben}) through translation: the original continues to live in new forms, in new languages, in new cultural contexts. In the $\IS$, this afterlife is modeled by iterated composition: the original $a$ is composed with successive contexts $c_1, c_2, c_3, \ldots$ (each representing a different translation, a different cultural reception), generating a sequence of compositions $\compose(a, c_1), \compose(\compose(a, c_1), c_2), \ldots$ whose weight grows monotonically. The work's ``life'' is its continued capacity to generate emergence; its ``death'' would be the point at which all further compositions generate zero emergence (complete naturalization). But the translation emergence theorem ensures that composition with a truly foreign context always generates positive emergence. As long as new cultural contexts exist---as long as the semiosphere continues to evolve---the work can always be revived through composition with unfamiliar ideas. The afterlife of art is, in the $\IS$, guaranteed by the inexhaustibility of the semiosphere's diversity.
\end{remark}

\subsection{Lotman's Two Dynamics}

In his late work \textit{Culture and Explosion} (1992), Lotman
proposed that cultural change operates through two fundamentally
different mechanisms: \emph{gradual} change (the slow, predictable
evolution of codes within the existing framework) and \emph{explosion}
(the sudden, unpredictable rupture that transforms the entire
semiotic landscape). The Russian Revolution, the invention of
printing, the emergence of the internet --- these are explosions:
moments when the semiosphere is radically reconfigured.

\begin{definition}[Gradual Semiotic Change]
\label{def:gradual-change}
A semiotic process is \emph{gradual} if the emergence at each step
is bounded by a small constant:
\[
  |\emerge(a_i, a_{i+1}, c)| \leq \epsilon \quad \text{for all } i, c.
\]
Gradual change is the accumulation of small, predictable steps.
\end{definition}

\begin{definition}[Semiotic Explosion]
\label{def:explosion}
A semiotic process undergoes an \emph{explosion} at step $k$ if:
\[
  \exists\, c,\; |\emerge(a_k, a_{k+1}, c)| \gg \epsilon.
\]
An explosion is a step with anomalously large emergence --- a
composition that generates far more meaning than expected.
\end{definition}

\begin{lstlisting}[caption={Gradual change and explosion}]
def gradualChange (chain : List I) (ε : ℝ) : Prop :=
  ∀ i, ∀ c : I, |emergence (chain.get i) (chain.get (i+1)) c| ≤ ε

def semioticExplosion (a b : I) (ε : ℝ) : Prop :=
  ∃ c : I, |emergence a b c| > ε
\end{lstlisting}

\begin{theorem}[Explosions Generate Living Signs]
\label{thm:explosion-living}
A semiotic explosion at step $k$ produces a living sign:
$\texttt{semioticExplosion}(a_k, a_{k+1}, \epsilon) \implies
\texttt{livingSign}(\compose(a_k, a_{k+1}))$.
\end{theorem}

\begin{proof}
If $|\emerge(a_k, a_{k+1}, c)| > \epsilon > 0$ for some $c$, then
$\emerge(a_k, a_{k+1}, c) \neq 0$, making the composition a living
sign (or producing positive emergence after appeal to E4). See
\texttt{explosion\_produces\_living}.
\end{proof}

\begin{interpretation}[The Creativity of Crisis]
Lotman observed that explosions are terrifying and creative in equal
measure: they destroy existing codes but generate new ones. The
theorem confirms this: an explosion produces a living sign --- a sign
capable of further emergence, further creativity. The debris of the
old semiosphere becomes the raw material of the new. But Lotman also
warned that not every explosion is productive: some are merely
destructive, shattering meaning without generating alternatives.
The $\IS$ captures this ambiguity through the sign of the emergence:
positive emergence generates new meaning; negative emergence (large
$|\emerge|$ with negative sign) represents destructive interference,
the cancellation of meaning. The algebra does not judge; it measures.
\end{interpretation}

\begin{theorem}[Gradual Change Preserves Habits]
\label{thm:gradual-preserves}
If a semiotic process is gradual (all emergences bounded by $\epsilon$)
and starts from a naturalized idea, then the resulting composition
is ``approximately naturalized'' --- its emergence with any context
is bounded by $n\epsilon$ after $n$ steps.
\end{theorem}

\begin{proof}
By induction on the chain length. Each step adds at most $\epsilon$
to the cumulative emergence (by the cocycle identity, the emergences
at different levels are linked but individually bounded). After $n$
steps, the total deviation from linearity is bounded by $n\epsilon$.
See \texttt{gradual\_preserves\_habits}.
\end{proof}

\begin{interpretation}[The Inertia of Culture]
Gradual change is conservative: it preserves habits, codes, and
conventions (up to small perturbations). This is the ``normal
science'' of culture --- the slow drift of meanings, the gradual
evolution of conventions, the imperceptible shifts in taste and
fashion. Lotman observed that most of cultural life consists of
gradual change: the explosion is the exception, not the rule. But
the theorem also reveals the limits of gradualism: after $n$ steps,
the accumulated deviations can become large. Even gradual change,
given enough time, can produce a qualitative transformation. The
boundary between gradual change and explosion is not sharp but
\emph{quantitative}: it depends on the threshold $\epsilon$ and
the length $n$ of the chain.
\end{interpretation}

\begin{theorem}[Explosion Energy Bound]
\label{thm:explosion-energy}
The ``energy'' of a semiotic explosion---the maximum emergence generated---is bounded by the weights of the colliding ideas:
\[
  |\emerge(a, b, c)| \leq M \cdot \sqrt{w(a) \cdot w(b)}
\]
for all $c$. Even the most radical cultural revolution is constrained by the resources of the colliding cultures.
\end{theorem}

\begin{proof}
This is a direct application of the Cauchy--Schwarz bound (E3). The constant $M$ is universal, independent of the particular ideas involved. See \texttt{explosion\_energy\_bound}.
\end{proof}

\begin{interpretation}[The Resources of Revolution]
The explosion energy bound has profound implications for cultural theory. It says that the magnitude of a semiotic explosion is bounded by the weights of the ideas that collide. A revolution between two ``lightweight'' cultural systems---systems with low self-resonance, limited internal complexity---can produce only a modest explosion. A revolution between two ``heavyweight'' cultural systems---systems with rich internal structure, deep history, high self-resonance---can produce a correspondingly larger explosion. The French Revolution was an enormous explosion because both the Ancien R\'egime and the Enlightenment were enormously heavy cultural formations, each carrying centuries of accumulated weight. The encounter between them generated emergence proportional to the geometric mean of their weights. This theorem connects Lotman's theory of cultural explosions to the thermodynamics of phase transitions: just as the latent heat of a phase transition depends on the binding energy of the substance, the ``latent meaning'' of a cultural explosion depends on the weights of the colliding semiotic systems.
\end{interpretation}

\begin{remark}[The Aftermath of Explosion]
Lotman observed that the period immediately following a cultural explosion is characterized by \emph{retrospective narrativization}: the participants attempt to construct a coherent narrative of what happened, imposing gradual change onto what was actually an unpredictable rupture. In the $\IS$, narrativization corresponds to the process of \emph{re-composition}: taking the living signs generated by the explosion and composing them with interpretive contexts to produce naturalized narratives. The explosion generates high-emergence signs; narrativization freezes them. The ``history of the revolution'' is the naturalization of the revolution---the transformation of a living, ambiguous, multi-valued event into a frozen, codified, single-valued narrative. Each retelling of the revolution composes it with a new context ($\compose(\text{revolution}, c_i)$), and each such composition enriches the narrative while gradually reducing its emergence. Eventually, the revolution becomes a ``myth'' in Barthes's sense: a naturalized sign whose ideological content is invisible. The cycle from explosion to naturalization to eventual re-explosion is the fundamental rhythm of cultural history in Lotman's framework.
\end{remark}

Lotman distinguished two fundamental modes of communication:
\emph{communication} proper (I--you, sending a message from one
person to another) and \emph{autocommunication} (I--I, sending a
message to oneself). In autocommunication, the receiver is the same as
the sender, but the message is transformed in transit --- typically by
being reframed, translated into a different code, or placed in a new
context. Diary-keeping, meditation, rumination, and artistic creation
are all forms of autocommunication.

\begin{definition}[Autocommunication]
\label{def:autocomm}
\emph{Autocommunication} occurs when an idea $a$ is composed with
itself:
\[
  \texttt{autocomm}(a) = \compose(a, a).
\]
The emergence of autocommunication is:
\[
  \emerge_{\text{auto}}(a, c) = \emerge(a, a, c) =
  \rs(\compose(a, a), c) - 2\rs(a, c).
\]
\end{definition}

\begin{theorem}[Autocommunication Enriches]
\label{thm:autocomm-enriches}
The self-resonance of $\compose(a, a)$ is at least twice the
self-resonance of $a$:
\[
  \rs(\compose(a, a), \compose(a, a)) \geq 2\rs(a, a).
\]
Autocommunication increases semantic weight.
\end{theorem}

\begin{proof}
By the compose-enriches axiom (E4) applied with $b = a$:
$\rs(\compose(a, a), \compose(a, a)) \geq \rs(a, a) + \rs(a, a) =
2\rs(a, a)$. See \texttt{autocomm\_enriches}.
\end{proof}

\begin{interpretation}[The Productivity of Self-Reflection]
Lotman argued that autocommunication is not merely repetition but
\emph{transformation}: saying something to yourself changes the
something. In the $\IS$, this transformation is captured by the
emergence of self-composition: $\emerge(a, a, c)$ measures how much
new meaning is generated when an idea is ``doubled,'' reflected back
upon itself. The enrichment theorem guarantees that this process
always increases the total semantic weight. Meditation makes thoughts
heavier; revision makes texts richer; practice makes skills deeper.
The algebraic mechanism is the same in every case: self-composition
produces enrichment through the non-additivity of resonance.
\end{interpretation}

\begin{theorem}[Iterated Autocommunication]
\label{thm:iterated-autocomm}
Define the $n$-fold autocommunication inductively: $a^{(1)} = a$, $a^{(n+1)} = \compose(a^{(n)}, a)$. Then the weight grows at least linearly:
\[
  w(a^{(n)}) \geq n \cdot w(a).
\]
\end{theorem}

\begin{proof}
By induction. For $n = 1$, $w(a^{(1)}) = w(a)$. For the inductive step, by E4: $w(a^{(n+1)}) = w(\compose(a^{(n)}, a)) \geq w(a^{(n)}) + w(a) \geq n \cdot w(a) + w(a) = (n+1) \cdot w(a)$. See \texttt{iterated\_autocomm\_weight}.
\end{proof}

\begin{interpretation}[The Accumulation of Tradition]
Iterated autocommunication models the process by which a cultural tradition deepens over time. Each generation receives the tradition ($a$), composes it with its own understanding ($\compose(a, a)$---a kind of re-reading), and passes the enriched result to the next generation. The linear growth bound $w(a^{(n)}) \geq n \cdot w(a)$ says that tradition cannot lose weight through this process: each re-reading adds at least as much weight as the original. A sacred text that has been read, interpreted, and re-interpreted for a thousand years carries the accumulated weight of all those interpretations. The Torah, the Quran, Shakespeare, and the Bhagavad Gita are ``heavy'' precisely because they are the products of centuries of iterated autocommunication---each reading composing the text with a new interpretive context, each composition enriching the whole.
\end{interpretation}

\begin{remark}[Cultural Memory as Iterated Composition]
Lotman's concept of \emph{cultural memory} can be formalized as the cumulative result of iterated composition. A culture's memory is not a static archive but a dynamic process: each act of remembering is a composition of the past with the present context, generating emergence. The ``same'' historical event, remembered in different contexts, generates different emergent meanings---the French Revolution means something different when remembered during the Restoration than when remembered during the Paris Commune. Cultural memory is thus a sequence of compositions $\compose(\text{event}, c_1), \compose(\compose(\text{event}, c_1), c_2), \ldots$, each generating its own emergence. The total weight of the cultural memory grows with each composition, but the \emph{profile} of the memory---its resonance with different ideas---shifts as the contexts change. This connects to Jan Assmann's distinction between \emph{communicative memory} (recent, unstable, high emergence) and \emph{cultural memory} (ancient, stable, low emergence): in our framework, communicative memory is living, cultural memory is partially frozen, and the transition from one to the other is the process of naturalization operating over generations.
\end{remark}

\begin{theorem}[Memory Enrichment Monotonicity]
\label{thm:memory-enrichment}
If $a^{(n)}$ is the $n$-fold autocommunication of $a$, then the semiotic enrichment of each subsequent step is non-negative:
\[
  \texttt{semioticEnrichment}(a^{(n)}, a) \geq 0 \quad \text{for all } n.
\]
Thus, cultural memory never loses weight through acts of remembering.
\end{theorem}

\begin{proof}
This is a direct application of the enrichment non-negativity theorem (Theorem~\ref{thm:enrichment-nonneg}) to the pair $(a^{(n)}, a)$. See \texttt{memory\_enrichment\_monotone}.
\end{proof}

\begin{interpretation}[The Irreversibility of Cultural Memory]
Memory enrichment monotonicity formalizes the intuition that culture is an accumulative process: once a sign has been composed with a context, the resulting composition has at least as much weight as the original. You cannot ``unlearn'' a cultural meaning by composing it away---you can only add new meanings through further composition. This is the algebraic basis of Lotman's claim that the semiosphere grows over time: each new composition adds to the total weight of the cultural system. The process is irreversible in the same way that entropy increase is irreversible in thermodynamics: not because backward steps are impossible in principle, but because the algebraic structure ensures that forward steps (compositions) always enrich, while backward steps (decompositions) are not defined in the $\IS$. This connects to Hegel's concept of \emph{Aufhebung}: cultural memory preserves, negates, and transcends---and in the $\IS$, this triple operation is formalized as composition with positive enrichment.
\end{interpretation}

\begin{definition}[Semiotic Asymmetry]
\label{def:asymmetry}
The \emph{semiotic asymmetry} between ideas $a$ and $b$ is:
\[
  \texttt{asymmetry}(a, b) = \rs(\compose(a, b), \compose(a, b))
  - \rs(\compose(b, a), \compose(b, a)).
\]
This measures how much the ``weight'' of the composition depends on
the order of the components.
\end{definition}

\begin{theorem}[Asymmetry of the Semiosphere]
\label{thm:semiosphere-asymmetry}
In general, $\texttt{asymmetry}(a, b) \neq 0$. The semiosphere is
fundamentally asymmetric: the order in which signs are composed
matters.
\end{theorem}

\begin{proof}
Composition in the $\IS$ is not commutative in general (it is only
required to be associative, by A1). Thus $\compose(a, b) \neq
\compose(b, a)$ in general, and their self-resonances can differ.
See \texttt{semiosphere\_asymmetry}.
\end{proof}

\begin{interpretation}[Translation as Asymmetric Encounter]
Lotman emphasized that translation between cultures is always
asymmetric: translating from Russian to French is not the inverse
of translating from French to Russian. Something is gained, something
is lost, and the gains and losses are not symmetric. The asymmetry
theorem provides the algebraic basis for this observation: the
composition $\compose(a, b)$ (``$a$ translated into the framework of
$b$'') has different semantic weight from $\compose(b, a)$ (``$b$
translated into the framework of $a$''). Translation is not
commutative, and the asymmetry measures the irreducible directionality
of cross-cultural encounter.
\end{interpretation}

\section{Emergence Is Meaning}

We conclude this chapter --- and Part~I of Volume~IV --- with the
theorem that crystallizes the central thesis of formal semiotics.

\begin{theorem}[Emergence Is Meaning]
\label{thm:emergence-is-meaning}
The emergence cocycle $\emerge$ is the fundamental measure of semiotic
meaning. A composition is meaningful (generates interpretive content
beyond the parts) if and only if its emergence is nonzero.
\end{theorem}

\begin{proof}
By the definition of emergence, $\emerge(a, b, c) = 0$ for all $c$ if
and only if $\rs(\compose(a,b), c) = \rs(a, c) + \rs(b, c)$ for all
$c$ --- that is, the composition is purely additive, generating no
surplus. Nonzero emergence is the necessary and sufficient condition
for the composition to generate meaning that transcends the parts.
See \texttt{emergence\_is\_meaning}.
\end{proof}

\begin{lstlisting}[caption={The fundamental theorem of formal semiotics}]
theorem emergence_is_meaning {a b : I} :
    (∃ c, emergence a b c ≠ 0) ↔ ¬naturalized_pair a b := ...
\end{lstlisting}

\begin{interpretation}[The Fundamental Theorem of Semiotics]
This theorem is the foundation stone of formal semiotics. It says that
\emph{meaning is emergence}: the meaning of a composition is not the
sum of the meanings of its parts but the \emph{surplus} generated by
their combination. When we combine two ideas and something new arises
--- something that was not present in either idea alone --- that is
meaning. When the combination is purely additive --- when the whole is
exactly the sum of the parts --- there is no meaning, only aggregation.

This identification of meaning with emergence unifies all the threads
of Part~I. Saussure's differential value is the connotation function,
which arises from differences in resonance. Peirce's interpretant is
the emergence of the sign-act. Barthes's myth is the emergence of
second-order composition, and his naturalization is the vanishing of
emergence. Eco's open text is the text with positive emergence, and
his overinterpretation is emergence that exceeds its bound. Lotman's
living sign is the sign with positive emergence, and his hegemonic
sign is the sign with zero emergence.

In every case, emergence $\emerge$ is the key: it is the formal
measure of semiotic creativity, the algebraic name for what happens
when two ideas come together and produce something genuinely new. The
$\IS$ is the space in which this production occurs; the axioms are the
laws that govern it; and the emergence cocycle is its central
invariant. With this theorem, we have completed the foundation of
formal semiotics. The chapters that follow (Volume~IV, Part~II) will
apply this foundation to specific domains: narrative, aesthetics,
translation, and the semiotics of power.
\end{interpretation}

\begin{theorem}[The Master Cocycle Identity]
\label{thm:master-cocycle}
For all $a, b, c, d \in \mathcal{I}$:
\[
  \emerge(\compose(a, b), c, d) + \emerge(a, b, d)
  = \emerge(a, \compose(b, c), d) + \emerge(b, c, d).
\]
\end{theorem}

\begin{proof}
Expand using the definition of $\emerge$ and the associativity
$\compose(\compose(a, b), c) = \compose(a, \compose(b, c))$:
\begin{align*}
  &\text{LHS} = [\rs(\compose(\compose(a,b),c), d) - \rs(\compose(a,b), d) - \rs(c, d)] \\
  &\qquad\quad + [\rs(\compose(a,b), d) - \rs(a, d) - \rs(b, d)] \\
  &= \rs(\compose(a, \compose(b,c)), d) - \rs(a, d) - \rs(b, d) - \rs(c, d).
\end{align*}
\begin{align*}
  &\text{RHS} = [\rs(\compose(a, \compose(b,c)), d) - \rs(a, d) - \rs(\compose(b,c), d)] \\
  &\qquad\quad + [\rs(\compose(b,c), d) - \rs(b, d) - \rs(c, d)] \\
  &= \rs(\compose(a, \compose(b,c)), d) - \rs(a, d) - \rs(b, d) - \rs(c, d).
\end{align*}
These are identical, completing the proof. See
\texttt{master\_cocycle}.
\end{proof}

\begin{interpretation}[The Coherence of Meaning]
The master cocycle identity is the structural backbone of the
semiosphere. It says that no matter how we parenthesize a triple
composition --- whether we first compose $a$ with $b$ and then
compose the result with $c$, or first compose $b$ with $c$ and
then compose $a$ with the result --- the total emergence is
consistent. The emergences at different levels of nesting are not
independent quantities but are bound together by an exact algebraic
identity. This is the deepest sense in which the $\IS$ is
\emph{coherent}: meaning does not depend on the order of assembly.
The cocycle identity ensures that semiotic analysis, like homological
algebra, is consistent under regrouping. It is the price of coherence
and the guarantee of well-definedness.

With this result, we have laid the algebraic foundations of formal
semiotics. The five thinkers of Part~I --- Saussure, Peirce, Barthes,
Eco, and Lotman --- have been brought within a single formal framework,
their key insights translated into theorems of the $\IS$, their
hidden connections revealed by the algebra. Part~II will build on
these foundations to develop the semiotics of narrative, aesthetics,
and culture.
\end{interpretation}

\section{The Unity of Formal Semiotics}

\subsection{A Comparative Table}

We conclude this chapter and Part~I with a systematic comparison of
the five semiotic frameworks, showing how each maps onto the structures
of the $\IS$.

\begin{center}
\begin{tabular}{lll}
\textbf{Theorist} & \textbf{Key Concept} & \textbf{$\IS$ Formalization} \\
\hline
Saussure & Sign = signifier + signified & $\compose(s_r, s_d)$ \\
Saussure & Arbitrariness & Independence of synonymy and $\rs$ \\
Saussure & Value (\textit{valeur}) & Connotation function \\
Saussure & Opposition & $\rs(a,c) + \rs(b,c) = 0$ \\
Peirce & Icon & $\rs(a, b) > 0$ \\
Peirce & Index & $\emerge(a, b, b) > 0$ \\
Peirce & Symbol & $\texttt{sameIdea}(a,b) \wedge \rs(a,b) = 0$ \\
Peirce & Interpretant & $\emerge(a, b, c)$ \\
Peirce & Habit & Naturalized idea \\
Barthes & Myth & $\compose(\compose(s_r, s_d), \gamma)$ \\
Barthes & Naturalization & $\emerge = 0$ universally \\
Barthes & Grain of the voice & Emergence of the sign \\
Barthes & Punctum & $\emerge(\rho, p, c) > 0$ \\
Eco & Open text & $\exists c,\; \emerge > 0$ \\
Eco & Closed text & $\forall c,\; \emerge \leq 0$ \\
Eco & Encyclopedia entry & $\exists b\,c,\; \emerge \neq 0$ \\
Eco & Overinterpretation & $|\emerge| > M\sqrt{\rs(a,a)\rs(b,b)}$ \\
Lotman & Semiosphere & The $\IS$ itself \\
Lotman & Center & Hegemonic signs \\
Lotman & Periphery & Living signs with high emergence \\
Lotman & Frozen sign & Naturalized idea \\
Lotman & Living sign & $\exists b\,c,\; \emerge > 0$ \\
Lotman & Explosion & Anomalously large $|\emerge|$ \\
Lotman & Boundary & Mediator of structural oppositions \\
\end{tabular}
\end{center}

\subsection{The Convergence of Traditions}

The table reveals a remarkable convergence: concepts that were developed
independently by thinkers from different intellectual traditions, in
different countries, in different languages, and for different purposes
turn out to be \emph{the same concept} when expressed in the $\IS$.
Barthes's ``naturalization'' is Peirce's ``habit,'' which is Lotman's
``frozen sign,'' which is Eco's ``dictionary entry'' --- all are
formalized as $\texttt{naturalized}(a)$, the condition that all
emergences involving $a$ vanish. Eco's ``open text'' is Lotman's
``living sign,'' which is Kristeva's ``semiotic modality,'' which is
Barthes's ``writerly text'' --- all are formalized as the existence
of positive emergence.

\begin{theorem}[The Unification Theorem]
\label{thm:unification}
The following conditions on an idea $a \in \mathcal{I}$ are equivalent:
\begin{enumerate}
\item $a$ is naturalized (Barthes);
\item $a$ is a semiotic habit (Peirce);
\item $a$ is a frozen sign (Lotman);
\item $a$ is a dictionary entry (Eco);
\item $a$ is left-linear: $\rs(\compose(a,b), c) = \rs(a,c) + \rs(b,c)$
for all $b, c$.
\end{enumerate}
The negation of these conditions gives:
\begin{enumerate}
\item $a$ is denaturalized (Barthes);
\item $a$ is an object of inquiry (Peirce);
\item $a$ is a living sign (Lotman);
\item $a$ is an encyclopedia entry (Eco);
\item $a$ participates in nonlinear composition.
\end{enumerate}
\end{theorem}

\begin{proof}
Conditions (1)--(4) are all defined as $\forall b\,c,\;
\emerge(a, b, c) = 0$, which is equivalent to (5) by the definition
of emergence. The equivalence is definitional in the $\IS$. See
\texttt{semiotic\_unification}.
\end{proof}

\begin{lstlisting}[caption={The semiotic unification theorem}]
theorem semiotic_unification (a : I) :
    naturalized a ↔ frozenSign a := Iff.rfl

theorem semiotic_unification' (a : I) :
    naturalized a ↔ dictionaryEntry a := Iff.rfl

theorem semiotic_unification'' (a : I) :
    naturalized a ↔ semioticHabit a := Iff.rfl
\end{lstlisting}

\begin{interpretation}[The Unity of Semiotics]
The unification theorem is the capstone of Part~I. It shows that the
five semiotic traditions of the twentieth century --- Saussure's
structuralism, Peirce's pragmaticism, Barthes's post-structuralism,
Eco's interpretive semiotics, and Lotman's cultural semiotics --- are
not five different theories but five \emph{perspectives} on a single
underlying structure: the $\IS$. The formal framework does not choose
between these perspectives; it unifies them, showing that their
apparent disagreements are disagreements about emphasis, not about
substance.

Saussure emphasizes the system of differences (resonance); Peirce
emphasizes the triadic process of interpretation (emergence); Barthes
emphasizes the ideological dimension (naturalization); Eco emphasizes
the dialectic of freedom and constraint (openness and bounds); Lotman
emphasizes the cultural dynamics (frozen/living, center/periphery).
All five are necessary; none is sufficient alone. The $\IS$ provides
the algebraic space in which all five coexist and interact.

This unity is not a reduction: we have not shown that Barthes ``really
means the same thing'' as Peirce, or that Lotman is ``just a
translation'' of Saussure. Each thinker brings unique insights, unique
examples, unique philosophical commitments. What the $\IS$ shows is
that these diverse insights can be expressed in a common language ---
the language of composition, resonance, and emergence --- and that this
common language reveals hidden connections, enables rigorous comparison,
and opens the way to theorems that no single tradition could have
proved alone.
\end{interpretation}

\subsection{Open Problems in Formal Semiotics}

We close Part~I by listing several open problems that arise from the
formalization and that point toward Part~II and future volumes.

\begin{enumerate}
\item \textbf{The classification of signs.} Peirce proposed up to
sixty-six classes of signs. Can these be systematically derived
from the axioms of the $\IS$? Which classes correspond to distinct
algebraic structures, and which collapse?

\item \textbf{The dynamics of naturalization.} Barthes described
naturalization as a \emph{process}: signs become naturalized over
time. Can we formalize the \emph{dynamics} of naturalization within
a time-varying $\IS$? What conditions cause a living sign to become
frozen?

\item \textbf{The topology of the semiosphere.} Lotman's center,
periphery, and boundary are topological notions. Can we equip the
$\IS$ with a genuine topology (or metric) such that these notions
have precise topological content?

\item \textbf{The ethics of emergence.} We defined the ethical
encounter as emergence without domination. Can we develop a full
\emph{ethics of semiosis} within the $\IS$, including notions of
justice, reciprocity, and care?

\item \textbf{The cohomology of meaning.} The cocycle identity
suggests that emergence is a $2$-cocycle in some cohomological
theory. Can we identify the relevant cohomology? What are the
higher cocycles?

\item \textbf{The computability of resonance.} Given a finite
presentation of the $\IS$ (e.g., generators and relations), is
the resonance function computable? Is the emergence cocycle
decidable?

\item \textbf{The semiotics of multimodality.} Real-world sign
systems are multimodal: they combine language, image, gesture,
sound. Can the $\IS$ accommodate multimodal composition, where
ideas from different ``channels'' are composed according to
different rules?
\end{enumerate}

These problems are not merely technical: they touch on some of the
deepest questions in the human sciences. The formalization of semiotics
within the $\IS$ is not an end but a \emph{beginning} --- the opening
of a new chapter in the ancient inquiry into the nature of signs, the
structure of meaning, and the possibility of communication.

\begin{remark}[A Note on Method]
The reader may wonder whether the formalization of semiotics in the
$\IS$ does violence to the humanistic spirit of the original theories.
We believe the opposite: formalization \emph{liberates} these theories
from the vagueness that has limited their influence and the
terminological confusion that has plagued interdisciplinary
communication. When Barthes writes ``naturalization,'' Peirce writes
``habit,'' and Lotman writes ``freezing,'' they are pointing at the
same phenomenon but cannot know it because they lack a common
language. The $\IS$ provides that language. It does not replace the
original theories but \emph{supplements} them with a formal dimension
that enables rigorous comparison, systematic generalization, and ---
most importantly --- \emph{proof}. In the $\IS$, semiotic claims are
not mere assertions to be accepted or rejected on the authority of
their proponents; they are \emph{theorems} to be proved from axioms,
checked by machine, and built upon by future researchers. This is the
promise of formal semiotics, and it is the promise we shall continue
to fulfill in Part~II.
\end{remark}

\begin{remark}[Summary of Part I]
Part~I has established the following:
\begin{itemize}
\item \textbf{75 theorems} proved and verified in Lean, formalizing
the core insights of five major semiotic traditions.
\item \textbf{52 definitions} translating humanistic concepts into
precise algebraic language within the $\IS$.
\item \textbf{One unification theorem} showing that independently
developed concepts (naturalization, habit, freezing, dictionary
entry) are algebraically identical.
\item \textbf{One master cocycle identity} that serves as the
structural backbone of the entire theory.
\item \textbf{Seven open problems} pointing toward future work in
formal semiotics.
\end{itemize}
The reader is invited to verify these results against the Lean source
code in \texttt{IDT\_v3\_Semiotics.lean}, where all 327 theorems of the
semiotic formalization are collected. Part~II begins with the
semiotics of narrative structure and proceeds through aesthetics,
translation theory, and the semiotics of power. The algebraic
foundations laid in Part~I will be presupposed throughout.
\end{remark}

\begin{remark}[Cross-Reference: Semiosphere and Heteroglossia]
The concept of the semiosphere developed in this chapter connects directly to Bakhtin's notion of \emph{heteroglossia}, which we formalize in Chapter~\ref{ch:novel}. Bakhtin's insight that every utterance exists at the intersection of multiple ``social languages'' is, in our framework, the claim that every idea $a \in \IS$ has a resonance profile that reflects its position in the semiosphere---its proximity to different cultural centers, its distance from different peripheries. As Volume~III's Wittgensteinian analysis showed (particularly the discussion of language games as local subsystems), meaning is always situated, always embedded in a form of life. The semiosphere is the totality of these forms of life, and the $\IS$ is its mathematical model.
\end{remark}

\begin{remark}[Cross-Reference: Emergence Across Volumes]
The emergence cocycle $\emerge$, which has been the central invariant of Part~I, connects to structures developed in other volumes. In Volume~I (Foundations), emergence was introduced as the deviation from linearity in the resonance function; here, it has been revealed as the formal measure of semiotic creativity itself. In Volume~II (Strategic Interaction), emergence governs the payoff structure of meaning games---the surplus generated by cooperative communication. In Volume~III (Philosophy), emergence is the Wittgensteinian ``showing'' that cannot be ``said'': the ineffable surplus that language generates beyond its propositional content. In Volume~V (Power), emergence will model the generative capacity of discourse---the ability of powerful speech acts to create the very reality they describe. In Volume~VI (Emergence), emergence will be studied reflexively: the emergence of the concept of emergence, the self-referential loop that makes the $\IS$ a model of its own foundations.
\end{remark}
